% Leçon 13 — Grammaire : Interrogation avec 怎么了et怎么 ; 不但……, 而且…… ; 因为……, 所以…… — Chinois Tle A — Cours signé OnBuch+
% Programme officiel MINESEC. Compiler avec : tectonic ZH13-grammaire-interrogation-avec-et.tex
\newcommand{\DOCMATIERE}{Chinois}
\newcommand{\DOCNIVEAU}{Tle A}
\newcommand{\DOCMODULE}{Module 2 : Environnement et santé}
\newcommand{\DOCLECON}{ZH13}
\newcommand{\DOCTITRE}{Leçon 13 — Grammaire : Interrogation avec 怎么了et怎么 ; 不但……, 而且…… ; 因为……, 所以……}
% OnBuch+ — préambule « cours signés » ENRICHI (compilé avec tectonic / XeLaTeX)
% Système visuel « Pop » d'OnBuch+ : fond crème (jamais blanc), bordures encre
% épaisses, ombres « dures » décalées, titres Archivo Black en capitales.
% Version MATHÉMATIQUES : graphes (pgfplots/tikz), boîtes pédagogiques
% (définition, propriété, méthode, attention, le savais-tu, exercice résolu,
% exercice, TP, bilan), page de garde et signature OnBuch+ ancrée en bas.
%
% Macros attendues dans main.tex AVANT \input{preamble} :
%   \DOCMATIERE  \DOCNIVEAU  \DOCMODULE  \DOCLECON (numéro)  \DOCTITRE
\documentclass[11pt]{article}

\usepackage{fontspec}
\usepackage[french]{babel} % français = langue principale ; le chinois via \zh{..} / textebox (xeCJK)
\usepackage{xeCJK}
\usepackage{pifont} % \ding{51} \ding{55} utilisés par les modèles
\usepackage[a4paper,margin=2cm,top=3.4cm,bottom=2.8cm,headheight=1.4cm,footskip=1.3cm]{geometry}
\usepackage{amsmath,amssymb,amsfonts}
\usepackage{mathtools} % \xmapsto, \coloneqq... souvent utilisés par les modèles
\usepackage{cancel} % simplifications barrées (bilans de forces)
% \abs{z} (module, valeur absolue) et \norm{u} (norme), utilisés par les modèles
\providecommand{\abs}{}\let\abs\relax\DeclarePairedDelimiter{\abs}{\lvert}{\rvert}
\providecommand{\norm}{}\let\norm\relax\DeclarePairedDelimiter{\norm}{\lVert}{\rVert}
\usepackage[table]{xcolor} % [table] : fournit aussi \rowcolors (alternance de lignes)
\usepackage{pagecolor}
\usepackage{tikz}
\usetikzlibrary{arrows.meta,positioning,calc,shapes.geometric,shapes.misc,shapes.arrows,decorations.pathmorphing,decorations.pathreplacing,decorations.markings,patterns,shadows,mindmap,trees,fit,backgrounds,matrix,intersections,angles,quotes}
\usepackage{pgfplots}
\usepackage[version=4]{mhchem} % équations nucléaires \ce{^{235}_{92}U}
\usepackage[european,siunitx]{circuitikz} % schémas électriques
\pgfplotsset{compat=1.18}
\usepgfplotslibrary{fillbetween,groupplots} % aires entre courbes (calcul intégral)
\usepackage{enumitem}
\usepackage{fancyhdr}
% [most] charge toutes les bibliothèques tcolorbox (listings, minted, raster,
% documentation...) et gonfle inutilement la mémoire TeX sur un document long
% (~300 boîtes) : on ne charge que ce qui est réellement utilisé.
\usepackage[skins,breakable]{tcolorbox}
\usepackage{graphicx}
\usepackage{colortbl}
\usepackage{booktabs}
\usepackage{tabularx}
\usepackage{diagbox} % case d'en-tête diagonale des tableaux à double entrée
% Colonnes extensibles centrée / gauche / droite (C, L, R), souvent utilisées par les modèles
\newcolumntype{C}{>{\centering\arraybackslash}X}
\newcolumntype{L}{>{\raggedright\arraybackslash}X}
\newcolumntype{R}{>{\raggedleft\arraybackslash}X}
\usepackage{array}
\usepackage{multicol}
\usepackage{multirow}
\usepackage{siunitx}
% parse-numbers=false : accepte aussi \SI{2,5 \times 10^{-3}}{...}
\sisetup{locale=FR,output-decimal-marker={,},group-minimum-digits=4,parse-numbers=false}
\DeclareSIUnit\ohm{\ensuremath{\symup{\Omega}}} % U+2126 absent de la police
\DeclareSIUnit\division{div} % réglages d'oscilloscope (ms/div, V/div)
\DeclareSIUnit\tr{tr} % tours (vitesse de rotation, ex. \SI{1500}{\tr\per\minute})
\DeclareSIUnit\molar{M} % concentration molaire (ex. \SI{3}{\molar}, \SI{1}{\micro\molar})
\DeclareSIUnit\an{an} % année (le modèle confond parfois avec \year, un compteur TeX) — ex. \SI{5}{\kilo\meter\per\an}
\DeclareSIUnit\FCFA{FCFA} % franc CFA (ex. \SI{500}{\FCFA\per\kilo\gram})
\DeclareSIUnit\degreeBrix{\textdegree Brix} % teneur en sucre d'un sirop/jus (réfractométrie)
\DeclareSIUnit\month{mois} % le modèle utilise parfois \month, un compteur TeX, comme unité de durée
\usepackage{qrcode}
% \degree : commande fréquemment utilisée par les modèles pour un angle ou une
% température en dehors de \SI{}{} (non fournie par les paquets chargés).
\providecommand{\degree}{^{\circ}}
% \qed (fin de démonstration), utilisé par les modèles sans amsthm
\providecommand{\qed}{\hfill$\square$}
% \barre : double barre des valeurs interdites dans un tableau de variations (tkz-tab)
\providecommand{\barre}{\ensuremath{\|}}
% environnement proof (amsthm non chargé) utilisé par les modèles
\ifdefined\proof\else\newenvironment{proof}[1][Démonstration]{\par\noindent\textit{#1.}\ }{\qed\par}\fi
\usepackage{lastpage}
\usepackage{hyperref}
\hypersetup{hidelinks}

% ============================================================
% Palette « Pop » OnBuch+ — mêmes hex que la classe P (plus_ui.dart)
% ============================================================
\definecolor{poporange}{HTML}{F59321}
\definecolor{popdark}{HTML}{E07A0C}
\definecolor{popcream}{HTML}{FFF6E8}
\definecolor{popink}{HTML}{1C1714}
\definecolor{poppurple}{HTML}{7A5AE0}
\definecolor{popgreen}{HTML}{2E9E6A}
\definecolor{poppink}{HTML}{E0457B}
\definecolor{popblue}{HTML}{2F7DE1}
\definecolor{popgold}{HTML}{C98A12}
\definecolor{popmuted}{HTML}{8A7F76}
\definecolor{popline}{HTML}{EADFD2}
\definecolor{popgray}{HTML}{8A7F76}
\colorlet{popgrey}{popgray}
% Alias tolérants (le modèle nomme parfois les couleurs différemment)
\colorlet{popwhite}{white}
\colorlet{popblack}{popink}
\colorlet{popred}{poppink}
\colorlet{popyellow}{popgold}
% Teintes claires (fonds de tableaux/encadrés) — le modèle nomme parfois
% les couleurs avec un suffixe L (ex. popblueL) sans les avoir définies.
\colorlet{poporangeL}{poporange!20}
\colorlet{popdarkL}{popdark!20}
\colorlet{poppurpleL}{poppurple!20}
\colorlet{popgreenL}{popgreen!20}
\colorlet{poppinkL}{poppink!20}
\colorlet{popblueL}{popblue!20}
\colorlet{popgoldL}{popgold!20}
\colorlet{popredL}{popred!20}
\colorlet{popgrayL}{popgray!20}
% Teintes claires pour fonds de tableaux / zones
\colorlet{popblueL}{popblue!10!white}
\colorlet{poporangeL}{poporange!14!white}
\colorlet{popgreenL}{popgreen!12!white}
\colorlet{poppurpleL}{poppurple!12!white}
\colorlet{poppinkL}{poppink!12!white}
\colorlet{popgoldL}{popgold!14!white}

\pagecolor{popcream}

% ============================================================
% Typographie
% ============================================================
\defaultfontfeatures{Ligatures=TeX}
\setmainfont{PlusJakartaSans-Medium.ttf}[Path=../../../fonts/,BoldFont=PlusJakartaSans-Bold.ttf]
% Symboles, API et emojis absents de la police principale : repli sur DejaVu Sans (caractères actifs)
\newfontface\symfont{DejaVuSans.ttf}[Path=../../../fonts/]
\makeatletter
\newcommand\symactive[1]{\catcode#1=13 \begingroup\lccode`\~=#1 \lowercase{\endgroup\def~}{{\symfont\char#1}}}
\newcommand\symblank[1]{\catcode#1=13 \begingroup\lccode`\~=#1 \lowercase{\endgroup\def~}{}}
\newcommand\symhyph[1]{\catcode#1=13 \begingroup\lccode`\~=#1 \lowercase{\endgroup\def~}{\mbox{-}}}
\newcount\symc
\newcommand\symrange[2]{\symc=#1 \@whilenum\symc<#2 \do{\expandafter\symactive\expandafter{\the\symc}\advance\symc by 1 }}
\newcommand\symblankrange[2]{\symc=#1 \@whilenum\symc<#2 \do{\expandafter\symblank\expandafter{\the\symc}\advance\symc by 1 }}
\makeatother
\symrange{"0250}{"02FF}\symactive{"2713}\symactive{"2714}\symactive{"2717}\symactive{"2718}\symactive{"2610}\symactive{"2611}\symactive{"2612}
\symrange{"2600}{"27BF}
\catcode"2705=13 \begingroup\lccode`\~="2705 \lowercase{\endgroup\def~}{{\symfont\char"2713}}\symblank{"26BD}\symblank{"26A1}
\symrange{"2460}{"257F}\symactive{"223C}\symactive{"2248}\symactive{"2260}
\symhyph{"2011}\symhyph{"2E1A}\symblankrange{"1F300}{"1FAFF}\symblank{"FE0F}
% Caractères chinois : WenQuanYi Zen Hei (sous-ensemble GB2312 + pinyin), gras/italique simulés
\setCJKmainfont{WenQuanYiZenHei-subset.ttf}[Path=../../../fonts/,AutoFakeBold=3,AutoFakeSlant=0.2]
\xeCJKsetup{CJKspace=false}
% Pinyin : voyelles à caron (ǎ ǐ ǒ ǔ ǖ ǘ ǚ ǜ) absentes de la police principale -> caractères actifs
\newfontface\pyfont{WenQuanYiZenHei-subset.ttf}[Path=../../../fonts/]
\newcommand\pyactive[1]{\catcode#1=13 \begingroup\lccode`\~=#1 \lowercase{\endgroup\def~}{{\pyfont\char#1}}}
\pyactive{"01CD}\pyactive{"01CE}\pyactive{"01CF}\pyactive{"01D0}\pyactive{"01D1}\pyactive{"01D2}\pyactive{"01D3}\pyactive{"01D4}\pyactive{"01D5}\pyactive{"01D6}\pyactive{"01D7}\pyactive{"01D8}\pyactive{"01D9}\pyactive{"01DA}\pyactive{"01DB}\pyactive{"01DC}
% Maths en sans-serif (Fira Math) avec chiffres et lettres droites de Plus
% Jakarta, pour que les formules se fondent dans le texte.
\usepackage[math-style=ISO,bold-style=ISO]{unicode-math}
\setmathfont{FiraMath-Regular.otf}[Path=../../../fonts/]
\setmathfont{PlusJakartaSans-Medium.ttf}[Path=../../../fonts/,range={up/num,up/{Latin,latin},\mathup}]
\usepackage{icomma} % virgule décimale sans espace en mode math
% Caractères mathématiques Unicode absents des polices de texte (Plus Jakarta,
% Archivo Black) : rendus par leur équivalent mathématique, en texte comme en
% formule — sinon ils s'affichent en carré vide (ex. « Arithmétique dans ℤ »).
\usepackage{newunicodechar}
\newunicodechar{ℝ}{\ensuremath{\mathbb{R}}}
\newunicodechar{ℤ}{\ensuremath{\mathbb{Z}}}
\newunicodechar{ℂ}{\ensuremath{\mathbb{C}}}
\newunicodechar{ℕ}{\ensuremath{\mathbb{N}}}
\newunicodechar{ℚ}{\ensuremath{\mathbb{Q}}}
\newunicodechar{𝔻}{\ensuremath{\mathbb{D}}}
\newunicodechar{α}{\ensuremath{\alpha}}
\newunicodechar{β}{\ensuremath{\beta}}
\newunicodechar{γ}{\ensuremath{\gamma}}
\newunicodechar{δ}{\ensuremath{\delta}}
\newunicodechar{ε}{\ensuremath{\varepsilon}}
\newunicodechar{θ}{\ensuremath{\theta}}
\newunicodechar{λ}{\ensuremath{\lambda}}
\newunicodechar{μ}{\ensuremath{\mu}}
\newunicodechar{σ}{\ensuremath{\sigma}}
\newunicodechar{τ}{\ensuremath{\tau}}
\newunicodechar{φ}{\ensuremath{\varphi}}
\newunicodechar{ψ}{\ensuremath{\psi}}
\newunicodechar{ω}{\ensuremath{\omega}}
\newunicodechar{Σ}{\ensuremath{\Sigma}}
% majuscules produites par \MakeUppercase dans les titres (α-aminés -> Α)
\newunicodechar{Α}{\ensuremath{\alpha}}
\newunicodechar{Β}{\ensuremath{\beta}}
\newunicodechar{Γ}{\ensuremath{\Gamma}}
\newunicodechar{Δ}{\ensuremath{\Delta}}
\newunicodechar{Ε}{\ensuremath{\varepsilon}}
\newunicodechar{Θ}{\ensuremath{\Theta}}
\newunicodechar{Λ}{\ensuremath{\Lambda}}
\newunicodechar{Μ}{\ensuremath{\mu}}
\newunicodechar{Τ}{\ensuremath{\tau}}
\newunicodechar{Φ}{\ensuremath{\Phi}}
\newunicodechar{Ψ}{\ensuremath{\Psi}}
\newunicodechar{Ω}{\ensuremath{\Omega}}
\newunicodechar{↦}{\ensuremath{\mapsto}}
\newunicodechar{∈}{\ensuremath{\in}}
\newunicodechar{∉}{\ensuremath{\notin}}
\newunicodechar{∧}{\ensuremath{\wedge}}
\newunicodechar{⇔}{\ensuremath{\Leftrightarrow}}
\newunicodechar{⟺}{\ensuremath{\Longleftrightarrow}}
\newunicodechar{⇒}{\ensuremath{\Rightarrow}}
\newunicodechar{⟹}{\ensuremath{\Longrightarrow}}
\newunicodechar{∘}{\ensuremath{\circ}}
\newunicodechar{∪}{\ensuremath{\cup}}
\newunicodechar{∩}{\ensuremath{\cap}}
\newunicodechar{≡}{\ensuremath{\equiv}}
\newunicodechar{⊂}{\ensuremath{\subset}}
\newunicodechar{⊥}{\ensuremath{\perp}}
\newunicodechar{∃}{\ensuremath{\exists}}
\newunicodechar{∀}{\ensuremath{\forall}}
\newunicodechar{∛}{\ensuremath{\sqrt[3]{\,}}}
\newunicodechar{∜}{\ensuremath{\sqrt[4]{\,}}}
\newunicodechar{⃗}{\ensuremath{{}^{\to}}}
\newunicodechar{ⁿ}{\ifmmode^{n}\else\textsuperscript{\ensuremath{n}}\fi}
\newunicodechar{ˣ}{\ifmmode^{x}\else\textsuperscript{\ensuremath{x}}\fi}
\newunicodechar{ᵇ}{\ifmmode^{b}\else\textsuperscript{\ensuremath{b}}\fi}
\newunicodechar{ʳ}{\ifmmode^{r}\else\textsuperscript{\ensuremath{r}}\fi}
\newunicodechar{ᵘ}{\ifmmode^{u}\else\textsuperscript{\ensuremath{u}}\fi}
\newunicodechar{ʸ}{\ifmmode^{y}\else\textsuperscript{\ensuremath{y}}\fi}
\newunicodechar{ᵃ}{\ifmmode^{a}\else\textsuperscript{\ensuremath{a}}\fi}
\newunicodechar{ᵉ}{\ifmmode^{e}\else\textsuperscript{\ensuremath{e}}\fi}
\newunicodechar{ᵏ}{\ifmmode^{k}\else\textsuperscript{\ensuremath{k}}\fi}
\newunicodechar{ᵖ}{\ifmmode^{p}\else\textsuperscript{\ensuremath{p}}\fi}
\newunicodechar{ᴺ}{\ifmmode^{N}\else\textsuperscript{\ensuremath{N}}\fi}
\newunicodechar{ᶜ}{\ifmmode^{c}\else\textsuperscript{\ensuremath{c}}\fi}
\newunicodechar{ʰ}{\ifmmode^{h}\else\textsuperscript{\ensuremath{h}}\fi}
\newunicodechar{ᵠ}{\ifmmode^{q}\else\textsuperscript{\ensuremath{q}}\fi}
\newunicodechar{ₙ}{\ifmmode_{n}\else\textsubscript{\ensuremath{n}}\fi}
\newunicodechar{ₐ}{\ifmmode_{a}\else\textsubscript{\ensuremath{a}}\fi}
\newunicodechar{ᵢ}{\ifmmode_{i}\else\textsubscript{\ensuremath{i}}\fi}
\newunicodechar{ᵣ}{\ifmmode_{r}\else\textsubscript{\ensuremath{r}}\fi}
\newunicodechar{ₖ}{\ifmmode_{k}\else\textsubscript{\ensuremath{k}}\fi}
\newunicodechar{ᵦ}{\ifmmode_{\beta}\else\textsubscript{\ensuremath{\beta}}\fi}
\newunicodechar{ₓ}{\ifmmode_{x}\else\textsubscript{\ensuremath{x}}\fi}
\newfontfamily\popdisplay{ArchivoBlack-Regular.ttf}[Path=../../../fonts/]
\newfontfamily\poplogo{SpaceGrotesk-Bold.ttf}[Path=../../../fonts/]
\newfontfamily\popsemi{PlusJakartaSans-SemiBold.ttf}[Path=../../../fonts/]
\newfontfamily\popxbold{PlusJakartaSans-ExtraBold.ttf}[Path=../../../fonts/]
\newfontfamily\popxboldls{PlusJakartaSans-ExtraBold.ttf}[Path=../../../fonts/,LetterSpace=3.0]

\setlist[enumerate]{leftmargin=1.7em,itemsep=5pt,topsep=4pt}
\setlist[itemize]{leftmargin=1.7em,itemsep=4pt,topsep=4pt,label={\color{poporange}\textbullet}}
\setlength{\parindent}{0pt}
\setlength{\parskip}{5pt}
\renewcommand{\arraystretch}{1.25}

% pgfplots : style Pop par défaut
\pgfplotsset{
  every axis/.append style={axis line style={popink,line width=1pt},tick style={popink},
    grid style={popline,line width=0.5pt},label style={font=\small},tick label style={font=\footnotesize}},
  popcurve/.style={poporange,line width=1.8pt,no markers,smooth},
}
% Même style disponible dans un \draw TikZ simple (hors pgfplots)
\tikzset{popcurve/.style={poporange,line width=1.8pt}}
\tikzset{popgrid/.style={popline,very thin}}
\colorlet{popcurve}{poporange} % le nom du style est parfois utilisé comme couleur

% ============================================================
% Pill / squiggle
% ============================================================
\newcommand{\poppill}[3]{%
  \tikz[baseline=-0.55ex]\node[draw=popink,line width=1.2pt,rounded corners=2.6pt,%
    fill=#2,inner xsep=6.5pt,inner ysep=3pt]{%
    \popxboldls\fontsize{7}{7}\selectfont\color{#3}\MakeUppercase{#1}};%
}
\newcommand{\popsquiggle}[1]{%
  \tikz[baseline=-0.4ex]{
    \draw[#1,line width=2.6pt,line cap=round]
      plot [smooth] coordinates {(0,0.09) (0.34,0.01) (0.68,0.21) (1.02,0.01) (1.36,0.10)};
  }%
}
% Mot-clé surligné (marqueur orange sous le texte)
\newcommand{\cle}[1]{{\popxbold\color{popdark}#1}}

% ============================================================
% En-tête / pied
% ============================================================
\pagestyle{fancy}
\fancyhf{}
\renewcommand{\headrulewidth}{0pt}
\renewcommand{\footrulewidth}{0pt}
\lhead{\raisebox{-0.15ex}{\poplogo\fontsize{13}{13}\selectfont\color{popink}OnBuch\color{poporange}+}}
\rhead{\poppill{\DOCMATIERE\ \textbullet\ \DOCNIVEAU\ \textbullet\ Leçon \DOCLECON}{white}{popink}}
\lfoot{\popxbold\fontsize{7.6}{7.6}\selectfont\color{popmuted}\MakeUppercase{Cours signé OnBuch+}\ \textbullet\ {\popsemi\DOCTITRE}}
\rfoot{\tikz[baseline=-0.6ex]\node[rounded corners=8.5pt,fill=popink,minimum height=17pt,minimum width=17pt,inner xsep=5pt,inner sep=0]{\popxbold\fontsize{8}{8}\selectfont\color{popcream}\thepage\,/\,\pageref*{LastPage}};}

% ============================================================
% Titres
% ============================================================
\newcommand{\chaptitre}[2]{%
  \begin{tcolorbox}[enhanced,colback=poporange,colframe=popink,boxrule=2.6pt,arc=11pt,
      drop shadow=popink,left=15pt,right=15pt,top=12pt,bottom=11pt,before skip=3pt,after skip=18pt]
    \poppill{#2}{popink}{popcream}\par\vspace{9pt}
    {\raggedright\hyphenpenalty=10000\exhyphenpenalty=10000\popdisplay\fontsize{22}{24}\selectfont\color{popcream}\MakeUppercase{#1}\par}
  \end{tcolorbox}
}
% Section numérotée : pastille orange avec le numéro + titre Archivo
\newcounter{popsec}
\newcommand{\coursec}[1]{%
  \stepcounter{popsec}\setcounter{popsub}{0}%
  \par\vspace{16pt}\noindent
  \begin{minipage}{\linewidth}
  \tikz[baseline=-0.7ex]\node[draw=popink,line width=1.6pt,fill=poporange,rounded corners=4pt,minimum size=22pt,inner sep=0]{\popdisplay\fontsize{12}{12}\selectfont\color{popcream}\thepopsec};%
  \hspace{8pt}{\popdisplay\fontsize{15}{17}\selectfont\color{popink}\MakeUppercase{#1}}\par
  \vspace{4pt}\noindent{\color{poporange}\rule{36pt}{3.6pt}}
  \end{minipage}\par\nopagebreak\vspace{8pt}
}
\newcounter{popsub}
\newcommand{\courssub}[1]{%
  \stepcounter{popsub}%
  \par\vspace{9pt}\noindent{\popxbold\fontsize{11.5}{13}\selectfont\color{popink}{\color{poporange}\thepopsec.\thepopsub}\hspace{6pt}#1}\par\nopagebreak\vspace{4pt}
}

% ============================================================
% Boîtes pédagogiques — même recette : fond blanc, bordure épaisse colorée,
% ombre dure encre, badge plein. Titre optionnel [#1].
% ============================================================
\tcbset{popbase/.style={breakable,enhanced,colback=white,boxrule=2.3pt,arc=9pt,
  shadow={3pt}{-3pt}{0pt}{popink},left=14pt,right=14pt,top=11pt,bottom=11pt,before skip=11pt,after skip=15pt,
  fontupper=\popsemi\fontsize{10.6}{14.5}\selectfont\color{popink},
  fontlower=\popsemi\fontsize{10.6}{14.5}\selectfont\color{popink}}}
% Coche dessinée (le glyphe ✓ n'existe pas dans les polices du document)
\newcommand{\popcheck}[1]{\tikz[baseline=-0.5ex]\draw[#1,line width=1.6pt,line cap=round,line join=round](-0.13,0.0)--(-0.04,-0.09)--(0.13,0.1);}
\providecommand{\coche}{\popcheck{popgreen}}
\providecommand{\resultat}[1]{\par\smallskip\noindent\fcolorbox{popgreen}{popgreenL}{\parbox{\dimexpr\linewidth-2\fboxsep-2\fboxrule\relax}{\textbf{Résultat.} #1}}\par\smallskip}
\newcommand{\popboxhead}[3]{\poppill{#1}{#2}{white}\ifx\relax#3\relax\else\hspace{6pt}{\popxbold\fontsize{10.6}{12}\selectfont\color{popink}#3}\fi\par\vspace{7pt}}

\newtcolorbox{aretenir}[1][]{popbase,colframe=popgreen,before upper={\popboxhead{À retenir}{popgreen}{#1}}}
% Texte en chinois (document de lecture, dialogue, extrait) : cadre
% d'étude. Un titre optionnel (source, auteur) s'affiche dans l'en-tête.
\newtcolorbox{textebox}[1][]{popbase,colframe=popblue,before upper={\popboxhead{文本 · Texte}{popblue}{#1}},after upper={}}
% Marques ✓ / ✗ (forme correcte / fautive) souvent utilisées par les modèles
\providecommand{\cross}{\textcolor{poppink}{\ensuremath{\times}}}
\providecommand{\xmark}{\cross}\providecommand{\crossmark}{\cross}\providecommand{\cmark}{\ensuremath{\checkmark}}
% Ligne de réponse / justification à compléter (inventée par les modèles)
\providecommand{\justification}[1]{\par\noindent\textit{Justification~:} \dotfill\par}
% \textipa (package tipa non chargé) : la police ne couvre pas l'API -> simple passage
\providecommand{\textipa}[1]{#1}
% Soulignement (ulem non chargé) : \uline, \ul
\providecommand{\uline}[1]{\underline{#1}}\providecommand{\ul}[1]{\underline{#1}}
% \glossaire{\item ...} inventée par les modèles : liste de glossaire
\providecommand{\glossaire}[1]{\begin{itemize}#1\end{itemize}}
% \alert (beamer) inventée par les modèles : texte mis en évidence
\providecommand{\alert}[1]{\textbf{\textcolor{poppink}{#1}}}
% variantes ASCII d'ordinaux écrites en macro
\providecommand{\iere}{\textsuperscript{re}}\providecommand{\ieme}{\textsuperscript{e}}\providecommand{\ere}{\textsuperscript{re}}\providecommand{\eme}{\textsuperscript{e}}\providecommand{\er}{\textsuperscript{er}}\providecommand{\ie}{\textsuperscript{e}}
% Ordinaux français écrits en macro par les modèles : 1\ière, 3\ième
\newcommand{\ière}{\textsuperscript{re}}\newcommand{\ième}{\textsuperscript{e}}
% \exemple (inventée par les modèles) : étiquette « Exemple : »
\providecommand{\exemple}{\textbf{Exemple~:}\ }
% \square utilisable aussi hors mode math (cases à cocher)
\AtBeginDocument{\let\squareMath\square\renewcommand{\square}{\ensuremath{\squareMath}}}
% Mot ou phrase en chinois à l'intérieur d'un paragraphe français
\newcommand{\zh}[1]{\ifmmode\text{#1}\else{#1}\fi}
\let\eng\zh \let\esp\zh \let\all\zh % alias tolérés
% flèches d'intonation inventées par les modèles
\providecommand{\textsearrow}{}\renewcommand{\textsearrow}{\ensuremath{\searrow}}\providecommand{\textnearrow}{}\renewcommand{\textnearrow}{\ensuremath{\nearrow}}
% \fr{..} : mot français (inventé par les modèles)
\providecommand{\fr}[1]{\textit{#1}}
% zhbox : encadré de texte chinois inventé par les modèles
\newenvironment{zhbox}{\begin{quote}}{\end{quote}}
% \vs : « versus » inventé par les modèles
\providecommand{\vs}{\textit{vs}\ }
% \blank : case à compléter inventée par les modèles
\providecommand{\blank}[1][]{\underline{\hspace{1.6em}}}
\newtcolorbox{exemplebox}[1][]{popbase,colframe=poporange,before upper={\popboxhead{Exemple}{poporange}{#1}}}
\newtcolorbox{definition}[1][]{popbase,colframe=popblue,before upper={\popboxhead{Définition}{popblue}{#1}}}
\newtcolorbox{propriete}[1][]{popbase,colframe=poppurple,before upper={\popboxhead{Propriété}{poppurple}{#1}}}
\newtcolorbox{methode}[1][]{popbase,colframe=popgold,before upper={\popboxhead{Méthode}{popgold}{#1}}}
\newtcolorbox{attention}[1][]{popbase,colframe=poppink,before upper={\popboxhead{Attention}{poppink}{#1}}}
\newtcolorbox{experience}[1][]{popbase,colframe=popblue,colback=popblueL,before upper={\popboxhead{Expérience}{popblue}{#1}}}
% « Le savais-tu ? » : carte violette pleine (contexte, culture, Cameroun)
\newtcolorbox{savaistu}[1][]{popbase,colback=poppurple,colframe=popink,
  fontupper=\popsemi\fontsize{10.4}{14.2}\selectfont\color{white},
  before upper={\poppill{Le savais-tu\,?}{white}{poppurple}\ifx\relax#1\relax\else\hspace{6pt}{\popxbold\color{white}#1}\fi\par\vspace{7pt}}}
% Exercice résolu : énoncé en haut, \tcblower puis la solution sur fond crème
\newcounter{popexo}\newcounter{popexr}
\newtcolorbox{exoresolu}[1][]{popbase,colframe=poporange,colbacklower=popcream,
  segmentation style={popink,line width=1.2pt,dashed},
  before upper={\stepcounter{popexr}\popboxhead{Exercice résolu \thepopexr}{poporange}{#1}},
  before lower={\poppill{Solution}{popink}{popcream}\par\vspace{6pt}}}
% Exercice (non résolu en place) — la correction est dans la partie Corrigés
\newtcolorbox{exercice}[1][]{popbase,colframe=popink,
  before upper={\stepcounter{popexo}\popboxhead{Exercice \thepopexo}{popink}{#1}}}
% Corrigé
\newtcolorbox{corrige}[1][]{popbase,colframe=popgreen,colback=popgreenL,
  before upper={\popboxhead{Corrigé}{popgreen}{#1}}}
% Objectifs / prérequis / bilan
\newtcolorbox{objectifs}{popbase,colframe=popink,colback=poporangeL,
  before upper={\popboxhead{Objectifs de la leçon}{popink}{}}}
\newtcolorbox{prerequis}{popbase,colframe=popink,colback=popblueL,
  before upper={\popboxhead{Prérequis}{popblue}{}}}
\newtcolorbox{bilan}{popbase,colframe=popink,colback=white,boxrule=2.8pt,
  before upper={\popboxhead{Fiche bilan}{poporange}{L'essentiel à connaître pour l'examen}}}
% Figure encadrée avec légende
\newtcolorbox{popfigure}[1][]{enhanced,breakable=false,colback=white,colframe=popline,boxrule=1.4pt,arc=8pt,
  left=10pt,right=10pt,top=10pt,bottom=8pt,before skip=10pt,after skip=12pt,halign=center,
  lower separated=false,halign lower=center,
  fontlower=\popsemi\fontsize{9}{11.5}\selectfont\color{popmuted},
  #1}
\newcommand{\legende}[1]{\par\vspace{6pt}{\popsemi\fontsize{9}{11.5}\selectfont\color{popmuted}#1}}

% ============================================================
% Signature OnBuch+
% ============================================================
\newcommand{\poplogoinline}[1]{{\poplogo\fontsize{#1}{#1}\selectfont\color{popink}OnBuch\color{poporange}+}}
\newcommand{\signatureonbuch}{%
  \begin{tcolorbox}[enhanced,colback=popink,colframe=popink,boxrule=2.2pt,arc=10pt,
      drop shadow=poporange,left=14pt,right=14pt,top=10pt,bottom=10pt,before skip=0pt,after skip=0pt]
    \tikz[baseline=-0.7ex]\node[circle,fill=popgreen,draw=popcream,line width=1.3pt,minimum size=22pt,inner sep=0]{\popcheck{white}};%
    \hspace{9pt}\begin{minipage}[c]{0.62\linewidth}
      {\popxbold\fontsize{12}{13}\selectfont\color{popcream}Signé\ \textbullet\ {\poplogo OnBuch\color{poporange}+}}\par\vspace{2pt}
      {\raggedright\popsemi\fontsize{8}{10.5}\selectfont\color{popline}\MakeUppercase{Équipe pédagogique OnBuch+}\\\MakeUppercase{Conforme au programme officiel MINESEC}\par}
    \end{minipage}\hfill
    {\poplogo\fontsize{22}{22}\selectfont\color{popcream}OnBuch\color{poporange}+}
  \end{tcolorbox}%
}

% ============================================================
% Page de garde
% #1 titre · #2 matière · #3 niveau · #4 module · #5 n° leçon · #6 durée ·
% #7 description / objectifs (texte ou itemize)
% ============================================================
\newcommand{\pagegarde}[7]{%
  \thispagestyle{empty}
  \newgeometry{a4paper,margin=1.7cm,top=1.6cm,bottom=1.4cm}%
  \begin{tikzpicture}[remember picture,overlay]
    \fill[poporange!18] ([xshift=-3.2cm,yshift=-3.6cm]current page.north east) circle (4.6cm);
    \fill[poppurple!14] ([xshift=2.4cm,yshift=7.6cm]current page.south west) circle (3.8cm);
  \end{tikzpicture}%
  {\poplogo\fontsize{30}{30}\selectfont\color{popink}OnBuch\color{poporange}+}\hfill
  \raisebox{0.6ex}{\poppill{Cours signé \textbullet\ Premium}{popink}{popcream}}\par
  \vspace{1pt}\popsquiggle{poppurple}\par
  \vspace{8pt}
  \begin{tcolorbox}[enhanced,colback=poporange,colframe=popink,boxrule=2.8pt,arc=13pt,
      drop shadow=popink,left=18pt,right=18pt,top=12pt,bottom=13pt,after skip=10pt]
    \poppill{#2 \textbullet\ #3}{popink}{popcream}\hspace{5pt}\poppill{#4}{white}{popink}\par\vspace{8pt}
    {\popdisplay\fontsize{14}{14}\selectfont\color{popink}LEÇON #5}\par\vspace{4pt}
    {\raggedright\hyphenpenalty=10000\exhyphenpenalty=10000\popdisplay\fontsize{25}{27}\selectfont\color{popcream}\MakeUppercase{#1}\par}
    \vspace{8pt}
    {\popxbold\fontsize{9.5}{9.5}\selectfont\color{popink}\MakeUppercase{Durée conseillée : #6}}
  \end{tcolorbox}
  \begin{tcolorbox}[enhanced,colback=white,colframe=popink,boxrule=2.2pt,arc=10pt,
      drop shadow=popink,left=16pt,right=16pt,top=10pt,bottom=10pt,after skip=8pt]
    \poppill{Dans cette leçon}{poppurple}{white}\par\vspace{5pt}
    \popsemi\fontsize{9.3}{12.6}\selectfont\color{popink}#7
  \end{tcolorbox}
  \noindent\begin{minipage}[t]{0.487\linewidth}
    \begin{tcolorbox}[enhanced,colback=popgreen,colframe=popink,boxrule=2.2pt,arc=10pt,
        drop shadow=popink,left=11pt,right=11pt,top=10pt,bottom=10pt,height=3.1cm,valign=center]
      \tcbox[colback=white,colframe=popink,boxrule=1.3pt,arc=5pt,boxsep=0pt,nobeforeafter,
        left=3pt,right=3pt,top=3pt,bottom=3pt,valign=center]{\qrcode[height=1.9cm]{https://play.google.com/store/apps/details?id=com.luvvix.onbuch&hl=fr}}%
      \hspace{8pt}\begin{minipage}[c]{0.5\linewidth}
        {\popdisplay\fontsize{11}{12.5}\selectfont\color{white}\MakeUppercase{Télécharge OnBuch}}\par\vspace{4pt}
        {\raggedright\popsemi\fontsize{8.4}{11}\selectfont\color{white}Gratuit \textbullet\ Google Play\\Tuteur IA, annales, concours, résultats\par}
      \end{minipage}
    \end{tcolorbox}
  \end{minipage}\hfill
  \begin{minipage}[t]{0.487\linewidth}
    \begin{tcolorbox}[enhanced,colback=poppurple,colframe=popink,boxrule=2.2pt,arc=10pt,
        drop shadow=popink,left=11pt,right=11pt,top=10pt,bottom=10pt,height=3.1cm,valign=center]
      \tikz[baseline=-0.7ex]\node[circle,fill=white,draw=popink,line width=1.3pt,minimum size=1.6cm,inner sep=0]{\popdisplay\fontsize{24}{24}\selectfont\color{poppurple}+};%
      \hspace{8pt}\begin{minipage}[c]{0.6\linewidth}
        {\popdisplay\fontsize{11}{12.5}\selectfont\color{white}\MakeUppercase{Passe à OnBuch+}}\par\vspace{4pt}
        {\raggedright\popsemi\fontsize{8.4}{11}\selectfont\color{white}Tous les cours signés, Tuteur IA illimité, fiches PDF — onglet OnBuch+ de l'app\par}
      \end{minipage}
    \end{tcolorbox}
  \end{minipage}
  \vspace{8pt}
  \popcontenu
  \vfill
  \signatureonbuch
  \restoregeometry
  \newpage
}

% Rangée « Ce cours contient » (page de garde)
\newcommand{\poptile}[3]{% #1 couleur #2 gros texte #3 légende
  \begin{tcolorbox}[enhanced,colback=#1,colframe=popink,boxrule=2pt,arc=9pt,shadow={2.5pt}{-2.5pt}{0pt}{popink},
      left=6pt,right=6pt,top=9pt,bottom=9pt,halign=center,valign=center,height=1.75cm,width=\linewidth,nobeforeafter]
    {\popdisplay\fontsize{12.5}{13}\selectfont\color{white}\MakeUppercase{#2}}\par\vspace{3pt}
    {\centering\popsemi\fontsize{7.8}{9.5}\selectfont\color{white}#3\par}
  \end{tcolorbox}}
\newcommand{\popcontenu}{%
  \par\noindent{\popxboldls\fontsize{7.5}{7.5}\selectfont\color{popmuted}CE COURS SIGNÉ CONTIENT}\par\vspace{7pt}
  \noindent\begin{minipage}[t]{0.235\linewidth}\poptile{popblue}{Cours}{complet, illustré, pas à pas}\end{minipage}\hfill
  \begin{minipage}[t]{0.235\linewidth}\poptile{poporange}{Méthodes}{et exercices résolus}\end{minipage}\hfill
  \begin{minipage}[t]{0.235\linewidth}\poptile{poppink}{Exercices}{type examen + corrigés}\end{minipage}\hfill
  \begin{minipage}[t]{0.235\linewidth}\poptile{popgreen}{Fiche bilan}{l'essentiel à retenir}\end{minipage}\par}

% Sommaire
\newtcolorbox{sommaire}{enhanced,colback=white,colframe=popink,boxrule=2.2pt,arc=10pt,
  drop shadow=popink,left=18pt,right=18pt,top=13pt,bottom=13pt,before skip=6pt,after skip=20pt,
  before upper={\poppill{Sommaire}{poppurple}{white}\par\vspace{9pt}\popsemi\fontsize{10.6}{14.5}\selectfont\color{popink}}}

% Signature de fin de cours — poussée tout en bas de la dernière page
\newcommand{\findecours}{%
  \par\vfill\nopagebreak
  \begin{center}\popsquiggle{poporange}\end{center}\vspace{4pt}
  \signatureonbuch
}
\AtBeginDocument{\def\llbracket{\mathopen{[\mkern-3.5mu[}}\def\rrbracket{\mathclose{]\mkern-3.5mu]}}} % intervalles d'entiers (glyphe ⟦ absent de la police)
% Glyphes absents de Fira Math : redéfinis à partir de glyphes disponibles
\AtBeginDocument{%
  \def\setminus{\mathbin{\backslash}}\let\smallsetminus\setminus
  \def\vdots{\mathord{\vbox{\baselineskip3.2pt\lineskiplimit0pt\kern4pt\hbox{.}\hbox{.}\hbox{.}}}}%
  \def\ddots{\mathinner{\mkern1mu\raise6.5pt\hbox{.}\mkern2mu\raise3.6pt\hbox{.}\mkern2mu\raise0.7pt\hbox{.}\mkern1mu}}%
  \def\triangle{\mathord{\scalebox{0.9}{$\symup{\Delta}$}}}%
  \def\nabla{\mathord{\raisebox{\depth}{\scalebox{1}[-1]{$\symup{\Delta}$}}}}%
  \def\checkmark{\popcheck{popdark}}%
  \def\star{\mathbin{\ast}}\def\bigstar{\mathord{\ast}}%
  \def\diamond{\mathbin{\scalebox{0.6}{$\Diamond$}}}%
  \def\bigcup{\mathop{\mathchoice{\raisebox{-0.3em}{\scalebox{1.7}{$\cup$}}}{\scalebox{1.25}{$\cup$}}{\cup}{\cup}}\displaylimits}%
  \def\bigcap{\mathop{\mathchoice{\raisebox{-0.3em}{\scalebox{1.7}{$\cap$}}}{\scalebox{1.25}{$\cap$}}{\cap}{\cap}}\displaylimits}%
  \def\lesseqqgtr{\mathrel{\lessgtr}}\def\bigoplus{\mathop{\oplus}\displaylimits}%
}
\providecommand{\ssi}{si et seulement si} % abréviation fréquente
\AtBeginDocument{\providecommand{\wideparen}{\overparen}} % arc de cercle

%% FIGFIX-BEGIN v5 — correctifs globaux des figures (docs/onbuch-generation/tools/figfix)
\usepackage{adjustbox}\usepackage{etoolbox}\tcbuselibrary{raster}
% toute figure TikZ/pgfplots plus large que la ligne est réduite à la largeur disponible
\BeforeBeginEnvironment{tikzpicture}{\begin{adjustbox}{max width=\linewidth}}
\AfterEndEnvironment{tikzpicture}{\end{adjustbox}}
% légendes pgfplots lisibles par-dessus les courbes
\pgfplotsset{every axis/.append style={legend style={fill=white,fill opacity=0.92,text opacity=1,draw=black!15,inner sep=3pt}}}
% étiquettes d'arêtes (syntaxe edge["..."]) avec fond blanc
\tikzset{every edge quotes/.append style={fill=white,inner sep=1.2pt}}
%% FIGFIX-END
\begin{document}

\pagegarde{Leçon 13 — Grammaire : Interrogation avec 怎么了et怎么 ; 不但……, 而且…… ; 因为……, 所以……}{Chinois}{Tle A}{Module 2 : Environnement et santé}{ZH13}{2 h}{Cette leçon de grammaire étudie trois structures essentielles du chinois : l'interrogation avec 怎么了 et 怎么, la corrélation additive 不但……而且…… et la corrélation causale 因为……所以……. Elle prépare directement aux épreuves de traduction et de production écrite du Baccalauréat.
\vspace{4pt}
\begin{multicols}{2}\small
\begin{itemize}
\item À la fin de cette leçon, je sais distinguer 怎么了 (que se passe-t-il ?) et 怎么 (comment ? pourquoi ?)
\item Je sais construire une question correcte avec 怎么了 et 怎么 en respectant l'ordre des mots
\item Je sais employer la corrélation 不但……而且…… pour exprimer une addition progressive
\item Je sais employer la corrélation 因为……所以…… pour exprimer la cause et la conséquence
\item Je sais placer correctement les corrélations par rapport au sujet selon que les deux propositions ont le même sujet ou non
\item Je sais éviter les erreurs typiques des francophones (double sujet mal placé, traduction mot à mot, oubli de 了)
\end{itemize}\end{multicols}}

\begin{sommaire}
\begin{multicols}{2}
\begin{enumerate}
\item Observation et découverte
\item Interroger avec 怎么了 et 怎么
\item La corrélation 不但……而且……
\item La corrélation 因为……所以……
\item Comparaison avec le français et pièges
\item Entraînement guidé et transformation
\item Activité d'intégration
\item Méthodes et exercices résolus
\item Exercices
\item Corrigés des exercices
\item Fiche bilan
\end{enumerate}
\end{multicols}
\end{sommaire}

\begin{objectifs}
\begin{itemize}
\item À la fin de cette leçon, je sais distinguer 怎么了 (que se passe-t-il ?) et 怎么 (comment ? pourquoi ?)
\item Je sais construire une question correcte avec 怎么了 et 怎么 en respectant l'ordre des mots
\item Je sais employer la corrélation 不但……而且…… pour exprimer une addition progressive
\item Je sais employer la corrélation 因为……所以…… pour exprimer la cause et la conséquence
\item Je sais placer correctement les corrélations par rapport au sujet selon que les deux propositions ont le même sujet ou non
\item Je sais éviter les erreurs typiques des francophones (double sujet mal placé, traduction mot à mot, oubli de 了)
\item Je sais utiliser ces trois structures dans des phrases et un court texte sur le thème de la santé
\item Je sais traduire du français vers le chinois des phrases contenant ces structures
\end{itemize}
\end{objectifs}

\begin{prerequis}
\begin{itemize}
\item L'ordre de base de la phrase chinoise Sujet–Verbe–Complément (Première)
\item Les pronoms personnels et les verbes courants (是, 有, 去, 吃, 学习)
\item Les mots interrogatifs simples : 什么, 谁, 哪儿, 为什么
\item La particule 了 de changement d'état (fin de phrase)
\item Le vocabulaire de la santé vu dans le Module 2 (生病, 感冒, 医院, 医生)
\end{itemize}
\end{prerequis}

\newpage

\coursec{Observation et découverte}

\courssub{Situation d'accroche : une consultation à l'hôpital de Bamenda}

Imaginez la scène. À l'hôpital de district de Bamenda, un médecin chinois coopérant reçoit un patient camerounais. La consultation commence par une question toute simple :

\zh{医生：你怎么了？} \emph{Yīshēng : Nǐ zěnme le ?} « Médecin : Qu'est-ce qui ne va pas ? »

Le patient explique son état :

\zh{病人：因为天气很冷，所以我感冒了。} \emph{Bìngrén : Yīnwèi tiānqì hěn lěng, suǒyǐ wǒ gǎnmào le.} « Patient : Parce qu'il fait très froid, j'ai attrapé un rhume. »

Le médecin termine par un conseil en deux parties :

\zh{医生：你不但要多喝水，而且要好好休息。} \emph{Yīshēng : Nǐ búdàn yào duō hē shuǐ, érqiě yào hǎohao xiūxi.} « Médecin : Non seulement tu dois boire beaucoup d'eau, mais encore tu dois bien te reposer. »

Trois phrases, trois structures grammaticales, une consultation complète : questionner avec \zh{怎么了}, expliquer une cause avec \zh{因为……所以……}, ajouter une information avec \zh{不但……而且……}. C'est exactement ce que le Bac attend de vous : savoir reconnaître, construire et employer ces trois outils dans un texte et dans une conversation.

\textbf{Problématique :} comment le chinois pose-t-il une question sur l'état ou la manière, exprime-t-il une addition d'idées et une relation de cause à conséquence, et comment éviter les pièges que le français nous tend ?

\courssub{Écoute et lecture à voix haute : six exemples modèles}

Lisez chaque exemple à voix haute, en respectant les tons, avant de regarder la traduction. Ces six phrases sont les modèles de toute la leçon.

\begin{exemplebox}[Six exemples modèles à lire à voix haute]
1. \zh{你怎么了？} \emph{Nǐ zěnme le ?} « Qu'est-ce que tu as ? / Qu'est-ce qui ne va pas ? »\\[2pt]
2. \zh{他今天怎么了？} \emph{Tā jīntiān zěnme le ?} « Qu'a-t-il aujourd'hui ? »\\[2pt]
3. \zh{这个字怎么写？} \emph{Zhège zì zěnme xiě ?} « Comment s'écrit ce caractère ? »\\[2pt]
4. \zh{你怎么不去上课？} \emph{Nǐ zěnme bú qù shàngkè ?} « Pourquoi ne vas-tu pas en cours ? »\\[2pt]
5. \zh{他不但会说汉语，而且会说英语。} \emph{Tā búdàn huì shuō Hànyǔ, érqiě huì shuō Yīngyǔ.} « Non seulement il sait parler chinois, mais encore il sait parler anglais. »\\[2pt]
6. \zh{因为下雨了，所以我不去学校。} \emph{Yīnwèi xià yǔ le, suǒyǐ wǒ bú qù xuéxiào.} « Parce qu'il pleut, je ne vais pas à l'école. »
\end{exemplebox}

\begin{methode}[Comment observer une structure grammaticale]
Face à une phrase chinoise nouvelle, suivez toujours ces quatre étapes :
\begin{enumerate}
\item \textbf{Lire} le chinois à voix haute, en respectant les tons, sans regarder la traduction.
\item \textbf{Souligner} les mots-outils : ici \zh{怎么}, \zh{了}, \zh{不但}, \zh{而且}, \zh{因为}, \zh{所以}.
\item \textbf{Comparer} avec la traduction française : quel mot français correspond à chaque mot-outil ? Y a-t-il des mots français sans équivalent chinois, ou l'inverse ?
\item \textbf{Dégager le schéma} : écrire la structure sous forme de formule, par exemple Sujet + 怎么了 ?
\end{enumerate}
Cette démarche d'observation est celle de l'épreuve de grammaire au Bac : on ne vous demande pas de réciter une règle, mais de la découvrir dans des phrases.
\end{methode}

\courssub{Observation guidée : souligner, repérer, comparer}

Reprenons les six exemples avec la loupe du grammairien. Pour chaque structure, soulignez mentalement les mots-outils et observez la place du sujet.

\textbf{Premier constat : la question avec \zh{怎么了}.} Dans \zh{你怎么了？}, le sujet \zh{你} vient d'abord, puis \zh{怎么}, puis la particule \zh{了} en fin de phrase. Le français dit « Qu'est-ce que tu as ? » avec un verbe ; le chinois n'a pas besoin de verbe : \zh{怎么了} fonctionne comme un bloc interrogatif complet. Dans \zh{他今天怎么了？}, le complément de temps \zh{今天} se place entre le sujet et \zh{怎么了}.

\textbf{Deuxième constat : la question avec \zh{怎么} seul.} Dans \zh{这个字怎么写？}, \zh{怎么} se place juste avant le verbe \zh{写} : il demande la manière (« comment ? »). Dans \zh{你怎么不去上课？}, le même mot devant \zh{不去} demande la raison (« pourquoi ne... pas ? »). Un seul mot, deux emplois : c'est la position devant le verbe qui reste constante.

\textbf{Troisième constat : la corrélation \zh{不但……而且……}.} Les deux moitiés encadrent deux informations de même nature : \zh{不但会说汉语，而且会说英语}. Le sujet \zh{他} est placé avant \zh{不但} et n'est pas répété après \zh{而且}. Le français dit « non seulement... mais encore... » : la logique est la même, mais l'ordre des mots diffère.

\textbf{Quatrième constat : la corrélation \zh{因为……所以……}.} La cause est introduite par \zh{因为}, la conséquence par \zh{所以}, et une virgule sépare les deux propositions : \zh{因为下雨了，所以我不去学校。} Grande différence avec le français : en français, on choisit « parce que » OU « donc », jamais les deux ensemble ; en chinois, les deux mots-outils se répondent et s'emploient volontiers dans la même phrase.

\begin{attention}[Le piège majeur des francophones]
En français, « parce qu'il pleut, donc je sors » est une faute. En chinois, c'est la construction normale : \zh{因为下雨了，所以我出去。} Ne supprimez jamais \zh{所以} par habitude du français. À l'inverse, en traduisant vers le français, il faudra choisir : soit « parce que... », soit « ...donc... ».
\end{attention}

\courssub{Tableau récapitulatif des trois structures observées}

\begin{center}
\begin{tabularx}{\linewidth}{|X|X|X|}
\hline
\rowcolor{popblueL} Structure & Exemple chinois + pinyin & Traduction française \\
\hline
Sujet + 怎么了？ & \zh{你怎么了？} Nǐ zěnme le ? & Qu'est-ce que tu as ? \\
\hline
Sujet + 怎么 + verbe ? & \zh{这个字怎么写？} Zhège zì zěnme xiě ? & Comment s'écrit ce caractère ? \\
\hline
Sujet + 怎么 + ne pas + verbe ? & \zh{你怎么不去上课？} Nǐ zěnme bú qù shàngkè ? & Pourquoi ne vas-tu pas en cours ? \\
\hline
Sujet + 不但 + A，而且 + B & \zh{他不但会说汉语，而且会说英语。} Tā búdàn huì shuō Hànyǔ, érqiě huì shuō Yīngyǔ. & Non seulement il parle chinois, mais encore il parle anglais. \\
\hline
因为 + cause，所以 + conséquence & \zh{因为下雨了，所以我不去学校。} Yīnwèi xià yǔ le, suǒyǐ wǒ bú qù xuéxiào. & Parce qu'il pleut, je ne vais pas à l'école. \\
\hline
\end{tabularx}
\end{center}

\begin{popfigure}
\begin{tikzpicture}[
  grow cyclic,
  mindmap,
  concept color=popblue,
  root concept/.append style={font=\large\bfseries, text=white},
  level 1 concept/.append style={sibling angle=120, font=\small, text=white, minimum size=3.2cm},
  every node/.style={concept}
]
\node[root concept]{3 structures 三种结构}
  child[concept color=poporange] { node[align=center]{Questionner\\ \zh{怎么了} / \zh{怎么}\\ \emph{zěnme}} }
  child[concept color=popgreen] { node[align=center]{Ajouter\\ \zh{不但……而且……}\\ \emph{búdàn... érqiě...}} }
  child[concept color=poppurple] { node[align=center]{Expliquer\\ \zh{因为……所以……}\\ \emph{yīnwèi... suǒyǐ...}} };
\end{tikzpicture}
\legende{Carte mentale : trois structures, trois fonctions de communication — questionner, ajouter, expliquer.}
\end{popfigure}

\courssub{Texte de découverte : mini-dialogue à l'hôpital}

Lisez ce dialogue deux fois : une première fois pour le sens global, une deuxième fois en soulignant les trois structures étudiées.

\begin{textebox}[À l'hôpital — 在医院]
A : 你怎么了？\\
\emph{A : Nǐ zěnme le ?}\\
« A : Qu'est-ce que tu as ? »\\[4pt]
B : 我头疼，而且发烧。\\
\emph{B : Wǒ tóu téng, érqiě fāshāo.}\\
« B : J'ai mal à la tête, et en plus j'ai de la fièvre. »\\[4pt]
A : 因为天气变了，所以很多人感冒了。你不但要多喝水，而且要好好休息。\\
\emph{A : Yīnwèi tiānqì biàn le, suǒyǐ hěn duō rén gǎnmào le. Nǐ búdàn yào duō hē shuǐ, érqiě yào hǎohao xiūxi.}\\
« A : Parce que le temps a changé, beaucoup de gens ont attrapé un rhume. Non seulement tu dois boire beaucoup d'eau, mais encore tu dois bien te reposer. »\\[4pt]
B : 好的，谢谢医生！\\
\emph{B : Hǎo de, xièxie yīshēng !}\\
« B : D'accord, merci docteur ! »
\end{textebox}

\textbf{Glossaire du dialogue.}

\begin{center}
\begin{tabularx}{\linewidth}{|X|X|X|X|}
\hline
\rowcolor{popblueL} Caractères & Pinyin & Nature & Traduction \\
\hline
头疼 & tóu téng & expression verbale & avoir mal à la tête \\
\hline
发烧 & fāshāo & verbe & avoir de la fièvre \\
\hline
天气 & tiānqì & nom & le temps (météo) \\
\hline
变 & biàn & verbe & changer \\
\hline
感冒 & gǎnmào & verbe / nom & attraper un rhume ; un rhume \\
\hline
喝水 & hē shuǐ & expression verbale & boire de l'eau \\
\hline
休息 & xiūxi & verbe & se reposer \\
\hline
\end{tabularx}
\end{center}

\textbf{Questions de compréhension.} 1) Quels sont les deux symptômes du patient ? 2) Pourquoi beaucoup de gens sont-ils malades ? 3) Quels sont les deux conseils du médecin ? 4) Retrouvez dans le dialogue les trois structures de la leçon et recopiez-les.

\begin{savaistu}[怎么, le champion des mots interrogatifs]
\zh{怎么} \emph{zěnme} est l'un des mots interrogatifs les plus fréquents du chinois parlé. On le retrouve partout : \zh{怎么办？} \emph{Zěnme bàn ?} « Que faire ? », expression de politesse et d'entraide qu'on entend dès qu'un problème surgit ; \zh{怎么样？} \emph{Zěnme yàng ?} « Comment ça va ? Qu'en penses-tu ? » ; \zh{怎么走？} \emph{Zěnme zǒu ?} « Comment y aller ? », la question du voyageur perdu. Maîtriser \zh{怎么}, c'est débloquer des dizaines de questions du quotidien, à Yaoundé comme à Pékin.
\end{savaistu}

\courssub{Premier exercice résolu : identifier la structure}

\begin{exoresolu}[Identifier la structure et sa fonction]
Énoncé. Pour chaque phrase, indiquez la structure utilisée et sa fonction (questionner l'état, questionner la manière, ajouter, expliquer).\\
a) \zh{老师今天怎么了？} \emph{Lǎoshī jīntiān zěnme le ?}\\
b) \zh{因为明天有考试，所以我今天晚上学习。} \emph{Yīnwèi míngtiān yǒu kǎoshì, suǒyǐ wǒ jīntiān wǎnshang xuéxí.}\\
c) \zh{杜阿拉不但很大，而且很热闹。} \emph{Dù'ālā búdàn hěn dà, érqiě hěn rènao.}\\
d) \zh{这个词怎么读？} \emph{Zhège cí zěnme dú ?}
\tcblower
Solution détaillée.\\
a) Structure : Sujet + 怎么了. Fonction : questionner l'état de quelqu'un. Traduction : « Qu'a le professeur aujourd'hui ? »\\
b) Structure : 因为 + cause，所以 + conséquence. Fonction : expliquer. Traduction : « Parce qu'il y a un examen demain, j'étudie ce soir. » On remarque la place du complément de temps \zh{今天晚上} entre le sujet et le verbe.\\
c) Structure : Sujet + 不但 + A，而且 + B. Fonction : ajouter une qualité à une autre. Traduction : « Douala est non seulement très grande, mais encore très animée. » Ici A et B sont des adjectifs (\zh{很大}, \zh{很热闹}) : les deux moitiés doivent être de même nature.\\
d) Structure : Sujet + 怎么 + verbe. Fonction : questionner la manière. Traduction : « Comment se lit ce mot ? »
\end{exoresolu}

\begin{attention}[Trois pièges repérés dès l'observation]
\begin{itemize}
\item Ne confondez pas \zh{怎么了} (état, « qu'est-ce qui ne va pas ? ») et \zh{怎么} + verbe (manière, « comment ? »). La particule finale \zh{了} fait toute la différence.
\item Dans \zh{不但……而且……}, si le sujet est le même pour les deux actions, il se place avant \zh{不但} et ne se répète pas : on dit \zh{他不但会说汉语，而且会说英语}, jamais \zh{他不但会说汉语，而且他会说英语} dans le style soigné.
\item L'ordre chinois est fixe : la cause (\zh{因为}) d'abord, la conséquence (\zh{所以}) ensuite. On ne peut pas inverser comme en français (« je ne vais pas à l'école parce qu'il pleut ») sans changer de structure.
\end{itemize}
\end{attention}

\begin{exercice}[À vous d'observer]
Soulignez les mots-outils dans les phrases suivantes, puis traduisez-les en français :\\
1. \zh{你怎么不买这个？}\\
2. \zh{我的同学不但喜欢足球，而且喜欢篮球。}\\
3. \zh{因为他生病了，所以今天没来上课。}\\
4. \zh{喀麦隆菜怎么做？}
\end{exercice}

\begin{corrige}[Exercice « À vous d'observer »]
1. Mots-outils : \zh{怎么} + négation \zh{不}. « Pourquoi n'achètes-tu pas ceci ? »\\
2. Mots-outils : \zh{不但……而且……}. « Mon camarade de classe aime non seulement le football, mais encore le basket-ball. »\\
3. Mots-outils : \zh{因为……所以……}. « Parce qu'il est malade, il n'est pas venu en cours aujourd'hui. »\\
4. Mot-outil : \zh{怎么} + verbe \zh{做}. « Comment prépare-t-on la cuisine camerounaise ? »
\end{corrige}

\coursec{Leçon 13 — Grammaire : Interrogation avec 怎么了 et 怎么 ; 不但……而且…… ; 因为……所以……}

\coursec{Observation et découverte}

\begin{textebox}[À la clinique universitaire de Yaoundé — Dialogue authentique]
A : 李明，你怎么了？你的脸色不太好。\\
\emph{Lǐ Míng, nǐ zěnme le? Nǐ de liǎnsè bù tài hǎo.}\\
« Li Ming, qu'as-tu ? Tu n'as pas bonne mine. »

B : 我头疼，可能感冒了。\\
\emph{Wǒ tóu téng, kěnéng gǎnmào le.}\\
« J'ai mal à la tête, j'ai peut-être attrapé un rhume. »

A : 医院怎么走？我陪你去吧。\\
\emph{Yīyuàn zěnme zǒu? Wǒ péi nǐ qù ba.}\\
« Comment va-t-on à l'hôpital ? Je t'accompagne. »

B : 坐出租车去吧。对了，你昨天怎么没来上课？\\
\emph{Zuò chūzūchē qù ba. Duì le, nǐ zuótiān zěnme méi lái shàngkè?}\\
« Prenons un taxi. Au fait, pourquoi n'es-tu pas venu en cours hier ? »

A : 因为我妈妈生病了，所以我在家照顾她。\\
\emph{Yīnwèi wǒ māma shēngbìng le, suǒyǐ wǒ zài jiā zhàogù tā.}\\
« Parce que ma mère est tombée malade, je suis resté à la maison pour m'occuper d'elle. »

B : 不但妈妈生病，而且妹妹也发烧了。\\
\emph{Búdàn māma shēngbìng, érqiě mèimei yě fāshāo le.}\\
« Non seulement maman est malade, mais en plus ma petite sœur a de la fièvre. »
\end{textebox}

Observons ce dialogue : trois structures grammaticales clés y apparaissent naturellement.
\begin{itemize}
\item L'\textbf{interrogation} avec \zh{怎么} sous trois formes : \zh{你怎么了？} (état), \zh{医院怎么走？} (manière), \zh{你怎么没来上课？} (raison).
\item La \textbf{corrélation additive} \zh{不但……而且……} « non seulement… mais encore… » (dernière réplique de B).
\item La \textbf{corrélation causale} \zh{因为……所以……} « parce que… donc… » (réplique d'A).
\end{itemize}
Cette leçon les analyse l'une après l'autre pour en maîtriser la construction, les nuances et les pièges fréquents pour un francophone.

\coursec{Interroger avec 怎么了 et 怎么}

\courssub{Première observation : trois questions, trois sens}

Dans le dialogue, \zh{怎么} apparaît trois fois avec trois fonctions distinctes :

\begin{center}
\begin{tabularx}{\linewidth}{|X|X|X|}
\hline
\rowcolor{popblueL} Question & Structure & Sens demandé \\
\hline
\zh{你怎么了？} \emph{Nǐ zěnme le?} & Sujet + \zh{怎么了} & l'état, le problème (« Qu'as-tu ? ») \\
\hline
\zh{医院怎么走？} \emph{Yīyuàn zěnme zǒu?} & \zh{怎么} + verbe & la manière, le moyen (« Comment ? ») \\
\hline
\zh{你怎么没来上课？} \emph{Nǐ zěnme méi lái shàngkè?} & \zh{怎么} + \zh{没} + verbe & la raison (« Pourquoi ? ») \\
\hline
\end{tabularx}
\end{center}

C'est toute la logique de cette section : la place de \zh{怎么} et les mots qui l'entourent déterminent le sens de la question. Étudions chaque schéma en détail.

\courssub{Schéma 1 : Sujet + 怎么了 ? — demander ce qui ne va pas}

\begin{propriete}[La structure 怎么了]
\textbf{Schéma :} Sujet + \zh{怎么} + \zh{了} ?

\begin{itemize}
\item Sans sujet : \zh{怎么了？} \emph{Zěnme le?} = « Qu'est-ce qui se passe ? Que se passe-t-il ? »
\item Avec sujet : \zh{你怎么了？} \emph{Nǐ zěnme le?} = « Qu'as-tu ? Qu'est-ce qui t'arrive ? »
\item La particule \zh{了} est \textbf{obligatoire} : elle marque le \cle{changement d'état} (quelque chose a changé par rapport à la normale : la personne est pâle, triste, absente, en retard…). Elle ne se traduit pas en français.
\item On n'emploie cette question que lorsqu'on \textbf{constate} un problème visible. On ne la dit jamais « pour saluer ».
\end{itemize}
\end{propriete}

\begin{popfigure}
\begin{tikzpicture}[
  bloc/.style={rectangle, rounded corners=4pt, draw=popblue, fill=popblueL, align=center, minimum height=0.9cm, font=\small},
  fleche/.style={-{Stealth[length=2.5mm]}, thick, popdark}
]
\node[bloc, align=center] (s) at (0,0){Sujet\\ \zh{你 / 他 / 老师}};
\node[bloc, fill=poporangeL, draw=poporange, align=center] (z) at (3.2,0){\zh{怎么}\\ comment / quoi};
\node[bloc, fill=poppinkL, draw=poppink, align=center] (l) at (6.2,0){\zh{了}\\ particule};
\node[bloc, fill=popgreenL, draw=popgreen] (q) at (8.6,0) {\zh{？}};
\draw[fleche] (s) -- (z);
\draw[fleche] (z) -- (l);
\draw[fleche] (l) -- (q);
\draw[fleche, poppink] (l.south) to[bend left=25] node[below, font=\footnotesize, align=center, text=popdark]{\zh{了} = changement d'état\\ (obligatoire, ne se traduit pas)} (z.south);
\end{tikzpicture}
\legende{Schéma de la structure Sujet + 怎么了 ? : le 了 final signale qu'un changement est survenu.}
\end{popfigure}

\begin{exemplebox}[怎么了 en situation]
\zh{你怎么了？你看起来很累。} \emph{Nǐ zěnme le? Nǐ kànqǐlai hěn lèi.} « Qu'as-tu ? Tu as l'air très fatigué. »

\zh{小王怎么了？他今天没来学校。} \emph{Xiǎo Wáng zěnme le? Tā jīntiān méi lái xuéxiào.} « Qu'a Xiao Wang ? Il n'est pas venu à l'école aujourd'hui. »

\zh{我的手机怎么了？打不开了。} \emph{Wǒ de shǒujī zěnme le? Dǎ bu kāi le.} « Qu'a mon téléphone ? Il ne s'allume plus. »

\zh{外面怎么了？怎么这么吵？} \emph{Wàimiàn zěnme le? Zěnme zhème chǎo?} « Qu'est-ce qui se passe dehors ? Pourquoi tant de bruit ? »
\end{exemplebox}

Remarquez le troisième exemple : le sujet de \zh{怎么了} peut être une \textbf{chose} (un téléphone, une voiture, le marché). La question signifie alors « Qu'est-ce qui ne va pas avec X ? ». C'est très utile dans la vie quotidienne à Douala comme à Pékin.

\begin{attention}[你怎么了 n'est pas une salutation !]
Erreur classique du francophone : vouloir traduire « Comment vas-tu ? » par \zh{你怎么了？}. C'est faux ! \zh{你怎么了？} suppose un problème visible (tu es pâle, tu boites, tu pleures) et signifie « Qu'as-tu ? Qu'est-ce qui t'arrive ? ». Pour saluer, la formule correcte est \zh{你好吗？} \emph{Nǐ hǎo ma?} ou, plus naturel, \zh{你怎么样？} \emph{Nǐ zěnmeyàng?} « Comment ça va ? ». Dire \zh{你怎么了？} à quelqu'un qui va bien revient à lui dire « Tu as l'air malade, qu'est-ce qui ne va pas ? » — très maladroit !
\end{attention}

\courssub{Schéma 2 : 怎么 + Verbe ? — demander la manière}

\begin{propriete}[怎么 + verbe = « comment ? »]
\textbf{Schéma :} (Sujet) + \zh{怎么} + Verbe (+ complément) ?

Cette structure demande la \cle{manière}, le \cle{moyen}, la \cle{procédure} : « comment faire ? ». \zh{怎么} se place toujours \textbf{immédiatement avant le verbe}, jamais en tête de phrase comme l'anglais \emph{how} ni en fin de phrase comme le français « on y va comment ? ».
\end{propriete}

\begin{exemplebox}[Demander la manière]
\zh{这个汉字怎么写？} \emph{Zhège Hànzì zěnme xiě?} « Comment s'écrit ce caractère ? »

\zh{汉语怎么说 « santé » ？} \emph{Hànyǔ zěnme shuō « jiànkāng »?} « Comment dit-on "santé" en chinois ? »

\zh{从杜阿拉怎么去雅温得？} \emph{Cóng Dù'ālā zěnme qù Yǎwēndé?} « Comment va-t-on de Douala à Yaoundé ? »

\zh{这道菜怎么做？} \emph{Zhè dào cài zěnme zuò?} « Comment prépare-t-on ce plat ? »

\zh{这个词怎么读？} \emph{Zhège cí zěnme dú?} « Comment se prononce ce mot ? »
\end{exemplebox}

Notez le modèle très fréquent \zh{……怎么说？} « Comment dit-on … ? », indispensable en classe de langue, et \zh{……怎么走？} « Comment va-t-on à … ? », indispensable pour demander son chemin.

\courssub{Schéma 3 : Sujet + 怎么 + 不/没 + Verbe ? — demander la raison}

\begin{propriete}[怎么 + 不/没 + verbe = « pourquoi ? »]
\textbf{Schéma :} Sujet + \zh{怎么} + \zh{不} / \zh{没} + Verbe ?

Cette structure demande la \cle{raison} d'une non-action, souvent avec un ton de \textbf{surprise} ou de \textbf{reproche doux} :

\begin{itemize}
\item \zh{怎么} + \zh{不} + verbe : pourquoi ne fais-tu pas X (présent, futur, général) ? \zh{你怎么不去？} \emph{Nǐ zěnme bù qù?} « Pourquoi n'y vas-tu pas ? »
\item \zh{怎么} + \zh{没} + verbe : pourquoi n'as-tu pas fait X (passé) ? \zh{你怎么没来？} \emph{Nǐ zěnme méi lái?} « Pourquoi n'es-tu pas venu ? »
\end{itemize}

Rappel de grammaire : on nie le passé avec \zh{没}(有), jamais avec \zh{不}.
\end{propriete}

\begin{exemplebox}[Demander la raison avec 怎么]
\zh{他怎么哭了？} \emph{Tā zěnme kū le?} « Pourquoi pleure-t-il ? » (surprise : il pleure, ce n'est pas normal)

\zh{你怎么不吃饭？} \emph{Nǐ zěnme bù chīfàn?} « Pourquoi ne manges-tu pas ? »

\zh{昨天的比赛，你怎么没来看？} \emph{Zuótiān de bǐsài, nǐ zěnme méi lái kàn?} « Pourquoi n'es-tu pas venu voir le match d'hier ? »

\zh{老师怎么还没来？} \emph{Lǎoshī zěnme hái méi lái?} « Comment se fait-il que le professeur ne soit pas encore arrivé ? »
\end{exemplebox}

\begin{aretenir}[怎么 (pourquoi) ou 为什么 (pourquoi) ?]
Les deux demandent une raison, mais la nuance est différente :

\begin{itemize}
\item \zh{为什么} : question \textbf{neutre}, objective, souvent écrite ou formelle. \zh{你为什么学汉语？} \emph{Nǐ wèishénme xué Hànyǔ?} « Pourquoi apprends-tu le chinois ? »
\item \zh{怎么} (+ \zh{不}/\zh{没} ou + \zh{了}) : question \textbf{chargée de surprise, d'inquiétude ou de reproche léger}. \zh{你怎么又迟到了？} \emph{Nǐ zěnme yòu chídào le?} « Pourquoi es-tu encore en retard ? » (le ton est celui d'un reproche).
\end{itemize}

En résumé : \zh{为什么} interroge la cause ; \zh{怎么} exprime en plus l'étonnement de celui qui parle.
\end{aretenir}

\begin{popfigure}
\begin{tikzpicture}[
  decision/.style={rectangle, rounded corners=4pt, draw=popblue, fill=popblueL, align=center, font=\small},
  result/.style={rectangle, rounded corners=4pt, draw=popgreen, fill=popgreenL, align=center, font=\small},
  fleche/.style={-{Stealth[length=2.5mm]}, thick, popdark}
]
\node[decision, fill=poporangeL, draw=poporange] (racine) at (0,0) {\zh{怎么} + … ?};
\node[decision, align=center] (m1) at (-4.2,-1.8){\zh{怎么} + verbe\\ \zh{怎么走？}};
\node[decision, align=center] (m2) at (0,-1.8){\zh{怎么} + \zh{不}/\zh{没} + verbe\\ \zh{怎么没来？}};
\node[decision, align=center] (m3) at (4.2,-1.8){\zh{怎么} + \zh{了}\\ \zh{怎么了？}};
\node[result, align=center] (r1) at (-4.2,-3.4){manière\\ « comment ? »};
\node[result, align=center] (r2) at (0,-3.4){raison\\ « pourquoi ? »};
\node[result, align=center] (r3) at (4.2,-3.4){état, problème\\ « qu'est-ce qui se passe ? »};
\draw[fleche] (racine) -- (m1);
\draw[fleche] (racine) -- (m2);
\draw[fleche] (racine) -- (m3);
\draw[fleche] (m1) -- (r1);
\draw[fleche] (m2) -- (r2);
\draw[fleche] (m3) -- (r3);
\end{tikzpicture}
\legende{Arbre de décision : la forme qui suit 怎么 détermine le sens de la question.}
\end{popfigure}

\courssub{Répondre aux questions en 怎么了 et 怎么}

\begin{exemplebox}[Réponses types]
À \zh{你怎么了？} on répond en décrivant l'état ou le problème :

\zh{我感冒了。} \emph{Wǒ gǎnmào le.} « J'ai attrapé un rhume. »

\zh{我有点头疼。} \emph{Wǒ yǒudiǎn tóu téng.} « J'ai un peu mal à la tête. »

\zh{没什么。} \emph{Méi shénme.} « Rien. / Ce n'est rien. » (réponse très fréquente pour rassurer)

À \zh{你怎么没来？} on répond avec la raison, souvent introduite par \zh{因为} :

\zh{因为我生病了，所以没来上课。} \emph{Yīnwèi wǒ shēngbìng le, suǒyǐ méi lái shàngkè.} « Parce que j'étais malade, je ne suis pas venu en cours. »

À \zh{医院怎么走？} on répond en indiquant le chemin :

\zh{一直走，然后往左拐。} \emph{Yìzhí zǒu, ránhòu wǎng zuǒ guǎi.} « Allez tout droit, puis tournez à gauche. »
\end{exemplebox}

\begin{attention}[Ne pas confondre 怎么 et 什么]
Les deux mots commencent par \zh{怎}/\zh{什} et se ressemblent à l'oreille, mais :

\begin{itemize}
\item \zh{什么} \emph{shénme} = « quoi, quel » : \zh{你说什么？} \emph{Nǐ shuō shénme?} « Qu'est-ce que tu dis ? »
\item \zh{怎么} \emph{zěnme} = « comment, pourquoi » : \zh{你怎么说？} \emph{Nǐ zěnme shuō?} « Comment le dis-tu ? / Pourquoi dis-tu cela ? »
\end{itemize}

Comparez aussi : \zh{这是什么？} « Qu'est-ce que c'est ? » $\neq$ \zh{这怎么用？} « Comment utilise-t-on ceci ? ». Une seule syllabe change tout le sens !
\end{attention}

\courssub{Tableau récapitulatif et comparaison avec le français}

\begin{center}
\begin{tabularx}{\linewidth}{|X|X|X|}
\hline
\rowcolor{popblueL} Structure & Exemple (caractères + pinyin) & Traduction \\
\hline
\zh{怎么了} (sans sujet) & \zh{怎么了？} \emph{Zěnme le?} & Qu'est-ce qui se passe ? \\
\hline
Sujet + \zh{怎么了} & \zh{你怎么了？} \emph{Nǐ zěnme le?} & Qu'as-tu ? \\
\hline
\zh{怎么} + verbe & \zh{怎么去学校？} \emph{Zěnme qù xuéxiào?} & Comment aller à l'école ? \\
\hline
\zh{怎么} + \zh{不} + verbe & \zh{你怎么不吃？} \emph{Nǐ zěnme bù chī?} & Pourquoi ne manges-tu pas ? \\
\hline
\zh{怎么} + \zh{没} + verbe & \zh{你怎么没来？} \emph{Nǐ zěnme méi lái?} & Pourquoi n'es-tu pas venu ? \\
\hline
\zh{为什么} + proposition & \zh{你为什么学汉语？} \emph{Nǐ wèishénme xué Hànyǔ?} & Pourquoi apprends-tu le chinois ? (neutre) \\
\hline
\end{tabularx}
\end{center}

Comparaison avec le français : le français place le mot interrogatif en tête (« Comment vas-tu à l'école ? », « Pourquoi pleures-tu ? »), parfois avec inversion du sujet. Le chinois, lui, garde \textbf{l'ordre normal de la phrase affirmative} et insère simplement \zh{怎么} à la place du mot demandé : \zh{你走路去学校。} « Tu vas à l'école à pied » devient \zh{你怎么去学校？} « Comment vas-tu à l'école ? ». Aucune inversion, aucun auxiliaire : c'est plus simple qu'en français, à condition de placer \zh{怎么} juste avant le verbe.

\begin{exoresolu}[Traduire et justifier]
Énoncé : traduisez en chinois les trois questions suivantes, en choisissant la bonne structure, puis justifiez votre choix. 1) Qu'as-tu ? Tu as l'air triste. 2) Comment dit-on « ndolé » en chinois ? 3) Pourquoi n'es-tu pas venu au match de football hier ?
\tcblower
Solution détaillée :

1) \zh{你怎么了？你看起来很难过。} \emph{Nǐ zěnme le? Nǐ kànqǐlai hěn nánguò.} — On constate un changement d'état (tristesse visible) : structure Sujet + \zh{怎么了}, avec le \zh{了} final obligatoire.

2) \zh{« Ndolé » 汉语怎么说？} \emph{« Ndolé » Hànyǔ zěnme shuō?} — On demande la manière de dire un mot : structure \zh{怎么} + verbe (\zh{说}), modèle \zh{……怎么说？}.

3) \zh{昨天的足球比赛，你怎么没来？} \emph{Zuótiān de zúqiú bǐsài, nǐ zěnme méi lái?} — On demande la raison d'une non-action passée : structure \zh{怎么} + \zh{没} + verbe. On choisit \zh{没} (et non \zh{不}) parce que l'action passée niée se nie toujours avec \zh{没}. Le ton exprime une légère surprise ; pour une question purement neutre, on aurait pu dire \zh{你为什么没来？}.
\end{exoresolu}

\begin{exercice}[Entraînement : Interrogation avec 怎么了 et 怎么]
1) Complétez avec \zh{怎么}, \zh{怎么了} ou \zh{什么} :
a) 你 \_\_\_\_ ？你的脸很红。
b) 这个字 \_\_\_\_ 写？
c) 你说 \_\_\_\_ ？我没听清楚。
d) 他昨天 \_\_\_\_ 没来上课？

2) Traduisez en français :
a) \zh{你怎么了？}
b) \zh{从你家到学校怎么走？}
c) \zh{她怎么哭了？}

3) Traduisez en chinois :
a) Qu'est-ce qui se passe ?
b) Comment va-t-on au marché de Mokolo ?
c) Pourquoi n'as-tu pas fait tes devoirs ?
\end{exercice}

\begin{corrige}[Exercice « Entraînement : Interrogation avec 怎么了 et 怎么 »]
1) a) \zh{你怎么了？} (état visible : le visage rouge) ; b) \zh{这个字怎么写？} (manière) ; c) \zh{你说什么？} (« quoi », pas « comment ») ; d) \zh{他昨天怎么没来上课？} (raison, passé nié avec \zh{没}).

2) a) « Qu'as-tu ? » ; b) « Comment va-t-on de chez toi à l'école ? » ; c) « Pourquoi pleure-t-elle ? »

3) a) \zh{怎么了？} \emph{Zěnme le?} ; b) \zh{Mokolo 市场怎么走？} \emph{Mokolo shìchǎng zěnme zǒu?} ; c) \zh{你怎么没做作业？} \emph{Nǐ zěnme méi zuò zuòyè?}
\end{corrige}

\begin{savaistu}[怎么 dans la culture chinoise]
En Chine, demander \zh{你怎么了？} à quelqu'un qui semble souffrir est un signe de sollicitude très apprécié : on montre qu'on a remarqué l'autre. Au Cameroun, la même attention existe — on demande volontiers « Ça va ? Tu as quoi ? » à un voisin ou un camarade. Dans les deux cultures, la réponse polie \zh{没什么} \emph{méi shénme} (« ce n'est rien ») permet de ne pas importuner son interlocuteur, exactement comme notre « ça va, ne t'inquiète pas ».
\end{savaistu}

\coursec{La corrélation 不但……而且……}

\courssub{Structure et sens}

\begin{propriete}[不但……而且…… — « non seulement… mais encore… »]
\textbf{Schéma :} Sujet + \zh{不但} + Proposition 1 ， \zh{而且} + Proposition 2 。

\begin{itemize}
\item \zh{不但} \emph{búdàn} = « non seulement » ; \zh{而且} \emph{érqiě} = « mais encore / de plus / en plus ».
\item La structure lie deux propositions qui vont \textbf{dans le même sens} (addition d'arguments, cumul de faits). La seconde proposition renforce ou étend la première.
\item \zh{而且} est \textbf{obligatoire} en réponse à \zh{不但} (on ne dit pas *\zh{不但……也……}).
\item Ponctuation : une virgule chinoise \zh{，} sépare généralement les deux propositions.
\end{itemize}
\end{propriete}

\begin{exemplebox}[不但……而且…… en contexte]
\zh{这道菜不但好吃，而且不贵。} \emph{Zhè dào cài búdàn hǎochī, érqiě bù guì.} « Ce plat est non seulement bon, mais en plus il n'est pas cher. »

\zh{他不但会说汉语，而且会写汉字。} \emph{Tā búdàn huì shuō Hànyǔ, érqiě huì xiě Hànzì.} « Il ne sait pas seulement parler chinois, mais il sait aussi écrire les caractères. »

\zh{昨天不但下雨，而且刮大风。} \emph{Zuótiān búdàn xiàyǔ, érqiě guā dàfēng.} « Hier, non seulement il a plu, mais en plus il a fait un grand vent. »

\zh{喀麦隆不但有足球，而且有美食。} \emph{Kāmàilóng búdàn yǒu zúqiú, érqiě yǒu měishí.} « Le Cameroun a non seulement le football, mais aussi une gastronomie. »
\end{exemplebox}

\courssub{Cas particulier : sujets différents}

Si les deux propositions ont des \textbf{sujets différents}, \zh{不但} se place \textbf{devant le premier sujet}, et \zh{而且} devant le second.

\begin{exemplebox}[Sujets différents]
\zh{不但老师来了，而且学生也来了。} \emph{Búdàn lǎoshī lái le, érqiě xuésheng yě lái le.} « Non seulement le prof est venu, mais les élèves sont venus aussi. »

\zh{不但哥哥喜欢看书，而且妹妹也喜欢。} \emph{Búdàn gēge xǐhuan kànshū, érqiě mèimei yě xǐhuan.} « Non seulement le grand frère aime lire, mais la petite sœur aussi. »
\end{exemplebox}

\begin{attention}[Piège : l'ordre des mots et le "aussi"]
En français, « aussi » / « également » se place souvent après le sujet ou en fin de phrase. En chinois, \zh{也} \emph{yě} (aussi) ou \zh{还} \emph{hái} (encore) se place \textbf{avant le verbe} dans la seconde proposition.
\begin{itemize}
\item Correct : \zh{他不但聪明，而且很用功。} (Sujet unique, pas de \zh{也} nécessaire).
\item Correct : \zh{不但他聪明，而且他弟弟也聪明。} (\zh{也} avant l'adjectif \zh{聪明}).
\item Incorrect : *\zh{而且他弟弟聪明也。}
\end{itemize}
\end{attention}

\courssub{Variante : 不但……还…… / 不但……也……}

À l'oral, \zh{而且} peut être remplacé par \zh{还} \emph{hái} (« en plus, encore ») ou \zh{也} \emph{yě} (« aussi »), mais \zh{而且} reste la forme standard à l'écrit et aux examens.

\begin{exemplebox}[Variantes orales]
\zh{这件衣服不但漂亮，还很便宜。} \emph{Zhè jiàn yīfu búdàn piàoliang, hái hěn piányi.} « Ce vêtement est non seulement joli, mais en plus il est très bon marché. »

\zh{我不但去过北京，也去过上海。} \emph{Wǒ búdàn qùguo Běijīng, yě qùguo Shànghǎi.} « Je ne suis pas seulement allé à Pékin, mais aussi à Shanghai. »
\end{exemplebox}

\begin{exoresolu}[Transformation : relier deux phrases avec 不但……而且……]
Énoncé : Reliez les deux phrases par \zh{不但……而且……}. 1) 这本书很有意思。这本书很有用。 2) 我会说英语。我会说法语。
\tcblower
Solution détaillée :

1) \zh{这本书不但很有意思，而且很有用。} \emph{Zhè běn shū búdàn hěn yǒu yìsi, érqiě hěn yǒu yòng.} (Sujet unique \zh{这本书}, placé avant \zh{不但}).

2) \zh{我不但会说英语，而且会说法语。} \emph{Wǒ búdàn huì shuō Yīngyǔ, érqiě huì shuō Fǎyǔ.} (Sujet unique \zh{我}, placé avant \zh{不但}). Attention : on ne répète pas le sujet devant \zh{而且} s'il est identique.
\end{exoresolu}

\begin{exercice}[Entraînement : 不但……而且……]
1) Reliez par \zh{不但……而且……} :
a) 今天很热。今天很湿。
b) 我哥哥会做饭。我姐姐也会做饭。
c) 学汉语不难。学汉语很有用。

2) Traduisez en chinois :
a) Non seulement il pleut, mais il fait aussi froid.
b) Mon père non seulement conduit bien, mais en plus il répare les voitures.
c) Le ndolé n'est pas seulement délicieux, il est aussi nutritif.
\end{exercice}

\begin{corrige}[Exercice « Entraînement : 不但……而且…… »]
1) a) \zh{今天不但很热，而且很湿。} ; b) \zh{不但我哥哥会做饭，而且我姐姐也会做饭。} (sujets différents) ; c) \zh{学汉语不但不难，而且很有用。}

2) a) \zh{不但下雨，而且很冷。} (sujet unique « il » implicite / météo) ; b) \zh{我爸爸不但开车开得好，而且会修车。} ; c) \zh{ Ndolé 不但好吃，而且有营养。} \emph{Ndolé búdàn hǎochī, érqiě yǒu yíngyǎng.}
\end{corrige}

\coursec{La corrélation 因为……所以……}

\courssub{Structure et sens}

\begin{propriete}[因为……所以…… — « parce que… donc… »]
\textbf{Schéma :} \zh{因为} + Cause ， \zh{所以} + Conséquence 。

\begin{itemize}
\item \zh{因为} \emph{yīnwèi} introduit la \cle{cause} (la raison).
\item \zh{所以} \emph{suǒyǐ} introduit la \cle{conséquence} (le résultat).
\item Les deux conjonctions sont \textbf{corrélées} : \zh{因为} appelle \zh{所以} (et inversement). On ne dit pas *\zh{因为……但是……} pour une relation cause-effet simple.
\item Ordre immuable : \textbf{Cause d'abord, Conséquence ensuite}. (Le français permet « Il est rentré parce qu'il pleuvait » = Conséquence + Cause ; le chinois impose \zh{因为下雨了，所以他回家了。})
\end{itemize}
\end{propriete}

\begin{popfigure}
\begin{tikzpicture}[
  bloc/.style={rectangle, rounded corners=4pt, draw=popblue, fill=popblueL, align=center, minimum height=0.8cm, font=\small},
  fleche/.style={-{Stealth[length=2.5mm]}, thick, popdark}
]
\node[bloc, fill=poporangeL, draw=poporange] (yin) at (0,0) {\zh{因为} + Cause};
\node[bloc, fill=popgreenL, draw=popgreen] (suo) at (5.5,0) {\zh{所以} + Conséquence};
\draw[fleche] (yin) -- node[above, font=\footnotesize] {entraîne} (suo);
\node[font=\footnotesize, text=popmuted] at (2.75, -1) {Ordre fixe : Cause $\rightarrow$ Conséquence};
\end{tikzpicture}
\legende{Schéma de la corrélation 因为……所以……}
\end{popfigure}

\begin{exemplebox}[因为……所以…… en situation]
\zh{因为下雨了，所以比赛取消了。} \emph{Yīnwèi xiàyǔ le, suǒyǐ bǐsài qǔxiāo le.} « Parce qu'il a plu, le match a été annulé. »

\zh{因为我喜欢足球，所以我每天看比赛。} \emph{Yīnwèi wǒ xǐhuan zúqiú, suǒyǐ wǒ měitiān kàn bǐsài.} « Parce que j'aime le foot, je regarde les matches tous les jours. »

\zh{因为生病了，所以没去上学。} \emph{Yīnwèi shēngbìng le, suǒyǐ méi qù shàngxué.} « Parce que j'étais malade, je ne suis pas allé à l'école. » (Sujet omis car contexte clair)
\end{exemplebox}

\courssub{Gestion du sujet : un ou deux sujets ?}

\begin{propriete}[Place du sujet avec 因为……所以……]
Deux cas de figure :
\begin{enumerate}
\item \textbf{Sujet unique (identique dans les deux propositions) :} Le sujet se place \textbf{avant 因为} (sujet commun) \textbf{OU} se répète devant \zh{因为} et \zh{所以} (insistance). Il ne se place \textbf{jamais} entre \zh{因为} et le verbe de la cause.
\item \textbf{Sujets différents :} Chaque sujet se place \textbf{après} sa conjonction respective (\zh{因为} Sujet₁ ..., \zh{所以} Sujet₂ ...).
\end{enumerate}
\end{propriete}

\begin{exemplebox}[Position du sujet]
\textbf{Sujet unique (avant 因为) :}
\zh{我因为累了，所以早早睡了。} \emph{Wǒ yīnwèi lèi le, suǒyǐ zǎozǎo shuì le.} « Moi, parce que j'étais fatigué, je me suis couché très tôt. »

\textbf{Sujet unique (répété pour insister) :}
\zh{因为我累了，所以我早早睡了。} \emph{Yīnwèi wǒ lèi le, suǒyǐ wǒ zǎozǎo shuì le.} « Parce que MOI j'étais fatigué, MOI je me suis couché tôt. »

\textbf{Sujets différents :}
\zh{因为下雨了，所以路湿了。} \emph{Yīnwèi xiàyǔ le, suǒyǐ lù shī le.} « Parce qu'il a plu, la route est mouillée. » (Sujet 1 = ciel/il ; Sujet 2 = route).

\zh{因为妈妈生病了，所以我请假了。} \emph{Yīnwèi māma shēngbìng le, suǒyǐ wǒ qǐngjià le.} « Parce que maman est malade, j'ai pris un congé. »
\end{exemplebox}

\begin{attention}[Erreur majeure : inverser Cause et Conséquence]
En français, on dit souvent : « Je n'ai pas réussi \textbf{parce que} je n'ai pas révisé. » (Conséquence + \emph{parce que} + Cause).
En chinois, c'est \textbf{interdit} avec \zh{因为……所以……}. On \textbf{doit} dire :
\zh{因为我没复习，所以我没考好。} \emph{Yīnwèi wǒ méi fùxí, suǒyǐ wǒ méi kǎo hǎo.}
(Cause + \zh{所以} + Conséquence).
Si vous voulez mettre le résultat en premier, utilisez \zh{……是因为……} « … est parce que … » : \zh{我没考好是因为我没复习。}
\end{attention}

\courssub{Omission de 因为 ou 所以}

À l'oral, l'une des deux conjonctions peut être omise (ellipse), mais la relation logique reste implicite.

\begin{exemplebox}[Ellipses courantes]
\zh{他没来，是因为生病了。} \emph{Tā méi lái, shì yīnwèi shēngbìng le.} (Pas de \zh{所以}, structure \zh{是因为}).

\zh{因为太贵，所以不买了。} \emph{Yīnwèi tài guì, suǒyǐ bù mǎi le.} (Complet).

\zh{太贵了，不买了。} \emph{Tài guì le, bù mǎi le.} (Les deux omis, relation cause-effet évidente par l'intonation).
\end{exemplebox}

\begin{exoresolu}[Correction d'erreurs : 因为……所以……]
Énoncé : Corrigez les phrases fautives.
1) \zh{所以我迟到了，因为堵车。}
2) \zh{因为他聪明，所以他考试考得好，而且他很努力。} (Mélange structures)
\tcblower
Solution détaillée :

1) \textbf{Erreur :} Ordre inversé (Conséquence + Cause).
\textbf{Correction :} \zh{因为堵车，所以我迟到了。} \emph{Yīnwèi dǔchē, suǒyǐ wǒ chídào le.} (Cause : embouteillage $\rightarrow$ Conséquence : retard).

2) \textbf{Erreur :} \zh{而且} n'appartient pas à la corrélation \zh{因为……所以……}. La phrase mélange cause-effet et addition.
\textbf{Correction :} Deux phrases distinctes ou restructuration.
Option A : \zh{因为他聪明，所以他考试考得好。不但聪明，而且很努力。}
Option B : \zh{他不但聪明，而且很努力，所以考试考得好。} (Cause complexe = addition de qualités $\rightarrow$ Conséquence).
\end{exoresolu}

\begin{exercice}[Entraînement : 因为……所以……]
1) Reliez par \zh{因为……所以……} (attention à l'ordre Cause $\rightarrow$ Conséquence) :
a) Je n'ai pas mangé. Je n'ai pas faim.
b) Il a eu 20/20. Il a beaucoup révisé.
c) La route est glissante. Il a plu.

2) Traduisez en chinois :
a) Parce qu'il y a un embouteillage, je suis en retard.
b) Comme (parce que) tu es mon ami, je t'aide.
c) Maman est malade, donc je rentre à la maison.
\end{exercice}

\begin{corrige}[Exercice « Entraînement : 因为……所以…… »]
1) a) \zh{因为我不饿，所以我没吃饭。} ; b) \zh{因为他复习了很多，所以他考了二十分。} (Cause = révisé ; Conséquence = note) ; c) \zh{因为下雨了，所以路滑。}

2) a) \zh{因为堵车，所以我迟到了。} ; b) \zh{因为你是我的朋友，所以我帮你。} ; c) \zh{因为妈妈生病了，所以我回家。}
\end{corrige}

\coursec{Comparaison avec le français et pièges — Synthèse}

\courssub{Tableau comparatif des trois structures}

\begin{center}
\begin{tabularx}{\linewidth}{|X|X|X|X|}
\hline
\rowcolor{popblueL} Structure & Fonction & Ordre des mots (FR vs CN) & Piège principal francophone \\
\hline
\zh{怎么了} & Constat d'anomalie & FR : « Qu'as-tu ? » (inversion) & Confondre avec salutation \zh{你好吗？} \\
\hline
\zh{怎么} + V & Manière & FR : « Comment + V ? » (tête) & Mettre \zh{怎么} en tête de phrase (*\zh{怎么你去？}) \\
\hline
\zh{怎么} + \zh{不/没} + V & Raison (surprise) & FR : « Pourquoi + V ? » & Confondre \zh{不}/\zh{没} (aspect) ; confondre avec \zh{为什么} (neutre) \\
\hline
\zh{不但} A \zh{而且} B & Addition (A+B) & FR : « Non seulement A mais B » & Oublier \zh{而且} ; mal placer \zh{也}/\zh{还} \\
\hline
\zh{因为} Cause \zh{所以} Conséqu. & Causalité & FR : Cause $\rightarrow$ Conséqu. OU Conséqu. $\rightarrow$ Cause & Inverser l'ordre (*\zh{所以……因为……}) \\
\hline
\end{tabularx}
\end{center}

\courssub{Exercices de discrimination et transformation}

\begin{exercice}[Discrimination : Choisissez la bonne structure]
Complétez avec : \zh{怎么了} / \zh{怎么} / \zh{什么} / \zh{不但……而且……} / \zh{因为……所以……}
1) \_\_\_\_ 他生病了，\_\_\_\_ 他没来上课。
2) 你 \_\_\_\_ ？脸色苍白。
3) 这道题 \_\_\_\_ 做？太难了。
4) 他 \_\_\_\_ 来了，\_\_\_\_ 带了礼物。
5) \_\_\_\_ 你不吃饭？饿不饿？
\end{exercice}

\begin{corrige}[Exercice « Discrimination »]
1) \zh{因为} 他生病了，\zh{所以} 他没来上课。 (Causalité)
2) 你 \zh{怎么了} ？ (Constat état)
3) 这道题 \zh{怎么} 做？ (Manière)
4) 他 \zh{不但} 来了，\zh{而且} 带了礼物。 (Addition, sujet unique avant \zh{不但})
5) 你 \zh{怎么} 不吃饭？ (Raison non-action présente, surprise)
\end{corrige}

\begin{exercice}[Transformation guidée]
Transformez les phrases selon le modèle indiqué.
Modèle : \zh{他很高。他很壮。} $\rightarrow$ \zh{他不但很高，而且很壮。}

1) \zh{下雨了。比赛停了。} (用 \zh{因为……所以……})
2) \zh{因为太累，所以睡觉。} (用 \zh{是因为……} — mettre le résultat en premier)
3) \zh{你为什么不去？} (Remplacer \zh{为什么} par \zh{怎么} pour exprimer la surprise)
4) \zh{这个词怎么读？} (Répondre : « Je ne sais pas le lire. » \zh{不会读})
5) \zh{我不但喜欢喝茶，而且喜欢喝咖啡。} (Garder le sens, changer structure : \zh{除了……以外，还……} — révision module 1)
\end{exercice}

\begin{corrige}[Exercice « Transformation guidée »]
1) \zh{因为下雨了，所以比赛停了。}
2) \zh{睡觉是因为太累了。} (ou \zh{我睡觉是因为太累了。})
3) \zh{你怎么不去？}
4) \zh{这个词我不会读。} (ou \zh{这个词怎么读？} $\rightarrow$ \zh{我不会读这个词。})
5) \zh{除了喜欢喝茶以外，我还喜欢喝咖啡。}
\end{corrige}

\coursec{Expression écrite et orale — Mise en situation}

\courssub{Production écrite : Message d'absence / Petit mot}

\begin{methode}[Rédiger un message d'absence formel (SMS / WeChat / Mot sur le bureau)]
\textbf{Contexte :} Vous êtes élève à Yaoundé ou étudiant à Pékin. Vous ne pouvez pas venir en cours / au travail. Écrivez un message à votre professeur / chef.
\textbf{Étapes :}
\begin{enumerate}
\item \textbf{Salutation polie} : \zh{老师好 / 王经理好。}
\item \textbf{Constat / Annonce} : \zh{我今天不能来上课 / 上班了。}
\item \textbf{Justification (Cause)} : \zh{因为……所以……} (ex: \zh{因为发烧了，所以去医院。})
\item \textbf{Information complémentaire (Addition)} : \zh{不但……而且……} (ex: \zh{不但发烧，而且头疼。})
\item \textbf{Demande / Engagement} : \zh{请假一天。明天一定来。}
\item \textbf{Formule de politesse} : \zh{谢谢您。/ 麻烦您了。}
\end{enumerate}
\end{methode}

\begin{exemplebox}[Modèle de message — Niveau HSK 3/4]
\shortstack[l]{\zh{陈老师好：}\\ \zh{因为我昨晚发高烧，所以今天早上去医院挂号了。}\\ \zh{不但发烧，而且嗓子也很疼。}\\ \zh{医生说要休息两天。}\\ \zh{请假两天（10月15日--16日），作业我会补上。}\\ \zh{麻烦您了，谢谢！}\\ \zh{学生：李明}\\ \zh{2025年10月14日 早上 8:00}}
\emph{Chén lǎoshī hǎo: Yīnwèi wǒ zuówǎn fā gāoshāo, suǒyǐ jīntiān zǎoshang qù yīyuàn guàhào le. Búdàn fāshāo, érqiě sǎngzi yě hěn téng. Yīshēng shuō yào xiūxi liǎng tiān. Qǐngjià liǎng tiān (10 yuè 15 rì--16 rì), zuòyè wǒ huì bǔ shàng. Máfan nín le, xièxie! Xuéshēng: Lǐ Míng. 2025 nián 10 yuè 14 rì zǎoshang 8:00.}
\end{exemplebox}

\begin{exercice}[Production écrite : À vous d'écrire]
Vous êtes \zh{王伟} (Wáng Wěi), élève en Terminale au Lycée de Buea. Vous ne pouvez pas venir en cours demain car vous devez accompagner votre père à l'hôpital à Douala (transfert). Écrivez un message à votre professeur principal \zh{张老师} (Zhāng Lǎoshī) en utilisant \textbf{au moins une fois} \zh{因为……所以……} et \textbf{au moins une fois} \zh{不但……而且……}. (80-100 caractères).
\end{exercice}

\courssub{Production orale : Jeu de rôle « À la clinique / Au commissariat »}

\begin{experience}[Jeu de rôle à 3 : Le patient, L'infirmier, L'accompagnateur]
\textbf{Scénario :} Un élève (Patient) arrive à la clinique scolaire. Il a mal quelque part. L'Infirmier l'interroge. L'Accompagnateur aide à décrire les symptômes et demande le chemin pour la pharmacie.
\textbf{Contraintes linguistiques obligatoires :}
\begin{itemize}
\item Patient : Utiliser \zh{我怎么了} (auto-question) ou décrire symptômes.
\item Infirmier : Utiliser \zh{你怎么了？} + \zh{怎么} + V (comment prendre le médicament) + \zh{因为……所以……} (explication posologie).
\item Accompagnateur : Utiliser \zh{怎么} + \zh{不/没} + V (pourquoi pas pris médoc hier ?) + \zh{不但……而且……} (symptômes multiples).
\item Tout le monde : \zh{药怎么吃？} \zh{药房怎么走？}
\end{itemize}
\textbf{Durée :} 5 min préparation, 3 min jeu par groupe. Évaluation par les pairs (grille : justesse structures, prononciation tons, fluidité).
\end{experience}

\coursec{Révision générale du Module 2 : Environnement et santé}

\begin{aretenir}[Check-list des acquis de la Leçon 13]
\begin{itemize}
\item \textbf{Interrogation :} Maîtriser les 3 visages de \zh{怎么} : \zh{怎么了} (état), \zh{怎么}+V (manière), \zh{怎么}+\zh{不/没}+V (raison surprise). Distinguer \zh{什么} et \zh{为什么}.
\item \textbf{Addition :} Construire \zh{不但 A，而且 B} (sujet unique avant \zh{不但} ; sujets différents après chaque conjonction). Placer \zh{也}/\zh{还} avant le verbe dans la 2e proposition.
\item \textbf{Causalité :} Construire \zh{因为 Cause，所以 Conséquence} (ordre fixe). Gérer le(s) sujet(s). Ne pas inverser l'ordre. Connaître l'ellipse \zh{是因为……}.
\item \textbf{Vocabulaire clé :} \zh{脸色、感冒、头疼、出租车、陪、照顾、比赛、取消、堵车、请假、挂号、嗓子、营养}.
\item \textbf{Écriture :} Savoir écrire dans l'ordre des traits : \zh{为} (wèi, 4 traits, point final), \zh{所} (suǒ, 8 traits, haut→bas), \zh{不} (bù, 4 traits), \zh{但} (dàn, 9 traits, radical \zh{人}), \zh{而} (ér, 6 traits), \zh{且} (qiě, 6 traits).
\end{itemize}
\end{aretenir}

\begin{exercice}[Bilan final — Type Examen Probatoire / Baccalauréat]
\textbf{Partie A : Grammaire (10 pts)}
1) Choisissez la bonne réponse (A, B, C, D) :
a) — \_\_\_\_ ? — 我感冒了。 A. 你什么 B. 你怎么 C. 你怎么了 D. 你为什么
b) 法语 \_\_\_\_ 说 « bonjour » ? A. 怎么 B. 什么 C. 为什么 D. 怎么样
c) \_\_\_\_ 他聪明，\_\_\_\_ 他努力，所以考得好。 A. 不但……而且 B. 因为……所以 C. 虽然……但是 D. 如果……就
d) \_\_\_\_ 下雨，\_\_\_\_ 地湿。 A. 因为……所以 B. 不但……而且 C. 怎么……了 D. 什么……了
e) 你 \_\_\_\_ 不吃肉？ A. 怎么了 B. 怎么 C. 什么 D. 为什么

2) Mettez les mots dans l'ordre :
a) 来 / 因为 / 所以 / 没 / 他 / 生病 / 了 / \_\_\_\_\_\_\_\_\_\_\_\_
b) 不但 / 会 / 而且 / 说 / 写 / 他 / 汉语 / 汉字 / \_\_\_\_\_\_\_\_\_\_\_\_

\textbf{Partie B : Traduction (10 pts)}
Traduisez en chinois :
1) Qu'est-ce qui t'arrive ? Tu as l'air inquiet.
2) Comment va-t-on de Yaoundé à Bafoussam ?
3) Pourquoi n'as-tu pas répondu à mon message ? (surprise)
4) Non seulement le ndolé est bon, mais le poulet DG est aussi délicieux.
5) Parce qu'il y avait un embouteillage, j'ai raté l'examen.

\textbf{Partie C : Expression écrite (10 pts)}
Sujet : « Mon ami \zh{小李} n'est pas venu en cours hier. Je lui envoie un message pour prendre de ses nouvelles et lui donner les devoirs. » (Utiliser \zh{怎么了}, \zh{因为……所以……}, \zh{不但……而且……}). 60-80 caractères.
\end{exercice}

\begin{corrige}[Bilan final — Leçon 13]
\textbf{Partie A :}
1) a) C (\zh{你怎么了？}) ; b) A (\zh{怎么} + V) ; c) A (\zh{不但……而且……} addition cause) ; d) A (\zh{因为……所以……}) ; e) B (\zh{怎么} + \zh{不} + V = raison surprise présent).
2) a) \zh{因为他生病了，所以没来。} (ou \zh{他因为生病了，所以没来。}) ; b) \zh{他不但会说汉语，而且会写汉字。}

\textbf{Partie B :}
1) \zh{你怎么了？你看起来很担心。}
2) \zh{从雅温得怎么去巴福萨姆？}
3) \zh{你怎么没回我的信息？}
4) \zh{不但 Ndolé 好吃，而且 鸡肉 DG 也很好吃。} (\zh{鸡肉 DG} = Poulet DG).
5) \zh{因为堵车，所以我错过了考试。} (\zh{错过} = rater/manquer).

\textbf{Partie C :} (Exemple)
\zh{小李，你怎么了？昨天不但没来上课，而且没回微信。因为生病了所以没来吗？今天作业是第13课课文和语法练习。好好休息，明天见！}
\end{corrige}

\coursec{La corrélation 不但……而且……}

Dans la section précédente, nous avons appris à \emph{poser des questions} avec \zh{怎么了} et \zh{怎么} : interroger, c'est chercher une information. Nous allons maintenant apprendre à \emph{enrichir} nos phrases : au lieu de dire une seule chose sur une personne ou une activité, nous dirons deux choses, la seconde étant \emph{plus forte} que la première. C'est exactement ce que fait le français avec « non seulement… mais encore… ». En chinois, cette corrélation s'exprime avec \zh{不但……而且……} (\emph{búdàn… érqiě…}).

\courssub{Observation : découvrir la structure dans un texte}

Lisons d'abord ce court portrait d'une élève de Terminale à Yaoundé. Repérez les mots qui se répètent.

\begin{textebox}[Un portrait : 我的同学玛丽]
我的同学玛丽不但学习好，而且体育也很好。\\
\emph{Wǒ de tóngxué Mǎlì búdàn xuéxí hǎo, érqiě tǐyù yě hěn hǎo.}\\
« Ma camarade Marie est non seulement bonne en études, mais encore très bonne en sport. »\\[4pt]
她不但会说汉语，而且会说英语和法语。\\
\emph{Tā búdàn huì shuō Hànyǔ, érqiě huì shuō Yīngyǔ hé Fǎyǔ.}\\
« Elle sait parler non seulement le chinois, mais encore l'anglais et le français. »\\[4pt]
不但老师喜欢她，而且同学们也都喜欢她。\\
\emph{Búdàn lǎoshī xǐhuan tā, érqiě tóngxuémen yě dōu xǐhuan tā.}\\
« Non seulement les professeurs l'aiment, mais les camarades l'aiment tous aussi. »
\end{textebox}

\begin{center}
\begin{tabularx}{\linewidth}{|X|X|X|X|}\hline
\rowcolor{popblueL} \textbf{Caractères} & \textbf{Pinyin} & \textbf{Nature} & \textbf{Traduction} \\ \hline
玛丽 & Mǎlì & n.pr. & Marie \\ \hline
同学 & tóngxué & n. & camarade de classe \\ \hline
学习 & xuéxí & v. & étudier \\ \hline
体育 & tǐyù & n. & sport, éducation physique \\ \hline
汉语 & Hànyǔ & n. & langue chinoise \\ \hline
英语 & Yīngyǔ & n. & anglais \\ \hline
法语 & Fǎyǔ & n. & français \\ \hline
老师 & lǎoshī & n. & professeur \\ \hline
同学们 & tóngxuémen & n. & camarades (pluriel) \\ \hline
不但 & búdàn & adv. & non seulement \\ \hline
而且 & érqiě & conj. & mais encore, de plus \\ \hline
也 & yě & adv. & aussi \\ \hline
还 & hái & adv. & encore, en plus \\ \hline
\end{tabularx}
\end{center}

Observons ensemble :

\begin{itemize}
\item Dans les deux premières phrases, le sujet (\zh{玛丽}, \zh{她}) est \cle{unique} et se place \emph{avant} \zh{不但}.
\item Dans la troisième phrase, il y a \cle{deux sujets différents} (\zh{老师} et \zh{同学们}) : cette fois, \zh{不但} se place \emph{avant} le premier sujet.
\item Dans tous les cas, la seconde information est \emph{plus forte} ou \emph{plus surprenante} que la première : c'est une \cle{addition progressive}.
\end{itemize}

\courssub{Sens et fonction de la corrélation}

\begin{definition}[不但……而且……]
La corrélation \zh{不但……而且……} (\emph{búdàn… érqiě…}) exprime une \cle{addition progressive} : « non seulement… mais encore… ». Elle relie deux informations portant sur une même réalité, et le \emph{second élément est toujours plus fort, plus remarquable ou plus complet} que le premier. Elle sert à \cle{valoriser} une personne, une activité ou une chose.
\end{definition}

\begin{popfigure}
\begin{tikzpicture}[>=Stealth, node distance=0.6cm]
\node[draw=popblue, fill=popblueL, rounded corners, minimum width=4.2cm, minimum height=1cm, align=center] (e1) at (0,0) {\textbf{Élément 1}\\ \zh{学习好} \emph{xuéxí hǎo}\\ bon en études};
\node[draw=poporange, fill=poporangeL, rounded corners, minimum width=4.2cm, minimum height=1cm, align=center, right=2.6cm of e1] (e2) {\textbf{Élément 2 (plus fort)}\\ \zh{体育也很好} \emph{tǐyù yě hěn hǎo}\\ bon en sport aussi};
\draw[->, very thick, poporange] (e1) -- node[above, align=center]{\zh{而且} : montée\\ d'intensité} (e2);
\end{tikzpicture}
\legende{Flèche d'intensité : le second élément dépasse le premier.}
\end{popfigure}

\begin{attention}[L'ordre des idées compte]
En chinois comme en français, on ne peut pas inverser librement les deux éléments : le plus fort doit venir \emph{après} \zh{而且}. On dira \zh{他不但会说汉语，而且会说三种语言。} (Il parle non seulement le chinois, mais encore trois langues), et non l'inverse, car « trois langues » est plus impressionnant que « le chinois ».
\end{attention}

\courssub{Schéma A : un seul sujet}

\begin{propriete}[Schéma A — même sujet]
Quand les deux informations concernent \cle{le même sujet}, le sujet se place \textbf{avant} \zh{不但} et n'est pas répété (ou est remplacé par un pronom) dans la seconde partie :
\begin{center}
\textbf{Sujet + 不但 + V1/Adj1，而且 + (也/还) + V2/Adj2}
\end{center}
\zh{而且} peut être renforcé par \zh{也} (« aussi ») ou \zh{还} (« encore, de plus »), placés \emph{avant} le verbe ou l'adjectif de la seconde partie.
\end{propriete}

\begin{exemplebox}[Schéma A — un seul sujet]
\zh{他不但学习汉语，而且学习法语。}\\
\emph{Tā búdàn xuéxí Hànyǔ, érqiě xuéxí Fǎyǔ.}\\
« Non seulement il étudie le chinois, mais encore le français. »\\[4pt]
\zh{运动不但对身体好，而且对学习也有帮助。}\\
\emph{Yùndòng búdàn duì shēntǐ hǎo, érqiě duì xuéxí yě yǒu bāngzhù.}\\
« Le sport est bon non seulement pour le corps, mais encore pour les études. »\\[4pt]
\zh{这道菜不但好吃，而且很便宜。}\\
\emph{Zhè dào cài búdàn hǎochī, érqiě hěn piányi.}\\
« Ce plat est non seulement délicieux, mais encore très bon marché. »\\[4pt]
\zh{她不但聪明，还很努力。}\\
\emph{Tā búdàn cōngming, hái hěn nǔlì.}\\
« Elle est non seulement intelligente, mais encore travailleuse. » (variante avec \zh{还} seul)
\end{exemplebox}

Remarquez le dernier exemple : \zh{而且} peut être omis quand \zh{还} ou \zh{也} est présent, mais au Bac il est conseillé de garder la corrélation complète \zh{不但……而且……}, plus claire et plus valorisée.

\courssub{Schéma B : deux sujets différents}

\begin{propriete}[Schéma B — sujets différents]
Quand les deux informations concernent \cle{deux sujets différents}, \zh{不但} se place \textbf{avant le premier sujet}, et le second sujet est placé juste après \zh{而且} (souvent suivi de \zh{也}) :
\begin{center}
\textbf{不但 + Sujet1 + prédicat，而且 + Sujet2 + 也 + prédicat}
\end{center}
\end{propriete}

\begin{exemplebox}[Schéma B — deux sujets différents]
\zh{不但老师会汉语，而且学生也会。}\\
\emph{Búdàn lǎoshī huì Hànyǔ, érqiě xuésheng yě huì.}\\
« Non seulement le professeur parle chinois, mais les élèves aussi. »\\[4pt]
\zh{不但我喜欢中国菜，而且我弟弟也喜欢。}\\
\emph{Búdàn wǒ xǐhuan Zhōngguó cài, érqiě wǒ dìdi yě xǐhuan.}\\
« Non seulement j'aime la cuisine chinoise, mais mon petit frère aussi. »\\[4pt]
\zh{不但喀麦隆人喜欢足球，而且中国人也喜欢。}\\
\emph{Búdàn Kāmàilóng rén xǐhuan zúqiú, érqiě Zhōngguó rén yě xǐhuan.}\\
« Non seulement les Camerounais aiment le football, mais les Chinois aussi. »
\end{exemplebox}

\begin{popfigure}
\begin{tikzpicture}[>=Stealth, node distance=1.1cm]
\node[draw=popblue, fill=popblueL, rounded corners, align=center, minimum width=5.6cm, minimum height=1.5cm] (a) at (0,0) {\textbf{Même sujet}\\ \textbf{S} + \zh{不但} + V1，\zh{而且} + V2\\ \zh{他不但学习汉语，而且学习法语。}};
\node[draw=poppurple, fill=poppurpleL, rounded corners, align=center, minimum width=5.6cm, minimum height=1.5cm, below=of a] (b) {\textbf{Sujets différents}\\ \zh{不但} + \textbf{S1} + V，\zh{而且} + \textbf{S2} + \zh{也} + V\\ \zh{不但老师会汉语，而且学生也会。}};
\node[draw=popdark, fill=popcream, rounded corners, align=center, left=1.6cm of a] (q) {Le sujet\\ est-il le\\ même ?};
\draw[->, thick, popblue] (q) -- node[above]{oui} (a);
\draw[->, thick, poppurple] (q.south) |- node[above, pos=0.75]{non} (b.west);
\end{tikzpicture}
\legende{Schéma comparatif : la place de \zh{不但} dépend du nombre de sujets.}
\end{popfigure}

\courssub{Focus caractères : 不 但 而 且}

\begin{definition}[Analyse des caractères clés]
\begin{itemize}
\item \zh{不} : clé \zh{一} (yī), 4 traits. Sens originel : « non, pas ». C'est la négation de base.
\item \zh{但} : clé \zh{人} (rén), 7 traits. Phonétique \zh{旦} (dàn, aube). Sens : « mais, cependant ». \zh{不但} = « non seulement » (litt. « pas seulement »).
\item \zh{而} : clé \zh{而} (ér), 6 traits. Caractère ancien représentant une barbe. Sens : « et, mais, cependant ». Utilisé comme conjonction de liaison.
\item \zh{且} : clé \zh{一} (yī), 6 traits. Sens originel : « provisoirement, un moment ». Comme conjonction : « et, de plus, ensuite ».
\end{itemize}
\emph{Ordre des traits (rappel) :} horizontal avant vertical (\zh{一} puis \zh{丨}) ; gauche avant droite ; enveloppe avant contenu ; trait de fermeture final.
\end{definition}

\courssub{Comparaison avec le français et pièges}

\begin{center}
\begin{tabularx}{\linewidth}{|X|X|X|}\hline
\rowcolor{popblueL} \textbf{Situation} & \textbf{Français} & \textbf{Chinois} \\ \hline
Même sujet & « Non seulement \emph{il} étudie le chinois… » (le sujet après « non seulement ») & \zh{他不但学习汉语……} (le sujet \textbf{avant} \zh{不但}) \\ \hline
Sujets différents & « Non seulement \emph{le professeur}…, mais \emph{les élèves} aussi » & \zh{不但老师……，而且学生也……} (\zh{不但} avant le sujet 1) \\ \hline
Renforcement & « mais \emph{aussi} / \emph{encore} » & \zh{而且……也/还……} \\ \hline
\end{tabularx}
\end{center}

\begin{attention}[Erreur fréquente des francophones]
En français, « non seulement » précède toujours le sujet : « \emph{Non seulement il} étudie le chinois ». Les élèves francophones transposent cet ordre et écrivent × \zh{不但他学习汉语，而且学习法语。} C'est \textbf{faux} quand le sujet est unique : en chinois, le sujet unique doit passer \emph{avant} \zh{不但} : \zh{他不但学习汉语，而且学习法语。} Réflexe à acquérir : \textbf{un seul sujet → sujet d'abord ; deux sujets → 不但 d'abord}. Autre piège : ne pas oublier \zh{也} dans la seconde partie du schéma B (\zh{而且学生也会}), car les deux sujets font la même action.
\end{attention}

\courssub{Méthode : transformer une phrase simple en corrélation}

\begin{methode}[Construire une phrase avec 不但……而且……]
Pour transformer deux informations simples en une phrase de valorisation :
\begin{enumerate}
\item \textbf{Identifier les deux informations} : ex. « Marie parle chinois » + « Marie parle anglais ».
\item \textbf{Classer par intensité} : l'information la plus forte ira après \zh{而且}.
\item \textbf{Vérifier le sujet} : est-ce le même dans les deux informations ?
\item \textbf{Choisir le schéma} : même sujet → S + \zh{不但}…\zh{而且}… ; sujets différents → \zh{不但} + S1…，\zh{而且} + S2 + \zh{也}…
\item \textbf{Vérifier} la place de \zh{也}/\zh{还} (avant le verbe) et la ponctuation chinoise (，entre les deux parties).
\end{enumerate}
\end{methode}

\begin{exoresolu}[Transformation guidée]
Énoncé : transformez en une seule phrase avec \zh{不但……而且……} : \zh{汉语很有意思。汉语很有用。} (\emph{Hànyǔ hěn yǒu yìsi. Hànyǔ hěn yǒuyòng.} — Le chinois est intéressant. Le chinois est utile.)
\tcblower
Solution détaillée :
\begin{enumerate}
\item Deux informations : « intéressant » et « utile ».
\item Intensité : « utile » est plus fort que « intéressant » → il ira après \zh{而且}.
\item Le sujet est le même : \zh{汉语} → schéma A.
\item Le sujet se place avant \zh{不但} ; on ajoute \zh{也} avant le second adjectif pour renforcer.
\end{enumerate}
Phrase finale : \zh{汉语不但很有意思，而且也很有用。}\\
\emph{Hànyǔ búdàn hěn yǒu yìsi, érqiě yě hěn yǒuyòng.}\\
« Le chinois est non seulement très intéressant, mais encore très utile. »
\end{exoresolu}

\courssub{Emploi au Bac : la phrase de valorisation}

Au baccalauréat, la corrélation \zh{不但……而且……} est un outil précieux pour \emph{valoriser} dans les productions écrites et orales :

\begin{itemize}
\item \textbf{Qualités d'une personne} : \zh{我们的老师不但知识丰富，而且对学生很有耐心。} \emph{Wǒmen de lǎoshī búdàn zhīshi fēngfù, érqiě duì xuésheng hěn yǒu nàixīn.} « Notre professeur a non seulement de grandes connaissances, mais encore beaucoup de patience avec les élèves. »
\item \textbf{Avantages d'une activité (le sport et la santé, thème du module)} : \zh{跑步不但不花钱，而且能让身体健康。} \emph{Pǎobù búdàn bù huā qián, érqiě néng ràng shēntǐ jiànkāng.} « La course ne coûte non seulement rien, mais encore elle rend le corps sain. »
\item \textbf{Avantages d'une langue} : \zh{学汉语不但能了解中国文化，而且能找到好工作。} \emph{Xué Hànyǔ búdàn néng liǎojiě Zhōngguó wénhuà, érqiě néng zhǎodào hǎo gōngzuò.} « Apprendre le chinois permet non seulement de connaître la culture chinoise, mais encore de trouver un bon emploi. »
\end{itemize}

\begin{savaistu}[而且 tout seul]
\zh{而且} peut aussi s'employer \emph{seul}, en début de phrase, pour ajouter une information à ce qui vient d'être dit, comme « de plus » en français : \zh{他是我的老师。而且，他还很年轻。} \emph{Tā shì wǒ de lǎoshī. Érqiě, tā hái hěn niánqīng.} « C'est mon professeur. De plus, il est encore jeune. » En revanche, \zh{不但} ne s'emploie jamais seul : il appelle toujours son partenaire \zh{而且} (ou \zh{也}/\zh{还}). Retenez aussi que dans la langue parlée, on entend souvent la forme raccourcie \zh{不光……而且……} (\emph{bùguāng}), qui a exactement le même sens.
\end{savaistu}

\begin{exercice}[À vous de jouer]
Transformez les paires de phrases suivantes avec \zh{不但……而且……}, en choisissant le bon schéma :
\begin{enumerate}
\item \zh{杜阿拉很大。杜阿拉很热闹。} (Douala est grande. Douala est animée.)
\item \zh{我喜欢音乐。我妹妹喜欢音乐。} (J'aime la musique. Ma petite sœur aime la musique.)
\item \zh{喝水对身体好。喝水不花钱。} (Boire de l'eau est bon pour la santé. Boire de l'eau ne coûte rien.)
\end{enumerate}
\end{exercice}

\begin{corrige}[Exercice « À vous de jouer »]
\begin{enumerate}
\item Même sujet (\zh{杜阿拉}) → schéma A : \zh{杜阿拉不但很大，而且很热闹。} \emph{Dùālā búdàn hěn dà, érqiě hěn rènao.} « Douala est non seulement grande, mais encore très animée. »
\item Sujets différents (\zh{我} / \zh{我妹妹}) → schéma B avec \zh{也} : \zh{不但我喜欢音乐，而且我妹妹也喜欢。} \emph{Búdàn wǒ xǐhuan yīnyuè, érqiě wǒ mèimei yě xǐhuan.} « Non seulement j'aime la musique, mais ma petite sœur aussi. »
\item Même sujet (\zh{喝水}) ; attention à l'intensité : « bon pour la santé » est plus fort → \zh{喝水不但不花钱，而且对身体好。} \emph{Hē shuǐ búdàn bù huā qián, érqiě duì shēntǐ hǎo.} « Boire de l'eau ne coûte non seulement rien, mais encore c'est bon pour la santé. »
\end{enumerate}
\end{corrige}

\begin{aretenir}[L'essentiel de la section]
\begin{itemize}
\item \zh{不但……而且……} = « non seulement… mais encore… » : addition progressive, le second élément est plus fort.
\item \textbf{Même sujet} : Sujet + \zh{不但} + V1，\zh{而且} + V2 (le sujet \emph{avant} \zh{不但}).
\item \textbf{Sujets différents} : \zh{不但} + S1…，\zh{而且} + S2 + \zh{也}… (\zh{不但} \emph{avant} le premier sujet).
\item Renforcement possible avec \zh{也} ou \zh{还} ; \zh{而且} peut aussi ouvrir une phrase seule (« de plus »).
\item Piège majeur du francophone : ne jamais écrire × \zh{不但他……} quand le sujet est unique.
\end{itemize}
\end{aretenir}

\coursec{La corrélation 因为……所以……}

Après la corrélation d'addition \zh{不但……而且……} qui permet d'accumuler deux qualités ou deux faits, nous abordons maintenant une corrélation logique encore plus fréquente : celle qui relie une \cle{cause} à sa \cle{conséquence}. En chinois, cette relation s'exprime grâce au couple \zh{因为……所以……}, littéralement « parce que… donc… ». C'est l'une des structures les plus utiles du niveau HSK 3-4 : elle permet de justifier, d'expliquer, de raisonner — autant d'opérations mentales indispensables à l'oral comme à l'écrit, notamment dans les réponses aux questions « pourquoi ? ».

\courssub{Observation : découvrir la structure dans un texte}

Lisons d'abord ce court texte original. Soulignez mentalement toutes les occurrences de \zh{因为} et de \zh{所以}.

\begin{textebox}[Le match annulé]
上个星期六，我们班本来要去杜阿拉踢足球比赛，但是因为下大雨，所以比赛取消了。\\
\emph{Shàng gè xīngqīliù, wǒmen bān běnlái yào qù Dù'ālā tī zúqiú bǐsài, dànshì yīnwèi xià dà yǔ, suǒyǐ bǐsài qǔxiāo le.}\\
« Samedi dernier, notre classe devait aller à Douala jouer un match de football, mais comme il a beaucoup plu, le match a été annulé. »\\[4pt]
同学们都很失望，因为大家准备了两个星期。\\
\emph{Tóngxuémen dōu hěn shīwàng, yīnwèi dàjiā zhǔnbèi le liǎng gè xīngqī.}\\
« Les camarades étaient tous très déçus, parce que tout le monde s'était préparé pendant deux semaines. »\\[4pt]
老师因为看到我们很难过，所以请我们吃了好吃的恩多莱。\\
\emph{Lǎoshī yīnwèi kàndào wǒmen hěn nánguò, suǒyǐ qǐng wǒmen chī le hǎochī de Ēnduōlái.}\\
« Le professeur, parce qu'il a vu que nous étions tristes, nous a offert un bon ndolé. »\\[4pt]
因为那天我们没踢比赛，所以这个星期六我们要去雅温得比赛。\\
\emph{Yīnwèi nà tiān wǒmen méi tī bǐsài, suǒyǐ zhè gè xīngqīliù wǒmen yào qù Yǎwēndé bǐsài.}\\
« Parce que ce jour-là nous n'avons pas joué, ce samedi nous irons jouer à Yaoundé. »
\end{textebox}

\textbf{Questions d'observation :}
\begin{enumerate}
\item Combien de fois \zh{因为} apparaît-il ? Et \zh{所以} ?
\item Dans chaque phrase, quelle partie exprime la cause ? Quelle partie exprime la conséquence ?
\item \zh{因为} et \zh{所以} peuvent-ils se trouver dans la même phrase ?
\end{enumerate}

On constate que \zh{因为} introduit toujours la \cle{cause} (la pluie, la préparation, la tristesse) et \zh{所以} introduit toujours la \cle{conséquence} (l'annulation, le repas offert, le nouveau match). Et surtout : les deux mots-outils cohabitent volontiers dans la même phrase, ce qui est impossible en français.

\begin{definition}[Vocabulaire clé de la section]
\begin{center}
\begin{tabularx}{\linewidth}{|X|X|X|X|}
\hline
\rowcolor{popblueL} \textbf{Caractères} & \textbf{Pinyin} & \textbf{Nature} & \textbf{Traduction} \\
\hline
因为 & yīnwèi & conj. & parce que, comme \\
\hline
所以 & suǒyǐ & adv. & donc, par conséquent \\
\hline
杜阿拉 & Dù'ālā & n.p. & Douala \\
\hline
雅温得 & Yǎwēndé & n.p. & Yaoundé \\
\hline
取消 & qǔxiāo & v. & annuler \\
\hline
失望 & shīwàng & adj. & déçu \\
\hline
准备 & zhǔnbèi & v. & préparer \\
\hline
难过 & nánguò & adj. & triste, peiné \\
\hline
请 & qǐng & v. & inviter (offrir un repas) \\
\hline
恩多莱 & Ēnduōlái & n. & ndolé (plat camerounais) \\
\hline
生病 & shēngbìng & v. & tomber malade \\
\hline
上课 & shàngkè & v. & avoir cours \\
\hline
天气 & tiānqì & n. & météo, temps qu'il fait \\
\hline
越来越多 & yuè lái yuè duō & loc. & de plus en plus \\
\hline
为什么 & wèishénme & adv. & pourquoi \\
\hline
通过 & tōngguò & v. & réussir (un examen), passer \\
\hline
考试 & kǎoshì & n. & examen \\
\hline
足球 & zúqiú & n. & football \\
\hline
累 & lèi & adj. & fatigué \\
\hline
睡觉 & shuìjiào & v. & dormir \\
\hline
\end{tabularx}
\end{center}
\end{definition}

\courssub{La règle : forme et emplois}

\begin{propriete}[Structure de la corrélation 因为……所以……]
\textbf{Schéma général :}

\zh{因为} + proposition de cause，\zh{所以} + proposition de conséquence。

\begin{itemize}
\item \zh{因为} (yīnwèi) = « parce que, comme, car » : il introduit la \cle{cause}.
\item \zh{所以} (suǒyǐ) = « donc, c'est pourquoi, par conséquent » : il introduit la \cle{conséquence}.
\item Contrairement au français, les deux mots-outils \textbf{peuvent et doivent souvent apparaître ensemble} dans la même phrase. C'est même la forme la plus naturelle et la plus soutenue en chinois.
\end{itemize}
\end{propriete}

\begin{exemplebox}[Exemple fondateur]
\zh{因为他病了，所以没来上课。}\\
\emph{Yīnwèi tā bìng le, suǒyǐ méi lái shàngkè.}\\
« Parce qu'il est malade, il n'est pas venu en cours. »\\[4pt]
Analyse : \zh{他病了} = cause (il est tombé malade) ; \zh{没来上课} = conséquence (il n'est pas venu en cours). En français, on dirait « Comme il est malade, il n'est pas venu » ou « Il est malade, donc il n'est pas venu », mais jamais « Parce qu'il est malade, donc il n'est pas venu ». En chinois, si : c'est la construction standard.
\end{exemplebox}

\begin{popfigure}
\begin{tikzpicture}[
  cause/.style={rectangle, rounded corners, draw=popblue, fill=popblueL, thick, align=center, minimum width=4.6cm, minimum height=1.5cm, font=\small},
  consq/.style={rectangle, rounded corners, draw=poporange, fill=poporangeL, thick, align=center, minimum width=4.6cm, minimum height=1.5cm, font=\small},
  fleche/.style={-{Stealth[length=4mm]}, very thick, popdark}
]
\node[cause, align=center] (c) at (0,0){\textbf{CAUSE}\\ \zh{因为} yīnwèi\\ « parce que »};
\node[consq, align=center] (q) at (8,0){\textbf{CONSÉQUENCE}\\ \zh{所以} suǒyǐ\\ « donc »};
\draw[fleche] (c) -- node[above, font=\small\itshape]{entraîne} (q);
\node[below=0.35cm of c, font=\footnotesize, align=center, text=popmuted]{\zh{因为他病了}\\ (il est malade)};
\node[below=0.35cm of q, font=\footnotesize, align=center, text=popmuted]{\zh{所以没来上课}\\ (il n'est pas venu)};
\end{tikzpicture}
\legende{Organigramme logique : la cause (\zh{因为}) entraîne la conséquence (\zh{所以}).}
\end{popfigure}

\courssub{Les trois emplois : corrélation complète ou mots-outils isolés}

La corrélation complète \zh{因为……所以……} est la forme de base, mais chaque mot-outil peut aussi être employé seul, selon le contexte de communication.

\textbf{1. Corrélation complète (forme la plus courante à l'écrit)}

\begin{exemplebox}[因为 + 所以 ensemble]
\zh{因为天气很热，所以我喝很多水。}\\
\emph{Yīnwèi tiānqì hěn rè, suǒyǐ wǒ hē hěn duō shuǐ.}\\
« Comme il fait très chaud, je bois beaucoup d'eau. »\\[4pt]
\zh{因为汉语很有意思，所以越来越多的喀麦隆学生学习汉语。}\\
\emph{Yīnwèi Hànyǔ hěn yǒu yìsi, suǒyǐ yuè lái yuè duō de Kāmàilóng xuésheng xuéxí Hànyǔ.}\\
« Parce que le chinois est très intéressant, de plus en plus d'élèves camerounais l'apprennent. »
\end{exemplebox}

\textbf{2. \zh{因为} seul : répondre à une question en \zh{为什么}}

Quand on répond à la question « pourquoi ? », on utilise \zh{因为} seul ; \zh{所以} est inutile dans une réponse courte.

\begin{exemplebox}[Dialogue : question 为什么 → réponse 因为]
— \zh{你为什么学汉语？}\\
\emph{Nǐ wèishénme xué Hànyǔ?}\\
« Pourquoi étudies-tu le chinois ? »\\[4pt]
— \zh{因为我想去中国。}\\
\emph{Yīnwèi wǒ xiǎng qù Zhōngguó.}\\
« Parce que je veux aller en Chine. »
\end{exemplebox}

\textbf{3. \zh{所以} seul : tirer une conclusion}

Quand la cause vient d'être mentionnée dans le discours, on peut enchaîner directement avec \zh{所以} pour conclure.

\begin{exemplebox}[所以 seul en conclusion]
\zh{明天下雨，所以我们不去市场了。}\\
\emph{Míngtiān xià yǔ, suǒyǐ wǒmen bù qù shìchǎng le.}\\
« Demain il pleut, donc nous n'irons pas au marché. »\\[4pt]
Ici, la cause est un simple constat sans \zh{因为} ; \zh{所以} annonce la décision qui en découle.
\end{exemplebox}

\courssub{La place du sujet : une subtilité importante}

Quand les deux propositions ont le \cle{même sujet}, celui-ci peut être placé \textbf{avant} \zh{因为}, et il n'est pas répété après \zh{所以}.

\begin{propriete}[Place du sujet dans la corrélation]
Deux schémas possibles quand le sujet est identique :
\begin{itemize}
\item Sujet + \zh{因为} + cause，\zh{所以} + conséquence (sujet non répété) ;
\item \zh{因为} + sujet + cause，\zh{所以} + (sujet) + conséquence.
\end{itemize}
Quand les sujets sont différents, chaque proposition garde son propre sujet.
\end{propriete}

\begin{exemplebox}[Même sujet, deux formulations correctes]
\zh{他因为生病，所以没来。}\\
\emph{Tā yīnwèi shēngbìng, suǒyǐ méi lái.}\\
« Parce qu'il était malade, il n'est pas venu. » (sujet \zh{他} placé avant \zh{因为})\\[4pt]
\zh{因为他生病，所以他没来。}\\
\emph{Yīnwèi tā shēngbìng, suǒyǐ tā méi lái.}\\
Même sens, sujet répété : correct aussi, mais plus lourd.\\[6pt]
Sujets différents :\\
\zh{因为路不好，所以公共汽车来晚了。}\\
\emph{Yīnwèi lù bù hǎo, suǒyǐ gōnggòng qìchē lái wǎn le.}\\
« Parce que la route est mauvaise, le bus est arrivé en retard. » (sujet 1 : la route ; sujet 2 : le bus)
\end{exemplebox}

\courssub{Comparaison avec le français : le grand piège}

\begin{center}
\begin{tabularx}{\linewidth}{|X|X|X|}
\hline
\rowcolor{popblueL} \textbf{Situation} & \textbf{Français} & \textbf{Chinois} \\
\hline
Cause + conséquence liées & « Comme il pleut, je reste. » (un seul mot-outil) & \zh{因为下雨，所以我待在家里。} (deux mots-outils) \\
\hline
« Parce que… donc… » & \textbf{Interdit} en français standard & \textbf{Naturel et recommandé} : \zh{因为……所以……} \\
\hline
Réponse à « pourquoi ? » & « Parce que… » & \zh{因为……} (seul) \\
\hline
Conclusion & « … donc … » & \zh{……所以……} (seul) \\
\hline
\end{tabularx}
\end{center}

\begin{attention}[Erreur n°1 des francophones : supprimer 所以]
Par calque du français, les élèves écrivent souvent : *\zh{因为他病了，没来上课。} en supprimant \zh{所以}. Cette phrase est compréhensible, mais en chinois soigné — et à l'examen — il faut \textbf{garder les deux mots-outils} : \zh{因为他病了，所以没来上课。} Retenez : ce qui est une faute en français (« parce que… donc… ») est une \textbf{obligation de style} en chinois.
\end{attention}

\begin{attention}[Erreur n°2 : inverser cause et conséquence]
\zh{因为} introduit \textbf{toujours la cause}, jamais le résultat. Ne dites pas *\zh{因为他没来上课，所以病了} pour dire « il n'est pas venu parce qu'il était malade » : vous affirmez alors l'inverse (il est tombé malade parce qu'il n'est pas venu !). Avant d'écrire, demandez-vous : « Qu'est-ce qui arrive \emph{d'abord} dans la réalité ? » C'est la cause, donc c'est \zh{因为}.
\end{attention}

\begin{savaistu}[Le ndolé et le football : Cameroun-Chine]
Le \zh{恩多莱} (ndolé) est le plat national du Cameroun, à base de feuilles de vernonia amères, d'arachides et de viande ou poisson. Dans le texte, le professeur l'offre pour consoler les élèves : un geste typique de l'hospitalité camerounaise. En Chine, le football (\zh{足球}) est le sport le plus suivi à la télévision, bien que le niveau de la ligue nationale (CSL) soit encore en développement. La Chine a participé à la Coupe du monde 2002 ; le Cameroun (les \emph{Lions Indomptables}) y a brillé dès 1990. Le sport crée un pont culturel fort entre Yaoundé et Pékin.
\end{savaistu}

\begin{methode}[Traduire « comme…, … » du français]
Pour traduire une phrase française de type « Comme il fait chaud, je bois beaucoup d'eau » :
\begin{enumerate}
\item Repérez la cause (« il fait chaud ») et la conséquence (« je bois beaucoup d'eau »).
\item Traduisez la cause en la faisant précéder de \zh{因为} : \zh{因为天气很热}.
\item Traduisez la conséquence en la faisant précéder de \zh{所以} : \zh{所以我喝很多水}.
\item Assemblez avec une virgule chinoise ：\zh{因为天气很热，所以我喝很多水。}
\item Vérifiez la place du sujet : s'il est unique, vous pouvez le placer avant \zh{因为} et ne pas le répéter.
\end{enumerate}
\end{methode}

\begin{exoresolu}[Traduction guidée]
Énoncé : traduisez en chinois « Comme le match a été annulé, les élèves sont rentrés à la maison. » (match = \zh{比赛} ; annuler = \zh{取消} ; rentrer = \zh{回家}).
\tcblower
Solution détaillée :
\begin{enumerate}
\item Cause : « le match a été annulé » → \zh{比赛取消了} ; avec l'outil de cause : \zh{因为比赛取消了}.
\item Conséquence : « les élèves sont rentrés à la maison » → \zh{同学们回家了} ; avec l'outil de conséquence : \zh{所以同学们回家了}.
\item Assemblage (sujets différents, on les garde tous les deux) :
\end{enumerate}
\zh{因为比赛取消了，所以同学们回家了。}\\
\emph{Yīnwèi bǐsài qǔxiāo le, suǒyǐ tóngxuémen huí jiā le.}\\
Vérification : \zh{因为} introduit bien la cause (l'annulation), \zh{所以} la conséquence (le retour à la maison). Aucune inversion.
\end{exoresolu}

\courssub{Entraînement oral guidé}

\begin{experience}[Jeu de rôle : « Pourquoi ? — Parce que… »]
\textbf{Consigne :} En binôme. L'élève A pioche une carte « Situation », l'élève B demande \zh{为什么？}, l'élève A répond avec \zh{因为……} (seul), puis ajoute une conséquence avec \zh{所以……} (corrélation complète).
\textbf{Cartes situations (exemples) :}
\begin{itemize}
\item Tu n'es pas venu en cours hier. (Tu étais malade.)
\item Tu bois beaucoup d'eau. (Il fait très chaud.)
\item Tu étudies le chinois tous les soirs. (Tu veux aller en Chine.)
\item Le match est annulé. (Il pleut.)
\item Tu es en retard. (Le bus est en panne / la route est mauvaise.)
\end{itemize}
\textbf{Modèle :}
A: \zh{你为什么没来上课？} \\
B: \zh{因为我病了。因为我病了，所以我没来上课。}
\end{experience}

\begin{exercice}[À vous de construire]
\begin{enumerate}
\item Traduisez : « Parce qu'elle étudie beaucoup, elle réussit à l'examen. » (réussir = \zh{通过} ; examen = \zh{考试})
\item Répondez à la question \zh{你为什么喜欢足球？} avec \zh{因为} seul.
\item Complétez : \zh{因为今天是星期一，所以……} (inventez une conséquence logique).
\item Corrigez la faute : *\zh{因为他很累，所以所以他想睡觉。}
\item Transformez en corrélation complète : \zh{明天有考试，我不去看电影。} (Ajoutez \zh{因为} et \zh{所以}).
\end{enumerate}
\end{exercice}

\begin{corrige}[Exercice « À vous de construire »]
\begin{enumerate}
\item \zh{因为她学习很努力，所以通过了考试。} \emph{Yīnwèi tā xuéxí hěn nǔlì, suǒyǐ tōngguò le kǎoshì.} (même sujet : on peut omettre \zh{她} après \zh{所以} → \zh{因为她学习很努力，所以通过了考试。}).
\item Exemple : \zh{因为足球很有意思。} \emph{Yīnwèi zúqiú hěn yǒu yìsi.} « Parce que le football est très intéressant. »
\item Exemple : \zh{因为今天是星期一，所以我们要上课。} \emph{Yīnwèi jīntiān shì xīngqīyī, suǒyǐ wǒmen yào shàngkè.} « Parce qu'aujourd'hui c'est lundi, nous avons cours. »
\item La faute est la répétition de \zh{所以}. Correction : \zh{因为他很累，所以他想睡觉。} \emph{Yīnwèi tā hěn lèi, suǒyǐ tā xiǎng shuìjiào.} « Parce qu'il est fatigué, il veut dormir. »
\item \zh{因为明天有考试，所以我不去看电影。} \emph{Yīnwèi míngtiān yǒu kǎoshì, suǒyǐ wǒ bù qù kàn diànyǐng.}
\end{enumerate}
\end{corrige}

\begin{aretenir}[L'essentiel de la corrélation 因为……所以……]
\begin{itemize}
\item \zh{因为} + cause，\zh{所以} + conséquence : les deux mots-outils coexistent naturellement, contrairement au français.
\item Réponse à \zh{为什么} ? → \zh{因为} seul. Conclusion → \zh{所以} seul.
\item Même sujet : il peut précéder \zh{因为} et ne pas être répété.
\item Jamais \zh{因为} devant la conséquence : vérifiez toujours l'ordre logique cause → effet.
\end{itemize}
\end{aretenir}

\coursec{Comparaison avec le français et pièges}

Nous venons d'étudier la corrélation \zh{因为……所以……} : vous avez vu qu'en chinois, les deux moitiés de la paire aiment voyager ensemble. Cette particularité n'est pas un détail : elle illustre une différence profonde entre le chinois et le français. Le français est une langue qui aime \emph{choisir} ses mots-outils (on dit « parce que » \emph{ou} « donc », rarement les deux), qui aime \emph{déplacer} ses mots (« non seulement il chante, mais encore il danse », avec inversion) ; le chinois, lui, est une langue de \emph{position fixe} et de \emph{paires obligatoires}. Cette section fait le point : pour chaque structure de la leçon, nous confrontons le chinois au français, puis nous dressons la liste des erreurs réellement relevées dans les copies camerounaises, avec leurs corrections.

\begin{definition}[Mots-outils et termes grammaticaux de la leçon]
\begin{center}
\begin{tabularx}{\linewidth}{|X|X|X|X|}
\hline
\rowcolor{popblueL} Caractères & Pinyin & Nature & Traduction \\
\hline
怎么了 & zěnme le & locution interrogative & Qu'est-ce qu'il y a ? / Qu'est-ce qui ne va pas ? \\
\hline
怎么 & zěnme & adv. interrogatif & Comment / Pourquoi (selon structure) \\
\hline
不但 & búdàn & conjonction (corrélatif) & Non seulement \\
\hline
而且 & érqiě & conjonction (corrélatif) & Mais aussi / De plus \\
\hline
因为 & yīnwèi & conjonction (corrélatif) & Parce que / Comme \\
\hline
所以 & suǒyǐ & conjonction (corrélatif) & Donc / C'est pourquoi \\
\hline
\end{tabularx}
\end{center}
\end{definition}

\courssub{Synthèse contrastive : \zh{怎么了} et \zh{怎么} ne veulent pas dire ce que vous croyez}

\begin{exemplebox}[Trois phrases à comparer avant toute règle]
\zh{你怎么了？} \emph{Nǐ zěnme le?} « Qu'est-ce que tu as ? / Qu'est-ce qui ne va pas ? »\\
\zh{这个字怎么写？} \emph{Zhège zì zěnme xiě?} « Comment s'écrit ce caractère ? »\\
\zh{他怎么没来？} \emph{Tā zěnme méi lái?} « Pourquoi n'est-il pas venu ? »
\end{exemplebox}

Observez : le même mot \zh{怎么} apparaît trois fois, mais il ne se traduit jamais pareil. Tout dépend de ce qui l'entoure.

\begin{propriete}[Les trois emplois de \zh{怎么} face au français]
\begin{enumerate}
\item \zh{怎么了} (avec \zh{了} final, souvent sans verbe) = s'enquérir d'un \cle{problème}, d'un changement inattendu : « Qu'est-ce qu'il y a ? Qu'est-ce que tu as ? ». Ce n'est \textbf{jamais} une formule de salutation.
\item \zh{怎么} + verbe = interroger sur la \cle{manière}, le procédé : « comment ». Structure : Sujet + \zh{怎么} + Verbe (+ complément).
\item \zh{怎么} + négation (\zh{不} / \zh{没}) = interroger sur la \cle{raison} d'un fait contraire à l'attente : « comment se fait-il que…, pourquoi… ». Structure : Sujet + \zh{怎么} + \zh{不/没} + Verbe.
\end{enumerate}
\end{propriete}

\begin{attention}[\zh{怎么了} n'est pas « Comment vas-tu ? »]
C'est l'erreur numéro un des francophones. « Comment vas-tu ? » se dit \zh{你好吗？} \emph{Nǐ hǎo ma?} ou \zh{你怎么样？} \emph{Nǐ zěnmeyàng?}. Si vous dites \zh{你怎么了？} à un ami en pleine forme, vous lui demandez en réalité « Qu'est-ce qui ne va pas chez toi ? » — comme s'il avait l'air malade ! Réservez \zh{怎么了} aux situations où quelque chose cloche visiblement : un camarade qui pleure, un téléphone en panne, un bruit étrange.
\end{attention}

\begin{exemplebox}[Même verbe, deux questions différentes]
\zh{你怎么学习汉语？} \emph{Nǐ zěnme xuéxí Hànyǔ?} « Comment étudies-tu le chinois ? » (manière : avec quelle méthode ?)\\
\zh{你怎么不学习汉语？} \emph{Nǐ zěnme bù xuéxí Hànyǔ?} « Pourquoi n'étudies-tu pas le chinois ? » (raison : c'est étonnant !)
\end{exemplebox}

\courssub{La corrélation double obligatoire : \zh{因为……所以……} contre le français}

\begin{exemplebox}[Observez la traduction mot à mot]
\zh{因为下雨了，所以足球比赛取消了。}\\
\emph{Yīnwèi xià yǔ le, suǒyǐ zúqiú bǐsài qǔxiāo le.}\\
Mot à mot : « Parce qu'il a plu, \emph{donc} le match de football a été annulé. »\\
Français correct : « Comme il a plu, le match a été annulé. »
\end{exemplebox}

En français, la grammaire scolaire interdit « parce que… donc… » dans la même phrase : on choisit l'un \emph{ou} l'autre. En chinois, c'est l'inverse : la paire \zh{因为……所以……} est la forme \cle{pleine}, attendue, notamment à l'écrit. Le chinois « voit » la relation cause-conséquence comme un pont à deux piliers ; en retirer un, c'est affaiblir la phrase.

\begin{propriete}[Règle de la paire causale]
Structure pleine : \zh{因为} + cause，\zh{所以} + conséquence。\\
À l'\textbf{écrit} (rédaction, épreuve du Bac), utilisez la paire complète. À l'\textbf{oral} rapide, on peut n'en garder qu'un (\zh{因为下雨，我不去。} ou \zh{下雨了，所以我不去。}), mais la forme complète n'est jamais fautive.
\end{propriete}

\begin{attention}[Le calque français qui coûte des points]
× \zh{因为下雨，我不去。} \emph{Yīnwèi xià yǔ, wǒ bù qù.} — acceptable à l'oral, mais \textbf{incomplet à l'écrit du Bac} : le correcteur attend la paire.\\
✓ \zh{因为下雨，所以我不去。} \emph{Yīnwèi xià yǔ, suǒyǐ wǒ bù qù.} « Parce qu'il pleut, je n'y vais pas. »\\
Erreur inverse, par calque de l'interdiction française : écrire \zh{所以} seul alors que la cause est exprimée. Si la cause est là, mettez \zh{因为} ; si la conséquence est là, mettez \zh{所以}.
\end{attention}

\courssub{La place du sujet dans \zh{不但……而且……} : la règle qui n'existe pas en français}

Voici la subtilité la plus délicate de la leçon. En français, « non seulement… mais aussi… » ne pose aucun problème de place de sujet. En chinois, tout dépend d'une question simple : \textbf{les deux moitiés de la phrase parlent-elles de la même personne ?}

\begin{propriete}[Sujet unique ou sujets différents : deux schémas]
\textbf{Cas 1 — même sujet dans les deux propositions} : le sujet se place \cle{avant} \zh{不但}, une seule fois, et n'est pas répété (ou remplacé par un pronom) après \zh{而且}.\\
Schéma : Sujet + \zh{不但} + prédicat 1，\zh{而且} + (pronom) + prédicat 2。\\
\zh{我不但会唱歌，而且会跳舞。} \emph{Wǒ búdàn huì chànggē, érqiě huì tiàowǔ.} « Non seulement je sais chanter, mais je sais aussi danser. »\\[4pt]
\textbf{Cas 2 — sujets différents} : \zh{不但} se place \cle{avant} le premier sujet, et le second sujet apparaît après \zh{而且}.\\
Schéma : \zh{不但} + Sujet 1 + prédicat 1，\zh{而且} + Sujet 2 + prédicat 2。\\
\zh{不但我会唱歌，而且我妹妹也会唱歌。} \emph{Búdàn wǒ huì chànggē, érqiě wǒ mèimei yě huì chànggē.} « Non seulement moi je sais chanter, mais ma petite sœur aussi. »
\end{propriete}

\begin{attention}[L'erreur de placement du sujet unique]
× \zh{不但我会唱歌，而且会跳舞。} — fautif : avec un sujet unique (\zh{我} pour les deux actions), le sujet doit \textbf{précéder} \zh{不但}. En plaçant \zh{不但} en tête, vous annoncez que deux sujets différents vont suivre… et le second n'arrive jamais.\\
✓ \zh{我不但会唱歌，而且会跳舞。}\\
Moyen mnémotechnique : comptez les sujets. \textbf{Un seul sujet} → il passe devant tout le monde : Sujet + \zh{不但}…\zh{而且}… \textbf{Deux sujets} → chacun reste dans son camp : \zh{不但} + S1…，\zh{而且} + S2…
\end{attention}

\begin{exemplebox}[Les deux cas côte à côte]
\zh{雅温得不但很大，而且很现代化。} \emph{Yǎwēndé búdàn hěn dà, érqiě hěn xiàndàihuà.} « Yaoundé est non seulement très grande, mais aussi très moderne. » (un seul sujet : Yaoundé)\\
\zh{不但学生喜欢这位老师，而且校长也喜欢他。} \emph{Búdàn xuésheng xǐhuan zhè wèi lǎoshī, érqiě xiàozhǎng yě xǐhuan tā.} « Non seulement les élèves aiment ce professeur, mais le proviseur l'aime aussi. » (deux sujets : élèves / proviseur)
\end{exemplebox}

\courssub{Ordre des mots figé : l'inversion française est impossible en chinois}

Le français peut « faire le spectacle » avec sa syntaxe : « \emph{Non seulement} il chante, \emph{mais encore} il danse » — avec inversion du sujet et du verbe, montée en tête de phrase, effet de style. Le chinois refuse tout cela : les mots-outils \zh{不但}, \zh{而且}, \zh{因为}, \zh{所以} ont une \cle{place fixe} dans la phrase, et l'ordre Sujet–Verbe–Complément ne bouge jamais.

\begin{center}
\begin{tabularx}{\linewidth}{|X|X|X|}
\hline
\rowcolor{popblueL} Structure & Exemple (caractères + pinyin) & Traduction et remarque \\
\hline
\zh{怎么了} & \zh{你的自行车怎么了？} \emph{Nǐ de zìxíngchē zěnme le?} & « Qu'est-ce que ta bicyclette a ? » — \zh{怎么了} reste en fin, jamais en tête. \\
\hline
\zh{怎么} + verbe & \zh{怎么去杜阿拉？} \emph{Zěnme qù Dùālā?} & « Comment aller à Douala ? » — \zh{怎么} juste avant le verbe. \\
\hline
\zh{不但……而且……} & \zh{他不但聪明，而且很努力。} \emph{Tā búdàn cōngming, érqiě hěn nǔlì.} & « Il est non seulement intelligent, mais aussi travailleur. » — pas d'inversion possible. \\
\hline
\zh{因为……所以……} & \zh{因为他病了，所以没来上课。} \emph{Yīnwèi tā bìng le, suǒyǐ méi lái shàngkè.} & « Parce qu'il était malade, il n'est pas venu en cours. » — cause d'abord, conséquence ensuite, toujours dans cet ordre. \\
\hline
\end{tabularx}
\end{center}

\begin{attention}[Trois déplacements impossibles]
\begin{itemize}
\item × Mettre \zh{所以} en tête de phrase pour « donc » français : « Donc, je reste » ≠ \zh{所以，我…} en début de texte sans cause exprimée ; en chinois, \zh{所以} introduit la conséquence d'une cause présente dans le contexte immédiat.
\item × Inverser cause et conséquence : \zh{所以我不去，因为下雨} est impossible ; le chinois dit toujours la cause d'abord.
\item × Placer \zh{怎么} après le verbe comme l'anglais ou le français placent parfois « comment » en fin de phrase familière : « Tu pars comment ? » se dit obligatoirement \zh{你怎么走？} \emph{Nǐ zěnme zǒu?}
\end{itemize}
\end{attention}

\courssub{Les cinq erreurs types des candidats camerounais}

Voici les fautes les plus fréquentes relevées dans les copies de Terminale, avec leur correction. Étudiez-les : ce sont exactement celles que le correcteur guette.

\begin{popfigure}
\begin{tikzpicture}
\node[fill=popink, text=white, font=\bfseries, minimum width=7.2cm, minimum height=0.8cm] (h1) at (0,0) {Erreur fréquente};
\node[fill=popgreen, text=white, font=\bfseries, minimum width=7.2cm, minimum height=0.8cm] (h2) at (7.6,0) {Correction};
\node[fill=popblueL, minimum width=7.2cm, minimum height=1.1cm, align=center] (a1) at (0,-1.05) {\zh{老师，你怎么了？}\\ \emph{(pour saluer le professeur)}};
\node[fill=popgreenL, minimum width=7.2cm, minimum height=1.1cm, align=center] (b1) at (7.6,-1.05) {\zh{老师，你好吗？}\\ \emph{Lǎoshī, nǐ hǎo ma?}};
\node[fill=popblueL, minimum width=7.2cm, minimum height=1.1cm, align=center] (a2) at (0,-2.35) {\zh{你学习怎么汉语？}\\ \emph{(ordre calqué du français)}};
\node[fill=popgreenL, minimum width=7.2cm, minimum height=1.1cm, align=center] (b2) at (7.6,-2.35) {\zh{你怎么学习汉语？}\\ \emph{Nǐ zěnme xuéxí Hànyǔ?}};
\node[fill=popblueL, minimum width=7.2cm, minimum height=1.1cm, align=center] (a3) at (0,-3.65) {\zh{不但我会唱歌，而且会跳舞。}\\ \emph{(sujet unique mal placé)}};
\node[fill=popgreenL, minimum width=7.2cm, minimum height=1.1cm, align=center] (b3) at (7.6,-3.65) {\zh{我不但会唱歌，而且会跳舞。}\\ \emph{Wǒ búdàn huì chànggē, érqiě huì tiàowǔ.}};
\node[fill=popblueL, minimum width=7.2cm, minimum height=1.1cm, align=center] (a4) at (0,-4.95) {\zh{因为下雨，我不去。}\\ \emph{(paire incomplète à l'écrit)}};
\node[fill=popgreenL, minimum width=7.2cm, minimum height=1.1cm, align=center] (b4) at (7.6,-4.95) {\zh{因为下雨，所以我不去。}\\ \emph{Yīnwèi xià yǔ, suǒyǐ wǒ bù qù.}};
\node[fill=popblueL, minimum width=7.2cm, minimum height=1.1cm, align=center] (a5) at (0,-6.25) {\zh{他所以没来，因为病了。}\\ \emph{(ordre cause/conséquence inversé)}};
\node[fill=popgreenL, minimum width=7.2cm, minimum height=1.1cm, align=center] (b5) at (7.6,-6.25) {\zh{因为他病了，所以没来。}\\ \emph{Yīnwèi tā bìng le, suǒyǐ méi lái.}};
\draw[popline, thick] (h1.south west) -- (a5.south west);
\draw[popline, thick] (h2.south west) -- (b5.south west);
\end{tikzpicture}
\legende{Tableau « Erreur → Correction » : les cinq fautes classiques des copies camerounaises.}
\end{popfigure}

\begin{exoresolu}[Corriger une phrase de copie réelle]
Énoncé : un candidat a écrit \zh{不但我的学校很大，而且很漂亮。} pour dire « Mon école est non seulement très grande, mais aussi très belle. » Relevez l'erreur, expliquez-la, corrigez.
\tcblower
Solution détaillée :
\begin{enumerate}
\item \textbf{Diagnostic} : les deux prédicats (\zh{很大} « très grande » et \zh{很漂亮} « très belle ») qualifient le \textbf{même} sujet, \zh{我的学校} « mon école ». Nous sommes donc dans le cas du sujet unique.
\item \textbf{Règle} : avec un sujet unique, celui-ci se place \emph{avant} \zh{不但} et n'apparaît qu'une fois.
\item \textbf{Correction} : \zh{我的学校不但很大，而且很漂亮。} \emph{Wǒ de xuéxiào búdàn hěn dà, érqiě hěn piàoliang.}
\item \textbf{Vérification} : la phrase du candidat serait correcte seulement si un second sujet suivait \zh{而且}, par exemple : \zh{不但我的学校很大，而且城里的中学也很大。} « Non seulement mon école est grande, mais les lycées de la ville aussi. »
\end{enumerate}
\end{exoresolu}

\courssub{Stratégie de vérification avant de rendre la copie}

\begin{methode}[L'auto-correction en quatre questions]
Avant de rendre votre copie, relisez chaque phrase contenant une structure de cette leçon et posez-vous, dans l'ordre, ces quatre questions :
\begin{enumerate}
\item \textbf{Le sujet} : dans \zh{不但……而且……}, y a-t-il un ou deux sujets ? Si un seul, est-il bien placé \emph{avant} \zh{不但} ?
\item \textbf{La paire de mots-outils} : si j'ai écrit \zh{因为}, ai-je écrit \zh{所以} ? Si j'ai écrit \zh{不但}, ai-je écrit \zh{而且} ? Lisez la phrase à voix haute : l'oreille francophone entend vite qu'il « manque » le deuxième pilier du pont.
\item \textbf{Le \zh{了} final} : dans une question en \zh{怎么了}, le \zh{了} est-il présent ? Dans une cause passée (\zh{因为……了，所以……}), ai-je marqué l'accompli là où il le faut ?
\item \textbf{La place de \zh{怎么}} : est-il bien \emph{avant} le verbe (manière) ou \emph{avant la négation} (raison), jamais après ?
\end{enumerate}
\end{methode}

\begin{exemplebox}[Application de la méthode à une phrase fautive]
Phrase de copie : \zh{因为天气很热，我所以喝很多水。}\\
Question 2 (paire) : \zh{因为} et \zh{所以} sont présents, mais… question 1 bis : la place ! \zh{所以} doit précéder le sujet de la conséquence ou le suivre immédiatement en tête de proposition, jamais se coincer après le sujet. Correction : \zh{因为天气很热，所以我喝很多水。} \emph{Yīnwèi tiānqì hěn rè, suǒyǐ wǒ hē hěn duō shuǐ.} « Comme il fait très chaud, je bois beaucoup d'eau. »
\end{exemplebox}

\begin{popfigure}
\begin{tikzpicture}[mindmap, grow cyclic, every node/.style={concept}, concept color=popblue, text=white,
level 1/.append style={level distance=3.4cm, sibling angle=120},
level 2/.append style={level distance=2.6cm, sibling angle=50}]
\node[concept color=popdark, font=\bfseries] {Grammaire Leçon 13}
child[concept color=poporange] { node {\zh{怎么}} 
  child { node[concept color=poporangeL, text=popdark, align=center] {\zh{怎么了？}\\ problème} }
  child { node[concept color=poporangeL, text=popdark, align=center] {\zh{怎么}+V\\ manière} }
  child { node[concept color=poporangeL, text=popdark, align=center] {\zh{怎么}+\zh{不/没}\\ raison} }
}
child[concept color=popgreen] { node {\zh{不但……而且……}}
  child { node[concept color=popgreenL, text=popdark, align=center] {1 sujet :\\ S+\zh{不但}…\zh{而且}…} }
  child { node[concept color=popgreenL, text=popdark, align=center] {2 sujets :\\ \zh{不但}+S1…\zh{而且}+S2…} }
}
child[concept color=poppurple] { node {\zh{因为……所以……}}
  child { node[concept color=poppurpleL, text=popdark, align=center] {cause d'abord,\\ conséquence ensuite} }
  child { node[concept color=poppurpleL, text=popdark, align=center] {paire complète\\ à l'écrit} }
};
\end{tikzpicture}
\legende{Carte mentale récapitulative : les trois structures de la leçon et leurs schémas.}
\end{popfigure}

\begin{aretenir}[L'essentiel de la section]
\begin{itemize}
\item \zh{怎么了} = « qu'est-ce qu'il y a ? » ; « comment vas-tu » = \zh{你好吗}.
\item \zh{怎么} + verbe = comment (manière) ; \zh{怎么} + \zh{不/没} = pourquoi (raison surprise).
\item Un sujet : Sujet + \zh{不但}…\zh{而且}… Deux sujets : \zh{不但} + S1…\zh{而且} + S2…
\item À l'écrit, la paire \zh{因为……所以……} est complète ; le français, lui, n'en garde qu'un.
\item Aucune inversion, aucun déplacement : les mots-outils chinois ont une place fixe.
\item Avant de rendre la copie : sujet ? paire ? \zh{了} ? place de \zh{怎么} ?
\end{itemize}
\end{aretenir}

\begin{exercice}[À vous de jouer]
Corrigez les phrases suivantes, toutes extraites de copies, puis traduisez-les en français :
\begin{enumerate}
\item \zh{你怎么了去雅温得？}
\item \zh{因为我很喜欢足球，所以我所以每天看比赛。}
\item \zh{不但喀麦隆的咖啡很好，而且中国的茶也很好。} (cette phrase est-elle fautive ? justifiez)
\item \zh{他怎么了不来学校？}
\end{enumerate}
\end{exercice}

\begin{corrige}[Exercice « À vous de jouer »]
\begin{enumerate}
\item \zh{你怎么去雅温得？} \emph{Nǐ zěnme qù Yǎwēndé?} « Comment vas-tu à Yaoundé ? » — \zh{了} n'a rien à faire ici : on demande la manière, pas un problème.
\item \zh{因为我很喜欢足球，所以每天看比赛。} ou mieux : \zh{因为我很喜欢足球，所以我每天看比赛。} — un seul \zh{所以} ! « Comme j'aime beaucoup le football, je regarde des matchs tous les jours. »
\item Phrase \textbf{correcte} : il y a deux sujets (\zh{喀麦隆的咖啡} / \zh{中国的茶}), donc \zh{不但} en tête et le second sujet après \zh{而且} : « Non seulement le café du Cameroun est très bon, mais le thé de Chine aussi. »
\item \zh{他怎么不来学校？} \emph{Tā zěnme bù lái xuéxiào?} « Pourquoi ne vient-il pas à l'école ? » — on interroge sur la raison (\zh{怎么} + négation), donc pas de \zh{了} ; \zh{怎么了} demanderait « qu'est-ce qu'il a ».
\end{enumerate}
\end{corrige}

\coursec{Entraînement guidé et transformation}

Après avoir comparé le chinois et le français et repéré les pièges les plus fréquents (ordre des mots, oubli du mot-outil, double emploi de \zh{因为} et \zh{所以}), il est temps de passer à la pratique. Cette dernière étape de la leçon suit une progression rigoureuse : de la simple reconnaissance des structures jusqu'à la production libre, en passant par la complétion, la transformation et la traduction. C'est exactement le cheminement exigé à l'épreuve de chinois du Baccalauréat.

\begin{popfigure}
\begin{tikzpicture}[node distance=0.35cm]
\node[fill=popblueL, rounded corners, minimum width=2.5cm, minimum height=1.1cm, align=center, font=\small\bfseries] (a) {1. Reconnaissance\\identifier};
\node[fill=popgreenL, rounded corners, minimum width=2.5cm, minimum height=1.1cm, align=center, font=\small\bfseries, right=of a] (b) {2. Complétion\\remplir les trous};
\node[fill=popgoldL, rounded corners, minimum width=2.5cm, minimum height=1.1cm, align=center, font=\small\bfseries, right=of b] (c) {3. Transformation\\relier, reformuler};
\node[fill=poporangeL, rounded corners, minimum width=2.5cm, minimum height=1.1cm, align=center, font=\small\bfseries, right=of c] (d) {4. Traduction\\thème et version};
\node[fill=poppinkL, rounded corners, minimum width=2.5cm, minimum height=1.1cm, align=center, font=\small\bfseries, right=of d] (e) {5. Production\\rédiger};
\draw[-{Stealth}, thick, popblue] (a) -- (b);
\draw[-{Stealth}, thick, popblue] (b) -- (c);
\draw[-{Stealth}, thick, popblue] (c) -- (d);
\draw[-{Stealth}, thick, popblue] (d) -- (e);
\end{tikzpicture}
\legende{Frise de progression des exercices : de la reconnaissance à la production libre.}
\end{popfigure}

\courssub{Exercice 1 : compléter avec 怎么了, 怎么 ou 为什么}

\begin{methode}[Choisir entre 怎么了, 怎么 et 为什么]
\begin{enumerate}
\item \textbf{怎么了} (zěnme le) : s'emploie seul ou après un nom/pronom pour demander « qu'est-ce qu'il y a ? », face à une situation anormale. Structure : \cle{Nom + 怎么了？}
\item \textbf{怎么} (zěnme) : précède un \textbf{verbe} pour demander la manière, la cause ou l'état : « comment ? », « pourquoi (se fait-il que) ? ». Structure : \cle{Sujet + 怎么 + Verbe ?}
\item \textbf{为什么} (wèishénme) : demande la raison, la cause ; attend une réponse avec \zh{因为}. Structure : \cle{Sujet + 为什么 + Verbe/Adjectif ?}
\end{enumerate}
\end{methode}

\begin{exoresolu}[Complétion guidée]
Complétez avec \zh{怎么了}, \zh{怎么} ou \zh{为什么} : \\
1. 你\_\_\_\_？脸色不太好。\\
2. 他今天\_\_\_\_没来上课？\\
3. 你\_\_\_\_不去看医生？
\tcblower
\textbf{Solution détaillée :}\\
1. 你怎么了？脸色不太好。 \emph{Nǐ zěnme le? Liǎnsè bù tài hǎo.} « Qu'est-ce que tu as ? Tu n'as pas bonne mine. » — On s'inquiète d'une situation anormale (le visage), donc \zh{怎么了}.\\
2. 他今天怎么没来上课？ \emph{Tā jīntiān zěnme méi lái shàngkè?} « Comment se fait-il qu'il ne soit pas venu en cours aujourd'hui ? » — \zh{怎么} + verbe \zh{来} exprime l'étonnement devant un fait.\\
3. 你为什么不去看医生？ \emph{Nǐ wèishénme bù qù kàn yīshēng?} « Pourquoi ne vas-tu pas voir le médecin ? » — On demande une raison ; la réponse naturelle serait \zh{因为……}.
\end{exoresolu}

\begin{exercice}[À vous]
Complétez avec \zh{怎么了}, \zh{怎么} ou \zh{为什么} :\\
1. 玛丽\_\_\_\_？她哭了。\\
2. 杜阿拉的空气\_\_\_\_这么差？\\
3. 你\_\_\_\_每天喝很多水？\\
4. 老师\_\_\_\_？他好像很累。\\
5. 你\_\_\_\_知道这家医院在哪儿？
\end{exercice}

\begin{corrige}[Exercice 1]
1. 玛丽怎么了？她哭了。 \emph{Mǎlì zěnme le? Tā kū le.} — situation anormale (elle pleure).\\
2. 杜阿拉的空气为什么这么差？ \emph{Dù'ālā de kōngqì wèishénme zhème chà?} — demande de raison, réponse attendue avec \zh{因为}.\\
3. 你为什么每天喝很多水？ \emph{Nǐ wèishénme měi tiān hē hěn duō shuǐ?} — demande de raison.\\
4. 老师怎么了？他好像很累。 \emph{Lǎoshī zěnme le? Tā hǎoxiàng hěn lèi.} — inquiétude face à un état.\\
5. 你怎么知道这家医院在哪儿？ \emph{Nǐ zěnme zhīdào zhè jiā yīyuàn zài nǎr?} — \zh{怎么} + verbe \zh{知道} : « comment sais-tu… ».
\end{corrige}

\courssub{Exercice 2 : relier avec 不但……而且……}

\begin{propriete}[Rappel de la règle]
La corrélation \cle{不但……而且……} (búdàn……érqiě……) exprime l'addition : « non seulement… mais aussi… ».\\
\textbf{Même sujet} : Sujet + 不但 + A + 而且 + B.\\
\textbf{Sujets différents} : 不但 + Sujet1 + A + 而且 + Sujet2 + B.\\
Le mot-outil \zh{而且} est \textbf{obligatoire} dans la seconde partie ; on ne peut pas l'omettre comme « aussi » en français. Le renfort \zh{也} (aussi) est fréquent devant le verbe de la seconde proposition si les sujets diffèrent.
\end{propriete}

\begin{exemplebox}[Modèles à imiter]
\zh{这种药不但便宜，而且有效。} \emph{Zhè zhǒng yào búdàn piányi, érqiě yǒuxiào.} « Ce médicament est non seulement bon marché, mais aussi efficace. » (même sujet)\\
\zh{不但医生来了，而且护士也来了。} \emph{Búdàn yīshēng lái le, érqiě hùshi yě lái le.} « Non seulement le médecin est venu, mais l'infirmière est venue aussi. » (sujets différents ; notez le \zh{也} de renfort)
\end{exemplebox}

\begin{exercice}[Reliez les deux phrases]
1. 这家医院很大。这家医院很干净。\\
2. 他会说法语。他会说汉语。\\
3. 病人吃药了。他的家人来看他了。
\end{exercice}

\begin{corrige}[Exercice 2]
1. 这家医院不但很大，而且很干净。 \emph{Zhè jiā yīyuàn búdàn hěn dà, érqiě hěn gānjìng.} — même sujet : \zh{不但} après le sujet.\\
2. 他不但会说法语，而且会说汉语。 \emph{Tā búdàn huì shuō Fǎyǔ, érqiě huì shuō Hànyǔ.}\\
3. 不但病人吃药了，而且他的家人也来看他了。 \emph{Búdàn bìngrén chī yào le, érqiě tā de jiārén yě lái kàn tā le.} — sujets différents : \zh{不但} avant le premier sujet, \zh{也} ajouté.
\end{corrige}

\courssub{Exercice 3 : transformer en 因为……所以……}

\begin{methode}[Transformer une phrase de cause]
\begin{enumerate}
\item Repérez la \textbf{cause} et la \textbf{conséquence} dans la phrase donnée.
\item Placez \zh{因为} devant la cause et \zh{所以} devant la conséquence : \cle{因为 + cause，所以 + conséquence。}
\item Vérifiez que l'ordre cause → conséquence est respecté (en chinois, la cause vient presque toujours en premier).
\end{enumerate}
\end{methode}

\begin{exoresolu}[Transformation guidée]
Transformez : 空气污染很严重。很多人咳嗽了。
\tcblower
Cause : la pollution de l'air est grave. Conséquence : beaucoup de gens toussent.\\
因为空气污染很严重，所以很多人咳嗽了。\\
\emph{Yīnwèi kōngqì wūrǎn hěn yánzhòng, suǒyǐ hěn duō rén késou le.}\\
« Comme la pollution de l'air est très grave, beaucoup de gens se sont mis à tousser. »
\end{exoresolu}

\begin{exercice}[Transformez]
1. 他生病了。他没来学校。\\
2. 雅温得的交通太堵了。我们迟到了。\\
3. 这种水果很有营养。孩子们应该多吃。
\end{exercice}

\begin{corrige}[Exercice 3]
1. 因为他生病了，所以他没来学校。 \emph{Yīnwèi tā shēngbìng le, suǒyǐ tā méi lái xuéxiào.}\\
2. 因为雅温得的交通太堵了，所以我们迟到了。 \emph{Yīnwèi Yǎwēndé de jiāotōng tài dǔ le, suǒyǐ wǒmen chídào le.}\\
3. 因为这种水果很有营养，所以孩子们应该多吃。 \emph{Yīnwèi zhè zhǒng shuǐguǒ hěn yǒu yíngyǎng, suǒyǐ háizimen yīnggāi duō chī.}
\end{corrige}

\courssub{Exercice 4 : thème (français → chinois)}

\begin{methode}[Réussir le thème au Bac]
\begin{enumerate}
\item \textbf{Traduisez d'abord la structure} : identifiez la charpente (怎么了 ? 不但……而且…… ? 因为……所以…… ?) et posez-la sur le papier.
\item \textbf{Insérez ensuite le vocabulaire} : placez les mots dans les emplacements prévus par la structure, en respectant l'ordre Sujet – Temps/Lieu – Verbe.
\item \textbf{Vérifiez enfin les tons et les caractères} : relisez chaque caractère (clé, nombre de traits), chaque pinyin mental, et chassez les mots-outils français en trop (« parce que… donc… » n'existe pas en chinois : on garde les deux, mais jamais \zh{因为} seul suivi d'une virgule sans suite logique).
\end{enumerate}
\end{methode}

\begin{exemplebox}[Corrigé type avec barème indicatif]
Phrase à traduire : « Comme il a couru sous la pluie, il est tombé malade. »\\
\textbf{Corrigé} : 因为他淋雨了，所以生病了。 \emph{Yīnwèi tā lín yǔ le, suǒyǐ shēngbìng le.}\\
\textbf{Barème indicatif (sur 3 pts)} : mot-outil correctement placé (\zh{因为}……\zh{所以}) : 1 pt ; ordre des mots et structure : 1 pt ; caractères exacts : 1 pt.\\
Autres phrases du thème :\\
1. « Qu'est-ce que tu as ? Tu tousses beaucoup. » → 你怎么了？你咳嗽得很厉害。 \emph{Nǐ zěnme le? Nǐ késou de hěn lìhai.}\\
2. « Le médecin lui a non seulement donné des médicaments, mais aussi conseillé de se reposer. » → 医生不但给他药，而且让他多休息。 \emph{Yīshēng búdàn gěi tā yào, érqiě ràng tā duō xiūxi.}\\
3. « Pourquoi l'eau de ce quartier est-elle sale ? » → 这个区的水为什么很脏？ \emph{Zhège qū de shuǐ wèishénme hěn zāng?}\\
4. « Comme il fait très chaud, les enfants boivent beaucoup d'eau. » → 因为天气很热，所以孩子们喝很多水。 \emph{Yīnwèi tiānqì hěn rè, suǒyǐ háizimen hē hěn duō shuǐ.}
\end{exemplebox}

\begin{attention}[Pièges du thème]
Ne traduisez jamais « non seulement… mais aussi » par deux phrases séparées : la corrélation \zh{不但……而且……} doit apparaître \textbf{dans la même phrase}. De même, « parce que… donc… » se traduit par \zh{因为……所以……} complet ; en français on n'utilise qu'un seul des deux mots, mais en chinois les deux ensemble sont la marque de bon niveau.
\end{attention}

\courssub{Exercice 5 : version (chinois → français)}

\begin{textebox}[À l'hôpital de Yaoundé]
因为空气污染很严重，所以很多人生病了。\\
\emph{Yīnwèi kōngqì wūrǎn hěn yánzhòng, suǒyǐ hěn duō rén shēngbìng le.}\\
医生不但给他们药，而且告诉他们要多休息。\\
\emph{Yīshēng búdàn gěi tāmen yào, érqiě gàosu tāmen yào duō xiūxi.}\\
一个病人问：“医生，我怎么了？”\\
\emph{Yí ge bìngrén wèn: “Yīshēng, wǒ zěnme le?”}\\
医生说：“你怎么不戴口罩？因为空气太脏了，所以你要小心。”\\
\emph{Yīshēng shuō: “Nǐ zěnme bù dài kǒuzhào? Yīnwèi kōngqì tài zāng le, suǒyǐ nǐ yào xiǎoxīn.”}
\end{textebox}

\begin{exercice}[Traduisez le texte ci-dessus en français]
Soulignez d'abord les trois structures étudiées, puis traduisez phrase par phrase.
\end{exercice}

\begin{corrige}[Exercice 5]
« Comme la pollution de l'air est très grave, beaucoup de gens sont tombés malades. Le médecin leur a non seulement donné des médicaments, mais leur a aussi dit de bien se reposer. Un malade demande : “Docteur, qu'est-ce que j'ai ?” Le médecin répond : “Pourquoi ne portez-vous pas de masque ? Comme l'air est très sale, vous devez faire attention.” »
\end{corrige}

\courssub{Exercice 6 : corriger les phrases fautives}

\begin{exoresolu}[Correction guidée]
Corrigez : 不但他很累，而且很饿。(même sujet)
\tcblower
Erreur : avec un même sujet, \zh{不但} se place \textbf{après} le sujet, pas avant.\\
Correction : 他不但很累，而且很饿。 \emph{Tā búdàn hěn lèi, érqiě hěn è.} « Il est non seulement fatigué, mais aussi affamé. »
\end{exoresolu}

\begin{exercice}[Corrigez les phrases]
1. 怎么了你？\\
2. 因为他病了，他没来。\\
3. 她不但会唱歌，会跳舞。\\
4. 为什么你不喝水因为天气很热？
\end{exercice}

\begin{corrige}[Exercice 6]
1. 你怎么了？ — le sujet \zh{你} précède \zh{怎么了}.\\
2. 因为他病了，所以他没来。 — avec \zh{因为}, il faut \zh{所以} dans la seconde partie (ou supprimer \zh{因为}).\\
3. 她不但会唱歌，而且会跳舞。 — le mot-outil \zh{而且} est obligatoire.\\
4. Phrase fautive : \zh{为什么你不喝水因为天气很热？} Erreur : ordre des mots incorrect (la cause précède la conséquence) et mélange structure interrogative (\zh{为什么}) et structure causative (\zh{因为}).\\
Corrections possibles : \zh{因为天气很热，所以你不喝水吗？} (question sur la conséquence) ou \zh{天气很热，你为什么不喝水？} (constat + question).
\end{corrige}

\courssub{Production libre guidée : la santé au Cameroun}

\begin{experience}[Rédaction en quatre phrases]
Rédigez un court paragraphe intitulé « 喀麦隆的健康 » (\emph{Kāmàilóng de jiànkāng}, « La santé au Cameroun ») contenant obligatoirement :\\
1. une question avec \zh{怎么了} ou \zh{怎么} ;\\
2. une phrase avec \zh{不但……而且……} ;\\
3. une phrase avec \zh{因为……所以……} ;\\
4. une phrase de conclusion libre.\\
\textbf{Exemple de production attendue} :\\
在喀麦隆，很多人怎么了？他们生病了。因为有些地方的空气和水不干净，所以人们容易生病。政府不但建了新的医院，而且请了很多医生。我们应该多喝水，多休息，保护健康。\\
\emph{Zài Kāmàilóng, hěn duō rén zěnme le? Tāmen shēngbìng le. Yīnwèi yǒuxiē dìfang de kōngqì hé shuǐ bù gānjìng, suǒyǐ rénmen róngyì shēngbìng. Zhèngfǔ búdàn jiàn le xīn de yīyuàn, érqiě qǐng le hěn duō yīshēng. Wǒmen yīnggāi duō hē shuǐ, duō xiūxi, bǎohù jiànkāng.}\\
« Au Cameroun, qu'arrive-t-il à beaucoup de gens ? Ils tombent malades. Comme dans certains endroits l'air et l'eau ne sont pas propres, les gens tombent facilement malades. Le gouvernement a non seulement construit de nouveaux hôpitaux, mais aussi recruté beaucoup de médecins. Nous devons boire beaucoup d'eau, bien nous reposer et protéger notre santé. »
\end{experience}

\begin{center}
\begin{tabularx}{\linewidth}{|X|X|X|X|}
\hline
\rowcolor{popblueL} Caractères & Pinyin & Nature & Traduction \\
\hline
咳嗽 & késou & verbe & tousser \\
\hline
口罩 & kǒuzhào & nom & masque (de protection) \\
\hline
营养 & yíngyǎng & nom & nutrition, éléments nutritifs \\
\hline
休息 & xiūxi & verbe & se reposer \\
\hline
政府 & zhèngfǔ & nom & gouvernement \\
\hline
保护 & bǎohù & verbe & protéger \\
\hline
容易 & róngyì & adjectif & facile, facilement \\
\hline
\end{tabularx}
\end{center}

\begin{savaistu}[Les corrélations au Bac]
Au Baccalauréat camerounais, les corrélations \zh{不但……而且……} et \zh{因为……所以……} rapportent des points précis de « maîtrise des structures » dans la grille officielle de correction : un candidat qui emploie correctement une corrélation complète dans sa production écrite est mieux noté qu'un candidat qui juxtapose des phrases simples. Entraînez-vous donc à placer au moins une corrélation dans chaque rédaction !
\end{savaistu}

\begin{aretenir}[Bilan de l'entraînement]
\begin{itemize}
\item \zh{怎么了} s'inquiète d'une situation ; \zh{怎么} + verbe s'étonne ou demande la manière ; \zh{为什么} demande la raison.
\item Même sujet : Sujet + \zh{不但}……\zh{而且}…… ; sujets différents : \zh{不但} + Sujet1……\zh{而且} + Sujet2……
\item \zh{因为} + cause，\zh{所以} + conséquence : la cause d'abord, toujours.
\item Au thème : structure d'abord, vocabulaire ensuite, vérification des caractères enfin.
\end{itemize}
\end{aretenir}

\newpage

\coursec{Activité d'intégration}

\coursec{Leçon 13 — Grammaire : Interrogation avec 怎么了 et 怎么 ; 不但……而且…… ; 因为……所以……}

\courssub{Objectifs de la leçon}

À l'issue de cette leçon, tu seras capable de :
\begin{itemize}
\item mobiliser dans une communication authentique, orale puis écrite, les trois structures de la leçon : l'interrogation avec \zh{怎么了} et \zh{怎么}, la corrélation \zh{不但……而且……} et la corrélation \zh{因为……所以……} ;
\item utiliser le lexique de la santé (symptômes, maladies, traitements) dans une situation de consultation médicale ;
\item interroger un patient, décrire des symptômes, expliquer une cause et proposer un traitement en chinois correct ;
\item rédiger un compte rendu écrit court (6 à 8 phrases) respectant l'ordre des mots chinois et la place exacte des mots-outils ;
\item évaluer la production de tes camarades à l'aide d'une grille de critères simples (exactitude des structures, prononciation des tons, richesse du vocabulaire).
\end{itemize}

\courssub{Observation et découverte}

\courssub{Situation de communication}

Tu es élève en Terminale A au lycée de Yaoundé. Dans le cadre d'un jumelage avec une école de Pékin, ton établissement organise une « journée santé ». Un médecin chinois volontaire, le docteur Wáng, anime un atelier de sensibilisation. Comme il ne parle ni français ni anglais, les élèves sinisants sont sollicités : certains jouent le rôle de patients, d'autres celui de médecin, d'autres encore celui d'interprète qui traduit en français pour le public.

Voici le modèle de dialogue affiché au tableau par ton professeur. Lis-le attentivement avant de commencer.

\begin{textebox}[Modèle de consultation médicale]
\textbf{医生：}你怎么了？\\
\emph{Yīshēng : Nǐ zěnme le ?}\\
« Médecin : Qu'est-ce que tu as ? »\\[4pt]
\textbf{病人：}我不但头疼，而且发烧，嗓子也很疼。\\
\emph{Bìngrén : Wǒ bùdàn tóu téng, érqiě fāshāo, sǎngzi yě hěn téng.}\\
« Patient : Non seulement j'ai mal à la tête, mais en plus j'ai de la fièvre, et j'ai aussi mal à la gorge. »\\[4pt]
\textbf{医生：}你怎么会这样？你昨天晚上做了什么？\\
\emph{Yīshēng : Nǐ zěnme huì zhèyàng ? Nǐ zuótiān wǎnshang zuò le shénme ?}\\
« Médecin : Comment se fait-il que tu sois comme ça ? Qu'as-tu fait hier soir ? »\\[4pt]
\textbf{病人：}因为昨天下大雨，我没有带伞，所以感冒了。\\
\emph{Bìngrén : Yīnwèi zuótiān xià dà yǔ, wǒ méiyǒu dài sǎn, suǒyǐ gǎnmào le.}\\
« Patient : Parce qu'hier il a beaucoup plu et que je n'avais pas de parapluie, j'ai attrapé un rhume. »\\[4pt]
\textbf{医生：}因为你太累了，而且着凉了，所以你要多喝水，多休息，还要吃这个药。\\
\emph{Yīshēng : Yīnwèi nǐ tài lèi le, érqiě zháoliáng le, suǒyǐ nǐ yào duō hē shuǐ, duō xiūxi, hái yào chī zhège yào.}\\
« Médecin : Parce que tu es trop fatigué et que tu as pris froid, tu dois boire beaucoup d'eau, bien te reposer, et aussi prendre ce médicament. »\\[4pt]
\textbf{病人：}谢谢医生！\\
\emph{Bìngrén : Xièxie yīshēng !}\\
« Patient : Merci, docteur ! »
\end{textebox}

\begin{attention}[Observer avant d'agir]
Repère dans le dialogue : (1) la question d'ouverture \zh{你怎么了？} ; (2) la question sur la cause avec \zh{怎么} seul : \zh{你怎么会这样？} ; (3) la description des symptômes avec \zh{不但……而且……} ; (4) l'explication de la cause et la prescription avec \zh{因为……所以……}. Ce sont exactement les quatre moments que tu devras reproduire.
\end{attention}

\begin{center}
\begin{tabularx}{\linewidth}{|X|X|X|X|}
\hline
\rowcolor{popblueL} Caractères & Pinyin & Nature & Traduction \\
\hline
医生 & yīshēng & nom & médecin \\
\hline
病人 & bìngrén & nom & patient, malade \\
\hline
头疼 & tóu téng & loc. verb. & avoir mal à la tête \\
\hline
发烧 & fāshāo & verbe & avoir de la fièvre \\
\hline
感冒 & gǎnmào & verbe / nom & attraper un rhume ; rhume \\
\hline
嗓子疼 & sǎngzi téng & loc. verb. & avoir mal à la gorge \\
\hline
咳嗽 & késou & verbe & tousser \\
\hline
吃药 & chī yào & loc. verb. & prendre un médicament \\
\hline
休息 & xiūxi & verbe & se reposer \\
\hline
下雨 & xià yǔ & loc. verb. & pleuvoir \\
\hline
带伞 & dài sǎn & loc. verb. & apporter un parapluie \\
\hline
累 & lèi & adjectif & fatigué \\
\hline
着凉 & zháoliáng & verbe & prendre froid, attraper froid \\
\hline
医院 & yīyuàn & nom & hôpital \\
\hline
检查 & jiǎnchá & verbe & examiner, faire un examen \\
\hline
疟疾 & nüèjí & nom & paludisme \\
\hline
力气 & lìqi & nom & force, énergie \\
\hline
住院 & zhùyuàn & verbe & être hospitalisé \\
\hline
\end{tabularx}
\end{center}

\courssub{Interroger avec 怎么了 et 怎么}

\begin{propriete}[Structure 1 : 怎么了 vs 怎么]
\textbf{A. \zh{怎么了} (zěnme le) — Question ouverte sur l'état ou l'événement}\\
\textbf{Schéma :} Sujet + \zh{怎么了} ？\\
\textbf{Emploi :} Demander « Qu'est-ce qui se passe ? », « Qu'est-ce que tu as ? », « Quel est le problème ? ». La particule \zh{了} marque un \textbf{changement d'état} (survenue d'un symptôme, d'un incident), non le passé.\\
\textbf{Exemples :}
\begin{itemize}
\item \zh{你怎么了？} \emph{Nǐ zěnme le ?} « Qu'est-ce que tu as ? / Qu'est-ce qui t'arrive ? »
\item \zh{他怎么了？} \emph{Tā zěnme le ?} « Qu'est-ce qu'il a ? »
\item \zh{手机怎么了？} \emph{Shǒujī zěnme le ?} « Qu'est-ce qu'il y a avec le téléphone ? »
\end{itemize}

\textbf{B. \zh{怎么} (zěnme) + Verbe — Question sur la manière, la cause ou le procédé}\\
\textbf{Schéma :} Sujet + \zh{怎么} + Verbe (+ Complément) ?\\
\textbf{Emploi :} Demander « Comment ? », « Pourquoi ? », « De quelle manière ? ». \zh{怎么} précède toujours le verbe.\\
\textbf{Exemples :}
\begin{itemize}
\item \zh{你怎么知道？} \emph{Nǐ zěnme zhīdào ?} « Comment le sais-tu ? » (manière)
\item \zh{你怎么会这样？} \emph{Nǐ zěnme huì zhèyàng ?} « Comment se fait-il que tu sois comme ça ? » (cause)
\item \zh{怎么去医院？} \emph{Zěnme qù yīyuàn ?} « Comment aller à l'hôpital ? » (procéder)
\item \zh{你怎么不来了？} \emph{Nǐ zěnme bù lái le ?} « Pourquoi ne viens-tu plus ? » (cause + changement d'état)
\end{itemize}

\textbf{Distinction clé :} \zh{怎么了} = question globale sur l'état (réponse : nom/adjectif/phrase). \zh{怎么}+V = question sur le processus/la cause (réponse : phrase explicative).
\end{propriete}

\begin{exemplebox}[Entraînement oral : Transformer les questions]
Transforme les questions ci-dessous en utilisant \zh{怎么了} ou \zh{怎么}+V selon le sens visé.
\begin{enumerate}
\item Tu vois un camarade qui boite. Demande : « Qu'est-ce qui t'arrive ? » $\rightarrow$ \zh{你怎么了？}
\item Tu veux savoir comment il s'est blessé. Demande : « Comment tu t'es fait mal ? » $\rightarrow$ \zh{你怎么受伤的？} (\zh{受伤} shòushāng : se blesser)
\item Le médecin demande la cause de la fièvre. $\rightarrow$ \zh{你怎么发烧了？}
\item Tu trouves la porte fermée. Demande : « Qu'est-ce qu'il y a avec la porte ? » $\rightarrow$ \zh{门怎么了？}
\end{enumerate}
\end{exemplebox}

\courssub{La corrélation 不但……而且……}

\begin{propriete}[Structure 2 : 不但……而且…… (bùdàn… érqiě…)]
\textbf{Schéma :} Sujet + \zh{不但} + Proposition 1，\zh{而且} + Proposition 2.\\
\textbf{Sens :} « Non seulement… mais encore… », « Non seulement… mais en plus… ». La deuxième proposition \textbf{renforce} ou \textbf{ajoute} un élément plus important, plus grave ou plus surprenant que le premier.\\
\textbf{Règle d'or (sujet commun) :} Si le sujet est identique dans les deux propositions, il se place \textbf{avant \zh{不但}} et n'est \textbf{pas répété} devant \zh{而且}.\\
\textbf{Exemples :}
\begin{itemize}
\item \zh{我不但头疼，而且发烧。} \emph{Wǒ bùdàn tóu téng, érqiě fāshāo.} (Sujet \zh{我} unique)
\item \zh{这道菜不但好吃，而且不贵。} \emph{Zhège cài bùdàn hǎochī, érqiě bù guì.} « Ce plat n'est pas seulement bon, en plus il n'est pas cher. »
\item \zh{不但下雨，而且刮风。} \emph{Bùdàn xià yǔ, érqiě guā fēng.} (Sujet nul / phénomène météo : pas de sujet explicite)
\end{itemize}

\textbf{Variante courante :} \zh{不但……还/也……} (moins formel, « non seulement… mais aussi… »). Ex. \zh{我不但头疼，还咳嗽。}
\end{propriete}

\begin{attention}[Piège francophone : Place du sujet]
\emph{Erreur :} \zh{不但我头疼，而且我发烧。} (Calque du français « Non seulement moi j'ai mal à la tête… »)
\emph{Correct :} \zh{我不但头疼，而且发烧。}
Le sujet \zh{我} « gouverne » les deux verbes. En chinois, l'économie syntaxique impose l'antéposition du sujet commun.
\end{attention}

\courssub{La corrélation 因为……所以……}

\begin{propriete}[Structure 3 : 因为……所以…… (yīnwèi… suǒyǐ…)]
\textbf{Schéma :} \zh{因为} + Cause (proposition complète)，\zh{所以} + Conséquence (proposition complète).\\
\textbf{Sens :} « Parce que… donc… », « Comme… par conséquent… ». Établit un lien logique \textbf{cause $\rightarrow$ conséquence}.\\
\textbf{Particularité cruciale vs Français :} En français, « parce que » et « donc » sont redondants (on dit « Parce qu'il pleut, je ne sors pas » \textbf{sans} « donc »). En chinois standard, \zh{因为} et \zh{所以} \textbf{s'utilisent ensemble} pour marquer explicitement la relation logique. L'omission de \zh{所以} est possible à l'oral familier, mais incorrecte à l'écrit et au niveau Terminale.\\
\textbf{Ordre des mots :} La clause \zh{因为} précède généralement la clause \zh{所以}. L'inversion (\zh{所以}… \zh{因为}…) est possible pour l'insistance sur le résultat mais plus rare.\\
\textbf{Exemples :}
\begin{itemize}
\item \zh{因为下雨了，所以我不去学校。} \emph{Yīnwèi xià yǔ le, suǒyǐ wǒ bú qù xuéxiào.}
\item \zh{因为他很努力，所以他考试通过了。} \emph{Yīnwèi tā hěn nǔlì, suǒyǐ tā kǎoshì tōngguò le.}
\item \zh{因为昨天下大雨，而且没带伞，所以我感冒了。} \emph{Yīnwèi zuótiān xià dà yǔ, érqiě méi dài sǎn, suǒyǐ wǒ gǎnmào le.} (Cause complexe avec \zh{而且})
\end{itemize}
\end{propriete}

\begin{exemplebox}[Cas particulier : Sujets différents]
Si la cause et la conséquence ont des sujets différents, chaque sujet reste dans sa proposition respective.
\zh{因为天气冷，所以河水结冰了。} \emph{Yīnwèi tiānqì lěng, suǒyǐ héshuǐ jiébīng le.} « Parce qu'il fait froid, l'eau de la rivière a gelé. » (Sujet 1 : \zh{天气}, Sujet 2 : \zh{河水})
\end{exemplebox}

\courssub{Comparaison avec le français et pièges fréquents}

\begin{attention}[Synthèse des erreurs fréquentes des francophones]
\begin{itemize}
\item \textbf{Ordre dans 怎么了 :} \cle{Jamais} \zh{你什么了？} $\rightarrow$ Toujours \zh{你怎么了？} (\zh{怎么} avant \zh{了}).
\item \textbf{Oubli de \zh{所以} :} \cle{Ne supprimez pas \zh{所以}} par habitude du français. \zh{因为下雨了，我不去。} $\leftarrow$ Phrase incomplète à l'écrit. Correct : \zh{因为下雨了，所以我不去。}
\item \textbf{Sujet répété dans \zh{不但……而且……} :} \zh{不但我头疼，而且我发烧。} $\leftarrow$ Lourd/incorrect. Correct : \zh{我不但头疼，而且发烧。}
\item \textbf{Nature de \zh{了} dans \zh{怎么了} :} Ce \zh{了} n'est \textbf{pas} le marqueur d'accompli (passé). C'est une particule modale de \textbf{changement d'état} / survenue nouvelle : « Qu'est-ce qui t'est arrivé (maintenant) ? ».
\item \textbf{Confusion \zh{怎么} / \zh{什么} :} \zh{怎么} = « comment / pourquoi » (adverbe, devant verbe). \zh{什么} = « quoi / quel » (pronom/adjectif, devant nom ou après verbe). Ex. \zh{你怎么不吃？} (Pourquoi) vs \zh{你不吃什么？} (Qu'est-ce que tu ne manges pas ?).
\end{itemize}
\end{attention}

\begin{center}
\begin{tabularx}{\linewidth}{|X|X|X|}
\hline
\rowcolor{popblueL} Structure chinoise & Équivalent français (structure) & Piège à éviter \\
\hline
\zh{你怎么了？} & Qu'est-ce que tu as ? & \zh{*你什么了？} \\
\hline
\zh{你怎么知道？} & Comment le sais-tu ? & \zh{*你什么知道？} \\
\hline
\zh{我不但累，而且饿。} & Non seulement je suis fatigué, mais j'ai faim. & \zh{*不但我累，而且我饿。} \\
\hline
\zh{因为下雨，所以我不去。} & Comme il pleut, je n'y vais pas. & \zh{*因为下雨，我不去。} (écrit) \\
\hline
\end{tabularx}
\end{center}

\courssub{Entraînement guidé et transformation}

\courssub{Activité intégrée : Simulation de consultation médicale (2 h)}

\begin{experience}[Consignes de l'atelier]
Travail par groupes de trois : un \textbf{médecin}, un \textbf{patient}, un \textbf{interprète}. Durée totale : 2 heures.

\begin{enumerate}
\item \textbf{Compréhension du modèle (10 min).} Lis le dialogue de la section 1 à voix haute dans ton groupe. Souligne les trois structures de la leçon et vérifie que chacun comprend la traduction. L'interprète s'entraîne à traduire chaque réplique en français.
\item \textbf{Choix du cas médical (5 min).} Choisissez une situation parmi :
\begin{itemize}
\item Le paludisme après la saison des pluies à Douala (\zh{疟疾} nüèjí, \zh{蚊子} wénzi).
\item Une intoxication alimentaire après un repas au marché (\zh{食物中毒} shíwù zhòngdú, \zh{肚子疼} dùzi téng, \zh{呕吐} ǒutù).
\item Une fatigue intense avant les examens du baccalauréat (\zh{太累} tài lèi, \zh{失眠} shīmián, \zh{压力大} yālì dà).
\item Une blessure au pied après un match de football (\zh{受伤} shòushāng, \zh{脚疼} jiǎo téng, \zh{骨折} gǔzhé).
\end{itemize}
\item \textbf{Préparation du jeu de rôle (20 min).} Rédigez ensemble un dialogue de 8 à 12 répliques. Contraintes obligatoires :
\begin{itemize}
\item Le médecin pose au moins deux questions : une avec \zh{怎么了} et une avec \zh{怎么} seul.
\item Le patient décrit ses symptômes avec \zh{不但……而且……}.
\item Le médecin explique la cause et donne le traitement avec \zh{因为……所以……}.
\end{itemize}
\item \textbf{Jeu de rôle devant la classe (5 min par groupe).} Jouez la consultation. L'interprète traduit chaque réplique en français juste après sa prononciation en chinois. Attention aux tons, surtout dans \zh{tóu téng} (2\textsuperscript{e} ton + 2\textsuperscript{e} ton $\rightarrow$ sandhi 3\textsuperscript{e}+2\textsuperscript{e} en pratique mais notation 2+2), \zh{fāshāo} (1\textsuperscript{er} + 1\textsuperscript{er}), \zh{zháoliáng} (2\textsuperscript{e} + 2\textsuperscript{e}), \zh{nüèjí} (4\textsuperscript{e} + 2\textsuperscript{e}).
\item \textbf{Rédaction du compte rendu (25 min).} Chaque groupe rédige par écrit le compte rendu de la consultation : 6 à 8 phrases en chinois, à la troisième personne (ex. \zh{病人不但头疼，而且发烧。}). Chacune des trois structures doit apparaître au moins une fois. Ajoutez le pinyin sous chaque phrase et la traduction française.
\item \textbf{Évaluation par les pairs (10 min).} Échangez votre compte rendu avec un autre groupe et remplissez la grille :
\begin{itemize}
\item Exactitude des trois structures (2 points chacune = 6 pts).
\item Prononciation des tons pendant le jeu de rôle (2 pts).
\item Richesse du vocabulaire santé (2 pts).
\end{itemize}
Notez une erreur à corriger et un point fort.
\end{enumerate}
\end{experience}

\courssub{Proposition de production et corrigé commenté}

\begin{exoresolu}[Compte rendu modèle : Consultation pour paludisme à Douala]
Énoncé : Rédige le compte rendu écrit (6 à 8 phrases) de la consultation jouée par ton groupe, en utilisant au moins une fois chaque structure de la leçon.
\tcblower
\textbf{Production modèle :}\\[4pt]
病人叫玛丽，她是杜阿拉的学生。\\
\emph{Bìngrén jiào Mǎlì, tā shì Dù'ālā de xuésheng.}\\
« La patiente s'appelle Marie, elle est élève à Douala. »\\[4pt]
医生问她：« 你怎么了？»\\
\emph{Yīshēng wèn tā : « Nǐ zěnme le ? »}\\
« Le médecin lui demande : "Qu'est-ce que tu as ?" »\\[4pt]
她说她不但发烧，而且头疼，也没有力气。\\
\emph{Tā shuō tā bùdàn fāshāo, érqiě tóu téng, yě méiyǒu lìqi.}\\
« Elle dit que non seulement elle a de la fièvre, mais en plus elle a mal à la tête et n'a pas de forces. »\\[4pt]
医生问她怎么病了。\\
\emph{Yīshēng wèn tā zěnme bìng le.}\\
« Le médecin lui demande comment elle est tombée malade. »\\[4pt]
因为杜阿拉最近天天下雨，蚊子很多，所以她得了疟疾。\\
\emph{Yīnwèi Dù'ālā zuìjìn tiāntiān xià yǔ, wénzi hěn duō, suǒyǐ tā dé le nüèjí.}\\
« Parce qu'à Douala il pleut tous les jours en ce moment et qu'il y a beaucoup de moustiques, elle a attrapé le paludisme. »\\[4pt]
因为她不但发烧，而且很累，所以医生让她住院检查。\\
\emph{Yīnwèi tā bùdàn fāshāo, érqiě hěn lèi, suǒyǐ yīshēng ràng tā zhùyuàn jiǎnchá.}\\
« Parce que non seulement elle a de la fièvre mais aussi elle est très fatiguée, le médecin lui demande d'être hospitalisée pour un examen. »\\[4pt]
她要吃药，多喝水，多休息。\\
\emph{Tā yào chī yào, duō hē shuǐ, duō xiūxi.}\\
« Elle doit prendre des médicaments, boire beaucoup d'eau et bien se reposer. »\\[4pt]
玛丽谢谢医生，就回家了。\\
\emph{Mǎlì xièxie yīshēng, jiù huí jiā le.}\\
« Marie remercie le médecin et rentre à la maison. »\\[8pt]
\textbf{Commentaire des choix de langue :}
\begin{itemize}
\item \textbf{Structure 1} : La question \zh{你怎么了？} est rapportée directement, comme dans le dialogue réel ; la seconde question utilise \zh{怎么} seul (\zh{怎么病了} « comment est-elle tombée malade »), ce qui montre la différence entre les deux emplois.
\item \textbf{Structure 2} : \zh{不但发烧，而且头疼} — le sujet \zh{她} est placé une seule fois avant \zh{不但}, conformément à la règle du sujet unique ; la virgule chinoise ， sépare les deux propositions. L'ajout de \zh{也没有力气} montre l'extension naturelle par \zh{也}.
\item \textbf{Structure 3} : Deux occurrences de \zh{因为……所以……}, dont une qui combine les trois structures dans la même phrase (phrase 6 : \zh{因为} [cause avec \zh{不但……而且……}] \zh{所以} [conséquence]) — c'est le niveau de complexité syntaxique attendu en Terminale A.
\item \textbf{Lexique santé} : \zh{发烧}, \zh{头疼}, \zh{力气}, \zh{疟疾} (paludisme), \zh{住院} (être hospitalisé), \zh{检查} : le contexte camerounais (Douala, pluies, moustiques) rend le vocabulaire mémorable et ancré dans la réalité des élèves.
\item \textbf{Particule \zh{了}} : Dans \zh{得了疟疾} et \zh{回家了}, \zh{了} marque l'accomplissement (aspect perfectif) ; dans \zh{怎么了} et \zh{病了}, il marque le changement d'état / survenue — ne pas les confondre.
\end{itemize}
\end{exoresolu}

\courssub{Bilan de l'activité}

Cette activité t'a permis de passer de la règle à l'usage réel : interroger avec \zh{怎么了} et \zh{怎么} pour ouvrir et orienter un dialogue, enrichir une description avec \zh{不但……而且……}, et construire un raisonnement logique (cause $\rightarrow$ conséquence) avec \zh{因为……所以……}. Tu as également mobilisé le lexique de la santé dans un contexte camerounais authentique (Douala, saison des pluies, paludisme) et pratiqué la médiation linguistique (traduire du chinois vers le français), une compétence précieuse pour le baccalauréat et au-delà.

\begin{aretenir}[Auto-évaluation : Je sais...]
\begin{itemize}
\item poser la question \zh{你怎么了？} et varier avec \zh{怎么} + verbe pour demander la cause ou la manière ;
\item décrire deux faits liés avec \zh{不但……而且……} en plaçant correctement le sujet commun avant \zh{不但} ;
\item expliquer une cause et sa conséquence avec \zh{因为……所以……} sans calquer la structure française (conserver \zh{所以}) ;
\item rédiger un court compte rendu de 6 à 8 phrases en respectant l'ordre des mots chinois et la ponctuation ， 。 ;
\item repérer et corriger les erreurs de structures (ordre, mots-outils) dans la production d'un camarade.
\end{itemize}
\end{aretenir}

\newpage

\coursec{Méthodes et exercices résolus}

\coursec{Leçon 13 — Grammaire : Interrogation avec 怎么了 et 怎么 ; 不但……而且…… ; 因为……所以……}

\begin{textebox}[Dialogue au dispensaire du lycée de Yaoundé]
\zh{护士：同学，你怎么了？} (Hùshì: Tóngxué, nǐ zěnme le?) \\
\zh{学生：护士阿姨，我头疼，而且喉咙痛。} (Xuésheng: Hùshì āyí, wǒ tóu téng, érqiě hóulóng tòng.) \\
\zh{护士：因为天气变化大，所以很容易感冒。你不但要吃药，而且要多喝水。} (Hùshì: Yīnwèi tiānqì biànhuà dà, suǒyǐ hěn róngyì gǎnmào. Nǐ búdàn yào chī yào, érqiě yào duō hē shuǐ.) \\
\zh{学生：怎么吃药？饭前还是饭后？} (Xuésheng: Zěnme chī yào? Fànqián háishi fànhòu?) \\
\zh{护士：饭后吃。好了，去休息吧。} (Hùshì: Fànhòu chī. Hǎo le, qù xiūxi ba.)
\tcblower
\textbf{Traduction :} \\
Infirmière : Élève, qu'est-ce que tu as ? \\
Élève : Tante infirmière, j'ai mal à la tête et en plus mal à la gorge. \\
Infirmière : Parce que le temps change beaucoup, il est facile d'attraper un rhume. Tu dois non seulement prendre des médicaments, mais aussi boire beaucoup d'eau. \\
Élève : Comment prendre le médicament ? Avant ou après le repas ? \\
Infirmière : Après le repas. Bon, va te reposer.
\end{textebox}

\begin{center}
\begin{tabularx}{\linewidth}{|X|X|X|X|}
\hline
\rowcolor{popblueL} \textbf{Caractères} & \textbf{Pinyin} & \textbf{Nature} & \textbf{Traduction} \\ \hline
\zh{怎么了} & zěnme le & locution & qu'est-ce qu'il y a ? / qu'est-ce qui ne va pas ? \\ \hline
\zh{怎么} & zěnme & adv. & comment / pourquoi (familier) \\ \hline
\zh{不但} & búdàn & conj. & non seulement \\ \hline
\zh{而且} & érqiě & conj. & mais aussi / de plus \\ \hline
\zh{因为} & yīnwèi & conj. & parce que \\ \hline
\zh{所以} & suǒyǐ & conj. & donc / par conséquent \\ \hline
\zh{头疼} & tóuténg & VO / adj. & avoir mal à la tête \\ \hline
\zh{喉咙痛} & hóulóng tòng & VO / adj. & avoir mal à la gorge \\ \hline
\zh{感冒} & gǎnmào & v. & attraper un rhume / rhume \\ \hline
\zh{发烧} & fāshāo & v. & avoir de la fièvre \\ \hline
\zh{吃药} & chī yào & VO & prendre des médicaments \\ \hline
\zh{喝水} & hē shuǐ & VO & boire de l'eau \\ \hline
\zh{环境} & huánjìng & n. & environnement \\ \hline
\zh{污染} & wūrǎn & v./n. & polluer / pollution \\ \hline
\zh{锻炼} & duànliàn & v. & faire de l'exercice, s'entraîner \\ \hline
\zh{医院} & yīyuàn & n. & hôpital \\ \hline
\zh{医生} & yīshēng & n. & médecin \\ \hline
\zh{护士} & hùshì & n. & infirmier(ère) \\ \hline
\zh{天气} & tiānqì & n. & temps, météo \\ \hline
\zh{变化} & biànhuà & v./n. & changer / changement \\ \hline
\zh{容易} & róngyì & adj. & facile (à) \\ \hline
\zh{饭前} & fànqián & n. & avant le repas \\ \hline
\zh{饭后} & fànhòu & n. & après le repas \\ \hline
\end{tabularx}
\end{center}

\courssub{Interroger avec 怎么了 et 怎么}

\begin{propriete}[Deux fonctions distinctes de 怎么]
\zh{怎么} (zěnme) est un mot interrogatif polyvalent. Son sens dépend de sa position et du contexte :
\begin{itemize}
\item \textbf{Manière / Moyen} (« comment ») : placé \textbf{avant le verbe} (ou l'adjectif). Ex: \zh{怎么去？} (Zěnme qù? « Comment y aller ? »).
\item \textbf{Cause / Raison} (« pourquoi », registre familier/oral) : placé \textbf{avant le verbe/l'adjectif}, souvent avec une nuance de reproche ou d'étonnement. Ex: \zh{你怎么不吃？} (Nǐ zěnme bù chī? « Pourquoi tu ne manges pas ? »).
\item \textbf{État / Problème} avec la particule \textbf{了} : \zh{怎么了} (zěnme le) = « Qu'est-ce qu'il y a ? / Qu'est-ce qui ne va pas ? ». Se place \textbf{après le sujet}.
\end{itemize}
\end{propriete}

\begin{methode}[Construire une question avec 怎么了 et 怎么]
\begin{enumerate}
\item \textbf{Identifier l'intention :} Constat d'anomalie (santé, humeur, objet cassé) $\rightarrow$ \zh{怎么了}. Demande de procédure $\rightarrow$ \zh{怎么} (manière). Demande de raison (oral) $\rightarrow$ \zh{怎么} (cause).
\item \textbf{Placer le mot interrogatif :}
\begin{itemize}
\item Sujet + \zh{怎么了} ? (Pas de verbe nécessaire). Ex: \zh{手机怎么了？} (Shǒujī zěnme le? « Qu'est-ce qu'a le téléphone ? »).
\item Sujet + \zh{怎么} + Verbe (+ Complément) ? Ex: \zh{这个中文怎么读？} (Zhège Zhōngwén zěnme dú? « Comment lit-on ce chinois ? »).
\item Sujet + \zh{怎么} + \zh{没/不} + Verbe ? Ex: \zh{你怎么没来？} (Nǐ zěnme méi lái? « Pourquoi n'es-tu pas venu ? »).
\end{itemize}
\item \textbf{Interdiction formelle :} Ne jamais ajouter la particule \zh{吗} à une phrase contenant déjà un mot interrogatif (\zh{怎么}, \zh{什么}, \zh{谁}, \zh{哪}...).
\item \textbf{Répondre :} Reprendre le sujet et donner l'information demandée (manière, raison, ou description de l'état). Ex: \zh{我感冒了。} (Wǒ gǎnmào le.)
\end{enumerate}
\end{methode}

\begin{exemplebox}[Exemples variés : 怎么了 vs 怎么]
\begin{itemize}
\item \zh{妹妹怎么了？} (Mèimei zěnme le?) — « Qu'est-ce qu'a ma petite sœur ? » (Constat : elle pleure).
\item \zh{这个字怎么写？} (Zhège zì zěnme xiě?) — « Comment s'écrit ce caractère ? » (Manière).
\item \zh{你怎么迟到了？} (Nǐ zěnme chídào le?) — « Pourquoi as-tu été en retard ? » (Cause, reproche léger).
\item \zh{怎么走？} (Zěnme zǒu?) — « Comment y aller ? / Quel chemin prendre ? » (Sujet omis = nous/vous).
\end{itemize}
\end{exemplebox}

\begin{exoresolu}[Application : Au dispensaire et en ville]
Énoncé — Traduisez en chinois (caractères + pinyin) :
1. Qu'est-ce qui ne va pas ? (Tu t'adresses à un ami).
2. Comment allez-vous à l'hôpital central de Yaoundé ?
3. Pourquoi n'as-tu pas fait tes devoirs ? (Registre familier).
4. Répondez à la question 1 : « J'ai mal au ventre. ».
\tcblower
\textbf{Raisonnement.}
1. Constat d'un problème $\rightarrow$ \zh{怎么了} après le sujet \zh{你}.
2. Demande de moyen/itinéraire $\rightarrow$ \zh{怎么} avant le verbe \zh{去}, complément de lieu \zh{雅温得中心医院} après le verbe.
3. Demande de cause (action non réalisée dans le passé) $\rightarrow$ \zh{怎么} + \zh{没} + V.
4. Réponse : Sujet + Verbe d'état + \zh{了} (changement d'état / apparition du symptôme).

\textbf{Réponses :}
1. \zh{你怎么了？} (Nǐ zěnme le?)
2. \zh{你们怎么去雅温得中心医院？} (Nǐmen zěnme qù Yǎwēndé Zhōngxīn Yīyuàn?)
3. \zh{你怎么没做作业？} (Nǐ zěnme méi zuò zuòyè?)
4. \zh{我肚子疼了。} (Wǒ dùzi téng le.) \emph{Note :} \zh{肚子疼} (dùzi téng) = mal au ventre.
\end{exoresolu}

\courssub{La corrélation 不但……而且…… (Gradation et addition)}

\begin{definition}[不但……而且……]
\zh{不但……而且……} (búdàn……érqiě……) est une corrélation coordonnante exprimant une \textbf{relation de gradation} : la deuxième proposition apporte une information \textbf{plus forte, plus importante ou supplémentaire} par rapport à la première. Équivalent français : « Non seulement… mais (aussi)… ».
\end{definition}

\begin{propriete}[Schémas structurels]
Deux cas de figure selon l'identité des sujets :
\begin{enumerate}
\item \textbf{Même sujet (cas le plus fréquent) :}
\zh{Sujet + 不但 + Predicat 1, 而且 + Predicat 2.}
\emph{Le sujet n'est pas répété.} Les prédicats peuvent être verbaux, adjectivaux ou nominaux (avec \zh{是}).
\item \textbf{Sujets différents :}
\zh{不但 + Sujet 1 + Predicat 1, 而且 + Sujet 2 + Predicat 2.}
\emph{Chaque sujet se place immédiatement après son corrélatif.}
\end{enumerate}
\textbf{Ponctuation :} Une virgule chinoise \zh{，} sépare obligatoirement les deux propositions.
\end{propriete}

\begin{methode}[Construire une phrase avec 不但……而且……]
\begin{enumerate}
\item \textbf{Vérifier la logique :} Y a-t-il une progression (gradation) ? « Il est grand, \textbf{en plus} il est fort » = OK. « Il est grand, \textbf{mais} il est petit » = NON (utiliser \zh{虽然……但是……}).
\item \textbf{Identifier les sujets :} Même sujet $\rightarrow$ Sujet avant \zh{不但}. Sujets différents $\rightarrow$ Sujet 1 après \zh{不但}, Sujet 2 après \zh{而且}.
\item \textbf{Assembler :} Conserver les mots-outils (négation \zh{不/没}, aspect \zh{了/过/着}, modaux \zh{会/能/要}) à l'intérieur de chaque proposition.
\item \textbf{Enrichir (optionnel) :} Ajouter \zh{也} (yě) ou \zh{还} (hái) après \zh{而且} pour renforcer l'addition : \zh{而且还……}, \zh{而且也……}.
\end{enumerate}
\end{methode}

\begin{exemplebox}[Sujet unique vs Sujets différents]
\begin{itemize}
\item \textbf{Même sujet :} \zh{这道菜不但好看，而且很好吃。} (Zhè dào cài búdàn hǎokàn, érqiě hěn hǎochī.) « Ce plat n'est pas seulement beau, il est en plus délicieux. » (Adjectifs prédicatifs, pas de \zh{是}).
\item \textbf{Même sujet + verbes :} \zh{他不但会说汉语，而且会写汉字。} (Tā búdàn huì shuō Hànyǔ, érqiě huì xiě Hànzì.)
\item \textbf{Sujets différents :} \zh{不但老师去北京，而且学生也去。} (Búdàn lǎoshī qù Běijīng, érqiě xuésheng yě qù.) « Non seulement le prof va à Pékin, mais les élèves y vont aussi. »
\item \textbf{Avec 也/还 :} \zh{我哥哥不但踢足球，而且还打篮球。} (Wǒ gēge búdàn tī zúqiú, érqiě hái dǎ lánqiú.)
\end{itemize}
\end{exemplebox}

\begin{exoresolu}[Transformation : Le sport à Douala]
Énoncé — 1. Combinez : \zh{他喜欢跑步。他喜欢游泳。} (Même sujet, addition).
2. Traduisez : « Non seulement ma sœur étudie le chinois, mais elle l'enseigne aussi. »
3. Combinez (sujets différents) : \zh{空气脏。水也脏。} (L'air est sale. L'eau aussi est sale.)
\tcblower
\textbf{Réponses :}
1. \zh{他不但喜欢跑步，而且喜欢游泳。} (Tā búdàn xǐhuan pǎobù, érqiě xǐhuan yóuyǒng.) \emph{Répéter le verbe \zh{喜欢} est plus sûr et plus clair.}
2. \zh{我妹妹不但学汉语，而且也教汉语。} (Wǒ mèimei búdàn xué Hànyǔ, érqiě yě jiāo Hànyǔ.) \emph{Sujet unique \zh{我妹妹}.}
3. \zh{不但空气脏，而且水也脏。} (Búdàn kōngqì zāng, érqiě shuǐ yě zāng.) \emph{Sujets différents (\zh{空气}, \zh{水}) placés après les corrélatifs.}
\end{exoresolu}

\courssub{La corrélation 因为……所以…… (Cause et conséquence)}

\begin{definition}[因为……所以……]
\zh{因为……所以……} (yīnwèi……suǒyǐ……) est une corrélation subordonnante marquant explicitement le lien de \textbf{causalité}.
\begin{itemize}
\item \zh{因为} (yīnwèi) introduit la \textbf{cause} (la raison).
\item \zh{所以} (suǒyǐ) introduit la \textbf{conséquence} (le résultat).
\end{itemize}
\end{definition}

\begin{propriete}[Différence majeure avec le français]
En français, l'énoncé « Parce qu'il pleut, \textbf{donc} je reste chez moi » est considéré comme une \textbf{faute de syntaxe} (redondance / anacoluthie). On choisit \emph{soit} « Parce qu'il pleut, je reste… » \emph{soit} « Il pleut, \textbf{donc} je reste… ».
\medskip
\cle{En chinois, la double marque 因为……所以…… est la norme standard, correcte et élégante.} Elle structure le raisonnement. L'omission de \zh{所以} est possible à l'oral ou dans l'écrit très fluide (\zh{因为下雨，我不去。}), mais \textbf{à l'écrit scolaire et au Bac, on garde les deux} pour garantir la clarté logique.
\end{propriete}

\begin{methode}[Employer 因为……所以……]
\begin{enumerate}
\item \textbf{Séparer Cause et Conséquence :} Identifier les deux événements et leur lien logique.
\item \textbf{Ordonner :} Cause d'abord (\zh{因为}...), Conséquence ensuite (\zh{所以}...). L'ordre inverse (\zh{所以}..., \zh{因为}...) est rare et marqué (insistance sur le résultat).
\item \textbf{Placer les sujets :}
\begin{itemize}
\item Sujets différents : \zh{因为 + S1 + V1, 所以 + S2 + V2.}
\item Même sujet : On peut le placer \textbf{avant \zh{因为}} (thème de la phrase) : \zh{S + 因为 + V1, 所以 + V2.} Ou le répéter dans chaque proposition.
\end{itemize}
\item \textbf{Gérer la négation :} Passé négatif $\rightarrow$ \zh{没(有)}. Présent/Futur/Habitude $\rightarrow$ \zh{不}.
\end{enumerate}
\end{methode}

\begin{exemplebox}[Variantes de position du sujet]
\begin{itemize}
\item \textbf{Sujets différents :} \zh{因为下雨了，所以比赛取消了。} (Yīnwèi xià yǔ le, suǒyǐ bǐsài qǔxiāo le.) « Parce qu'il a plu, le match a été annulé. »
\item \textbf{Même sujet (devant 因为) :} \zh{我因为生病了，所以没去上课。} (Wǒ yīnwèi shēngbìng le, suǒyǐ méi qù shàngkè.) « Moi, parce que j'étais malade, je ne suis pas allé en cours. » (Topicalisation).
\item \textbf{Même sujet (répété) :} \zh{因为我生病了，所以我没去上课。} (Yīnwèi wǒ shēngbìng le, suǒyǐ wǒ méi qù shàngkè.) Plus lourd mais clair.
\item \textbf{Sans 所以 (oral) :} \zh{因为太贵了，我没买。} (Yīnwèi tài guì le, wǒ méi mǎi.)
\end{itemize}
\end{exemplebox}

\begin{exoresolu}[Cause, conséquence et santé]
Énoncé — 1. Traduisez : « Parce que l'air de la ville est pollué, beaucoup de gens tombent malades. »
2. Complétez : \zh{\_\_\_\_\_\_ 他每天锻炼, \_\_\_\_\_\_ 他身体很好。}
3. Corrigez l'erreur : \zh{因为他感冒了，他没来学校。} (Phrase correcte mais incomplète pour l'examen).
\tcblower
\textbf{Réponses :}
1. \zh{因为城市的空气被污染了，所以很多人生病了。} (Yīnwèi chéngshì de kōngqì bèi wūrǎn le, suǒyǐ hěn duō rén shēngbìng le.) \emph{Vocabulaire module : \zh{污染} (pollution), \zh{生病} (tomber malade).}
2. \zh{因为他每天锻炼，所以他身体很好。} (Yīnwèi tā měitiān duànliàn, suǒyǐ tā shēntǐ hěn hǎo.) \emph{Attention : \zh{每天} appelle souvent \zh{都} (dōu) : \zh{因为他每天都锻炼……}}
3. Version attendue au Bac : \zh{因为他感冒了，\textbf{所以}他没来学校。} (Yīnwèi tā gǎnmào le, \textbf{suǒyǐ} tā méi lái xuéxiào.) \emph{Négliger \zh{所以} fait perdre le point de structure.}
\end{exoresolu}

\courssub{Comparaison avec le français et pièges fréquents}

\begin{attention}[Les 5 erreurs qui coûtent des points au Bac]
\begin{enumerate}
\item \textbf{Ajouter \zh{吗} à une question avec \zh{怎么} / \zh{怎么了} :} \zh{✗ 你怎么了吗？} $\rightarrow$ \zh{✅ 你怎么了？} Règle générale : \textbf{Mot interrogatif = Pas de \zh{吗}}.
\item \textbf{Supprimer \zh{所以} par réflexe francophone :} \zh{✗ 因为下雨了，我不去。} (Acceptable à l'oral, \textbf{pénalisé à l'écrit formel}). $\rightarrow$ \zh{✅ 因为下雨了，所以我不去。} Le chinois \textbf{exige} souvent les deux bornes.
\item \textbf{Mal placer le sujet avec \zh{不但} (sujets différents) :} \zh{✗ 我不但喜欢，而且他也喜欢。} (Sens : « Moi non seulement j'aime, mais lui aussi » — structure bancale). $\rightarrow$ \zh{✅ 不但我喜欢，而且他也喜欢。} (Sujet 1 après \zh{不但}, Sujet 2 après \zh{而且}).
\item \textbf{Confondre caractères et tons proches :}
\begin{itemize}
\item \zh{怎} (zěn, clé \zh{心} en bas) vs \zh{什} (shén/shí, clé \zh{人} à gauche). \zh{怎么} s'écrit avec \zh{怎}.
\item \zh{而} (ér, 2\textsuperscript{e} ton) vs \zh{儿} (ér, suffixe / enfant). \zh{而且} s'écrit \zh{而}.
\item \zh{为} (wèi, 4\textsuperscript{e} ton dans \zh{因为}) vs \zh{为} (wéi, 2\textsuperscript{e} ton dans \zh{为了} wèile / \zh{作为} zuòwéi).
\end{itemize}
\item \textbf{Oublier \zh{没} pour le passé négatif dans une conséquence :} \zh{✗ 因为他病了，所以他不来学校。} (\zh{不来} = présent/habitude « il ne vient pas »). $\rightarrow$ \zh{✅ 因为他病了，所以他\textbf{没}来学校。} (\zh{没来} = passé « il n'est pas venu »).
\end{enumerate}
\end{attention}

\begin{center}
\begin{tabularx}{\linewidth}{|X|X|X|}
\hline
\rowcolor{poporangeL} \textbf{Structure} & \textbf{Français (Faute vs Correct)} & \textbf{Chinois (Standard)} \\ \hline
Interrogation « Pourquoi » & Pourquoi n'es-tu pas venu ? & \zh{你怎么没来？} (Pas de \zh{吗}) \\ \hline
Interrogation « Qu'est-ce qui ne va pas ? » & Qu'est-ce qu'il a ? & \zh{他怎么了？} (Sujet + 怎么了) \\ \hline
Gradation (Même sujet) & Non seulement il court, mais il nage. & \zh{他不但跑步，而且游泳。} (Sujet avant \zh{不但}) \\ \hline
Gradation (Sujets diff.) & Non seulement moi, mais lui aussi. & \zh{不但我去，而且他也去。} (Sujets après corrélatifs) \\ \hline
Cause/Conséquence & \textbf{Faute :} Parce que... donc... & \zh{因为……所以……} (Les deux obligatoires) \\ \hline
Négation passé (Conséquence) & Il n'est PAS venu (hier). & \zh{他没来} (Pas \zh{不来}) \\ \hline
\end{tabularx}
\end{center}

\courssub{Entraînement guidé : Transformation et combinaison}

\begin{methode}[Transformer des phrases simples en phrases complexes]
\begin{enumerate}
\item \textbf{Analyser la relation logique :} Addition/Gradation $\rightarrow$ \zh{不但……而且……} ; Cause/Effet $\rightarrow$ \zh{因为……所以……} ; Constat d'état $\rightarrow$ Question \zh{怎么了}.
\item \textbf{Déterminer les sujets :} Identiques ou différents ? $\rightarrow$ Place du sujet.
\item \textbf{Fusionner :} Supprimer le sujet répété (cas même sujet). Conserver \zh{了}, \zh{没}, \zh{过}, modaux.
\item \textbf{Ponctuer :} Virgule \zh{，} entre les deux membres de la corrélation. Point \zh{。} final.
\item \textbf{Relire :} Vérifier l'ordre \zh{不但}... \zh{而且} / \zh{因为}... \zh{所以}. Vérifier l'absence de \zh{吗}.
\end{enumerate}
\end{methode}

\begin{exoresolu}[Atelier de transformation complet]
Énoncé — Transformez selon la structure indiquée :
1. \zh{这种药很贵。这种药很有效。} $\rightarrow$ \zh{不但……而且……}
2. \zh{他感冒了。他没来学校。} $\rightarrow$ \zh{因为……所以……}
3. \zh{你看起来很累。} $\rightarrow$ Question avec \zh{怎么了}
4. \zh{我哥哥会做中国菜。我弟弟也会做。} $\rightarrow$ \zh{不但……而且……} (Sujets différents)
5. \zh{今天天气很热。我们不去爬山。} $\rightarrow$ \zh{因为……所以……}
\tcblower
\textbf{Réponses :}
1. \zh{这种药不但很贵，而且很有效。} (Zhè zhǒng yào búdàn hěn guì, érqiě hěn yǒuxiào.) \emph{Classificateur \zh{种} conservé. Adjectifs prédicatifs.}
2. \zh{因为他感冒了，所以他没来学校。} (Yīnwèi tā gǎnmào le, suǒyǐ tā méi lái xuéxiào.) \emph{\zh{没来} passé.}
3. \zh{你怎么了？} (Nǐ zěnme le?) \emph{Constat visuel $\rightarrow$ 怎么了.}
4. \zh{不但我哥哥会做中国菜，而且我弟弟也会做。} (Búdàn wǒ gēge huì zuò Zhōngguó cài, érqiě wǒ dìdi yě huì zuò.) \emph{Sujets différents.}
5. \zh{因为今天天气很热，所以我们不去爬山。} (Yīnwèi jīntiān tiānqì hěn rè, suǒyǐ wǒmen bù qù páshān.) \emph{\zh{不去} présent/futur proche.}
\end{exoresolu}

\courssub{Expression écrite : Rédiger un paragraphe cohérent}

\begin{methode}[Rédiger un court texte mobilisant les 3 structures]
\begin{enumerate}
\item \textbf{Planifier (4-5 phrases) :}
\begin{itemize}
\item Phrase 1 : Situation initiale (Contexte : temps, lieu, personnage).
\item Phrase 2 : Question / Constat avec \zh{怎么了} (discours direct ou indirect).
\item Phrase 3 : Explication causale avec \zh{因为……所以……}.
\item Phrase 4 : Gradation / Aggravation avec \zh{不但……而且……}.
\item Phrase 5 : Conclusion / Souhait / Conseil.
\end{itemize}
\item \textbf{Varier les sujets} mais contrôler la syntaxe de chaque structure.
\item \textbf{Mobiliser le lexique du Module 2 :} \zh{环境, 空气, 污染, 生病, 发烧, 头疼, 医生, 药, 锻炼, 休息}.
\item \textbf{Relire (Check-list) :} Tons (pinyin) ? Caractères (clés, traits) ? Ordre des mots (Sujet / Corrélatifs / Verbe / Complément) ? Ponctuation chinoise ? Pas de \zh{吗} superflu ?
\end{enumerate}
\end{methode}

\begin{exoresolu}[Rédaction guidée : « Mon ami et la pollution »]
Énoncé — Rédigez un texte de 5 phrases intitulé \zh{环境与健康} (Environnement et santé) en utilisant \textbf{au moins une fois} : \zh{怎么了}, \zh{因为……所以……}, \zh{不但……而且……}.
\tcblower
\textbf{Texte proposé :}
\zh{昨天下午，我的朋友小李没来打篮球。我问他：“小李，你怎么了？”他说：“因为昨天空气污染很严重，所以我生病了。”他不但咳嗽，而且发烧了。医生说他要多喝水，好好休息。我希望环境变好，大家都健康。}
\\
(Zuótiān xiàwǔ, wǒ de péngyou Xiǎo Lǐ méi lái dǎ lánqiú. Wǒ wèn tā: « Xiǎo Lǐ, nǐ zěnme le? » Tā shuō: « Yīnwèi zuótiān kōngqì wūrǎn hěn yánzhòng, suǒyǐ wǒ shēngbìng le. » Tā búdàn késòu, érqiě fāshāo le. Yīshēng shuō tā yào duō hē shuǐ, hǎohǎo xiūxi. Wǒ xīwàng huánjìng biàn hǎo, dàjiā dōu jiànkāng.)

\textbf{Traduction :} Hier après-midi, mon ami Xiao Li n'est pas venu jouer au basket. Je lui ai demandé : « Xiao Li, qu'est-ce que tu as ? » Il a dit : « Parce qu'hier la pollution de l'air était grave, je suis tombé malade. » Il tousse non seulement, mais en plus il a de la fièvre. Le médecin a dit qu'il doit boire beaucoup d'eau et bien se reposer. J'espère que l'environnement s'améliorera et que tout le monde sera en bonne santé.

\textbf{Analyse grammaticale :}
\begin{itemize}
\item Phrase 2 : \zh{你怎么了？} (Sujet + 怎么了, discours direct).
\item Phrase 3 : \zh{因为……所以……} complet, sujet \zh{我} répété (clair), \zh{生病了} (changement d'état).
\item Phrase 4 : \zh{不但……而且……}, sujet unique \zh{他} (omis dans la 2\textsuperscript{e} proposition), verbes \zh{咳嗽} (késòu, tousser) et \zh{发烧}.
\item Vocabulaire Module 2 : \zh{空气污染, 严重, 生病, 咳嗽, 发烧, 医生, 健康}.
\end{itemize}
\end{exoresolu}

\begin{savaistu}[Santé et Environnement : Regards croisés Cameroun - Chine]
\begin{itemize}
\item \textbf{Cameroun :} Les grandes villes (Douala, Yaoundé) font face à la pollution de l'air (poussière harmattan, émissions véhicules, brûlage déchets) et de l'eau. Le paludisme (\zh{疟疾} nüèjí) reste endémique. Le système de santé combine médecine moderne (hôpitaux \zh{医院}, centres de santé) et médecine traditionnelle (\zh{中医} zhōngyī très prisée, plantes médicinales).
\item \textbf{Chine :} Lutte massive contre la pollution de l'air (\zh{雾霾} wùmái = smog/brume polluante) via restriction circulation, énergies propres. Système de santé universel (\zh{医保} yībǎo) très développé. Médecine traditionnelle chinoise (MTC : \zh{针灸} zhēnjiǔ acupuncture, \zh{拔罐} báguàn ventouses, pharmacopée) intégrée aux hôpitaux modernes.
\item \textbf{Point commun :} La prise de conscience écologique (\zh{环保} huánbǎo) grandit chez les jeunes des deux pays (tri déchets, vélos, végétarisme). Le sport (\zh{锻炼} duànliàn) est promu comme prévention santé n°1.
\end{itemize}

\end{savaistu}

\begin{experience}[Jeu de rôle : « Au dispensaire » / Débat : « Pollution et santé »]
\begin{enumerate}
\item \textbf{Binôme A (Infirmier) / B (Élève malade) :} B entre, A demande \zh{怎么了？}. B décrit symptômes (\zh{头疼, 发烧, 咳嗽}). A diagnostique : \zh{因为……所以……} (ex: \zh{因为天气变化，所以感冒}). A prescrit : \zh{不但……而且……} (ex: \zh{不但吃药，而且多喝水}). Échanger les rôles.
\item \textbf{Débat classe entière (Terminale) :} « \zh{保护环境就是保护健康} » (Protéger l'environnement, c'est protéger la santé). Deux camps : \emph{Gouvernement/Entreprises} vs \emph{Citoyens/ONG}. Utiliser obligatoirement \zh{因为……所以……} pour argumenter les causes/effets et \zh{不但……而且……} pour lister les mesures. Ex: \zh{因为工厂排放废气，所以空气脏。政府不但要罚款，而且要关厂。}
\end{enumerate}

\end{experience}

\begin{exercice}[Entraînement personnel — Devoir / Examen blanc]
\textbf{Exercice 1 : Questions (Traduction vers le chinois)}
1. Qu'est-ce qui t'arrive ? (Tu vois un ami pâle).
2. Comment écrirez-vous ce nouveau caractère ?
3. Pourquoi n'ont-ils pas nettoyé la salle de classe ? (Registre familier).
4. Pourquoi est-ce que tu ne manges pas de viande ? (Registre standard \zh{为什么} vs familier \zh{怎么} — donnez les deux formes).

\textbf{Exercice 2 : Transformation \zh{不但……而且……}}
Combinez chaque paire de phrases. Attention aux sujets.
a) \zh{北京很大。北京很美。}
b) \zh{我喜欢吃水饺。我妹妹也喜欢吃水饺。}
c) \zh{这个学生很聪明。这个学生很用功。} (\zh{用功} yònggōng = travailleur/assidu).

\textbf{Exercice 3 : Transformation \zh{因为……所以……}}
Combinez. Gardez les deux corrélatifs. Gérez la négation.
a) \zh{昨天下大雨。比赛取消了。} (\zh{取消} qǔxiāo = annuler).
b) \zh{他没复习。他考试没考好。} (\zh{复习} fùxí = réviser ; \zh{考好} kǎohǎo = réussir un examen).
c) \zh{空气污染严重。很多人得了呼吸道疾病。} (\zh{呼吸道疾病} hūxīdào jíbìng = maladies respiratoires).

\textbf{Exercice 4 : Correction d'erreurs (Corrigez et expliquez l'erreur)}
a) \zh{你怎么了吗？}
b) \zh{因为他累了，他不来参加聚会。} (Contexte : hier soir).
c) \zh{我不但喜欢喝茶，而且咖啡也喜欢。} (Ordre des mots dans 2\textsuperscript{e} proposition).
d) \zh{不但老师来了，学生来了。} (Manque \zh{而且}).

\textbf{Exercice 5 : Expression écrite (Type Bac)}
Sujet : \zh{我的健康生活} (Ma vie saine). Rédigez un paragraphe de 60-80 caractères chinois.
Contraintes : Utiliser \textbf{au moins} 1 fois \zh{怎么了} (discours direct ou rapporté), 1 fois \zh{因为……所以……}, 1 fois \zh{不但……而且……}. Vocabulaire : \zh{锻炼, 吃水果, 睡觉, 早睡早起, 身体好}.
\end{exercice}

\begin{corrige}[Exercices 1 à 5]
\textbf{Exercice 1 :}
1. \zh{你怎么了？} (Nǐ zěnme le?)
2. \zh{这个新字怎么写？} (Zhège xīn zì zěnme xiě?)
3. \zh{他们怎么没打扫教室？} (Tāmen zěnme méi dǎsǎo jiàoshì?)
4. Familier : \zh{你怎么不吃肉？} (Nǐ zěnme bù chī ròu?) / Standard : \zh{你为什么不吃肉？} (Nǐ wèishénme bù chī ròu?)

\textbf{Exercice 2 :}
a) \zh{北京不但很大，而且很美。} (Běijīng búdàn hěn dà, érqiě hěn měi.) — Sujet unique.
b) \zh{不但我喜欢吃水饺，而且我妹妹也喜欢吃水饺。} (Búdàn wǒ xǐhuan chī shuǐjiǎo, érqiě wǒ mèimei yě xǐhuan chī shuǐjiǎo.) — Sujets différents.
c) \zh{这个学生不但很聪明，而且很用功。} (Zhège xuésheng búdàn hěn cōngming, érqiě hěn yònggōng.) — Sujet unique, adjectifs.

\textbf{Exercice 3 :}
a) \zh{因为昨天下大雨了，所以比赛取消了。} (Yīnwèi zuótiān xià dà yǔ le, suǒyǐ bǐsài qǔxiāo le.)
b) \zh{因为他没复习，所以他考试没考好。} (Yīnwèi tā méi fùxí, suǒyǐ tā kǎoshì méi kǎohǎo.) — Double négation \zh{没} conservée.
c) \zh{因为空气污染很严重，所以很多人得了呼吸道疾病。} (Yīnwèi kōngqì wūrǎn hěn yánzhòng, suǒyǐ hěn duō rén déle hūxīdào jíbìng.) — \zh{得了} (déle) = contracter (maladie), aspect accompli.

\textbf{Exercice 4 :}
a) \zh{✗ 你怎么了吗？} $\rightarrow$ \zh{✅ 你怎么了？} \emph{Erreur : \zh{吗} interdit avec mot interrogatif.}
b) \zh{✗ ...所以他不来...} $\rightarrow$ \zh{✅ ...所以他\textbf{没}来参加聚会。} \emph{Erreur : Passé négatif = \zh{没} (+ V), pas \zh{不}.}
c) \zh{✗ ...而且咖啡也喜欢。} $\rightarrow$ \zh{✅ ...而且\textbf{也喜欢喝咖啡} / \textbf{也喜欢咖啡}.} \emph{Erreur : Ordre mots. Structure 2\textsuperscript{e} prop. = Sujet (omis) + 而且 + (也/还) + Verbe + Objet. Objet \zh{咖啡} ne peut pas précéder le verbe \zh{喜欢} ici.}
d) \zh{✗ 不但老师来了，学生来了。} $\rightarrow$ \zh{✅ 不但老师来了，\textbf{而且}学生也来了。} \emph{Erreur : Corrélatif \zh{而且} manquant. Ajout de \zh{也} naturel.}

\textbf{Exercice 5 : Proposition de rédaction}
\zh{我每天早上六点起床，跑步锻炼。因为早睡早起，所以身体很好。上个星期我朋友问我：“你怎么这么健康？”我说：不但我每天锻炼，而且吃很多水果、蔬菜，不吃垃圾食品。希望大家都注意健康。}
\\
(Wǒ měitiān zǎoshang liù diǎn qǐchuáng, pǎobù duànliàn. Yīnwèi zǎoshuìzǎoqǐ, suǒyǐ shēntǐ hěn hǎo. Shàng ge xīngqī wǒ péngyou wèn wǒ: « Nǐ zěnme zhème jiànkāng? » Wǒ shuō: Búdàn wǒ měitiān duànliàn, érqiě chī hěn duō shuǐguǒ, shūcài, bù chī lājī shípǐn. Xīwàng dàjiā dōu zhùyì jiànkāng.)
\emph{Vérification :} \zh{怎么了} absent (remplacé par \zh{怎么} « comment » dans \zh{你怎么这么健康} — acceptable car \zh{怎么} + adj = « comment/quoi si... »). Pour respecter la contrainte stricte \zh{那些了}, modifier : \zh{上个星期我生病了，朋友问：“你怎么了？”... 因为……所以…… 好转了。不但……而且……}.
\emph{Version corrigée contrainte :} \zh{上周我生病了，朋友问：“你怎么了？”因为我按时吃药、休息，所以好了很快。现在我不但天天锻炼，而且早睡早起，身体真好。}
\end{corrige}

\newpage

\coursec{Exercices}

\courssub{Je vérifie mes connaissances}

\begin{exercice}[QCM : la bonne structure]
Pour chaque phrase, choisis la bonne réponse (A, B ou C).

1. 你怎么了？—— 我头疼，\zh{……} 我今天不能去上课。\\
A. 因为 \quad B. 所以 \quad C. 而且

2. 他\zh{……}会说汉语，\zh{……}会写汉字。\\
A. 因为……所以…… \quad B. 不但……而且…… \quad C. 怎么……怎么了……

3. \zh{……}昨天下雨了，\zh{……}足球比赛取消了。\\
A. 不但……而且…… \quad B. 因为……所以…… \quad C. 怎么……了……

4. 你的脸色不太好，你\zh{……}？\\
A. 怎么 \quad B. 怎么了 \quad C. 为什么了
\end{exercice}

\begin{exercice}[Vrai ou faux — justifie ta réponse]
Dis si chaque phrase est correcte (V) ou incorrecte (F), puis justifie en une phrase en français.

1. \zh{不但他学习努力，而且成绩很好。}\\
\emph{Bùdàn tā xuéxí nǔlì, érqiě chéngjì hěn hǎo.}

2. \zh{因为他病了，所以他没来上课。}\\
\emph{Yīnwèi tā bìng le, suǒyǐ tā méi lái shàngkè.}

3. \zh{你怎么了？你看起来很累。}\\
\emph{Nǐ zěnme le? Nǐ kànqǐlai hěn lèi.}

4. \zh{因为天气很热，我不想去市场。}\\
\emph{Yīnwèi tiānqì hěn rè, wǒ bù xiǎng qù shìchǎng.}
\end{exercice}

\begin{exercice}[Vocabulaire : complète le tableau]
Recopie et complète le tableau suivant en donnant le pinyin et la traduction française de chaque mot.

\begin{center}
\begin{tabularx}{\linewidth}{|X|X|X|X|}
\hline
\rowcolor{popblueL} Caractères & Pinyin & Nature & Traduction \\
\hline
怎么 & \ldots\ldots & \ldots\ldots & \ldots\ldots \\
\hline
而且 & \ldots\ldots & \ldots\ldots & \ldots\ldots \\
\hline
因为 & \ldots\ldots & \ldots\ldots & \ldots\ldots \\
\hline
所以 & \ldots\ldots & \ldots\ldots & \ldots\ldots \\
\hline
取消 & \ldots\ldots & \ldots\ldots & \ldots\ldots \\
\hline
\end{tabularx}
\end{center}
\end{exercice}

\begin{exercice}[Associe le pinyin aux caractères]
Associe chaque groupe de caractères (1--5) au bon pinyin (A--E).

\begin{center}
\begin{tabularx}{\linewidth}{|X|X|}
\hline
\rowcolor{popblueL} Caractères & Pinyin \\
\hline
1. 怎么了 & A. suǒyǐ \\
\hline
2. 不但 & B. zěnme le \\
\hline
3. 而且 & C. érqiě \\
\hline
4. 因为 & D. bùdàn \\
\hline
5. 所以 & E. yīnwèi \\
\hline
\end{tabularx}
\end{center}
\end{exercice}

\courssub{Je m'entraîne}

\begin{exercice}[Traduis en chinois]
Traduis les phrases suivantes en chinois (caractères obligatoires, pinyin souhaité).

1. Qu'est-ce que tu as ? Tu as l'air malade.
2. Non seulement il parle français, mais il parle aussi anglais.
3. Parce qu'il a beaucoup plu hier, je ne suis pas allé à l'école.
4. Non seulement le sport est bon pour la santé, mais il aide aussi à mieux étudier.
\end{exercice}

\begin{exercice}[Combine les phrases]
Relie chaque paire de phrases avec la corrélation appropriée (\zh{不但……而且……} ou \zh{因为……所以……}).

1. \zh{他会说汉语。/ 他会写汉字。}\\
\emph{Tā huì shuō Hànyǔ. / Tā huì xiě Hànzì.}

2. \zh{昨天下雨了。/ 比赛取消了。}\\
\emph{Zuótiān xiàyǔ le. / Bǐsài qǔxiāo le.}

3. \zh{雅温得的水果很便宜。/ 雅温得的水果很新鲜。}\\
\emph{Yǎwēndé de shuǐguǒ hěn piányi. / Yǎwēndé de shuǐguǒ hěn xīnxiān.}

4. \zh{我弟弟生病了。/ 他今天不能去学校。}\\
\emph{Wǒ dìdi shēngbìng le. / Tā jīntiān bù néng qù xuéxiào.}
\end{exercice}

\begin{exercice}[Texte à trous : les mots-outils]
Complète le texte suivant avec : \zh{怎么了 / 不但 / 而且 / 因为 / 所以}.

\begin{textebox}[Au dispensaire de Buea]
医生：小王，你\zh{……}？\\
\emph{Yīshēng: Xiǎo Wáng, nǐ \ldots?}\\
小王：\zh{……}我昨天晚上淋雨了，\zh{……}今天头疼。\\
\emph{Xiǎo Wáng: \ldots wǒ zuótiān wǎnshang lín yǔ le, \ldots jīntiān tóu téng.}\\
医生：你\zh{……}头疼，\zh{……}发烧。你要多休息，多喝水。\\
\emph{Yīshēng: Nǐ \ldots tóu téng, \ldots fāshāo. Nǐ yào duō xiūxi, duō hē shuǐ.}
\end{textebox}
\end{exercice}

\begin{exercice}[Remets les mots dans l'ordre]
Remets les mots de chaque groupe dans le bon ordre pour former une phrase correcte, puis traduis-la en français.

1. 学习 / 他 / 不但 / 努力 / 而且 / 很好 / 成绩

2. 因为 / 所以 / 病了 / 他 / 没来 / 上课 / 他

3. 了 / 你 / 怎么 / ？ / 不舒服 / 我

4. 杜阿拉 / 不但 / 大 / 而且 / 很 / 热闹 / 很
\end{exercice}

\courssub{Je me prépare au Bac}

\begin{exercice}[Compréhension de texte et questions de langue]
Lis le texte suivant, puis réponds aux questions.

\begin{textebox}[La santé des élèves]
小李是雅温得一所中学的学生。最近他看起来很累，老师问他：「小李，你怎么了？」\\
\emph{Xiǎo Lǐ shì Yǎwēndé yì suǒ zhōngxué de xuésheng. Zuìjìn tā kànqǐlai hěn lèi, lǎoshī wèn tā: «Xiǎo Lǐ, nǐ zěnme le?»}\\
小李说：「因为我每天晚上睡得太晚，所以白天上课没有精神。」\\
\emph{Xiǎo Lǐ shuō: «Yīnwèi wǒ měitiān wǎnshang shuì de tài wǎn, suǒyǐ báitiān shàngkè méiyǒu jīngshen.»}\\
老师告诉他：「睡眠不但对身体好，而且能帮助学习。你应该早点睡觉，多做运动。」\\
\emph{Lǎoshī gàosu tā: «Shuìmián bùdàn duì shēntǐ hǎo, érqiě néng bāngzhù xuéxí. Nǐ yīnggāi zǎo diǎn shuìjiào, duō zuò yùndòng.»}\\
小李听了老师的话。现在他不但每天早睡早起，而且常常跑步，身体越来越好了。\\
\emph{Xiǎo Lǐ tīng le lǎoshī de huà. Xiànzài tā bùdàn měitiān zǎo shuì zǎo qǐ, érqiě chángcháng pǎobù, shēntǐ yuè lái yuè hǎo le.}
\end{textebox}

\textbf{Questions de compréhension} (réponds en français) :

1. Pourquoi le professeur pose-t-il une question à Xiao Li ?
2. Quelle est la cause de la fatigue de Xiao Li ?
3. Quels conseils le professeur lui donne-t-il ?
4. Quels changements observe-t-on chez Xiao Li à la fin du texte ?

\textbf{Questions de langue} :

5. Releve dans le texte une phrase contenant la structure \zh{因为……所以……} et traduis-la en français.
6. Releve dans le texte les deux phrases contenant \zh{不但……而且……} et traduis-les.
7. Construis une phrase avec \zh{怎么了} pour demander à un camarade ce qu'il a.
\end{exercice}

\begin{exercice}[Thème et version]
\textbf{Thème.} Traduis en chinois (caractères obligatoires) :

1. Non seulement le sport est bon pour la santé, mais il aide aussi à mieux étudier.
2. Parce qu'il a beaucoup plu hier, je ne suis pas allé à l'école.
3. Qu'est-ce que tu as ? Tu ne manges rien ce matin.

\textbf{Version.} Traduis en français le texte suivant :

\begin{textebox}[Version]
小王怎么了？他今天没有来学校。因为他感冒了，所以他在家休息。\\
\emph{Xiǎo Wáng zěnme le? Tā jīntiān méiyǒu lái xuéxiào. Yīnwèi tā gǎnmào le, suǒyǐ tā zài jiā xiūxi.}\\
小王不但学习努力，而且喜欢帮助同学，我们都很喜欢他。\\
\emph{Xiǎo Wáng bùdàn xuéxí nǔlì, érqiě xǐhuan bāngzhù tóngxué, wǒmen dōu hěn xǐhuan tā.}
\end{textebox}
\end{exercice}

\begin{exercice}[Expression écrite : au marché, un ami fatigué]
\textbf{Sujet.} Tu rencontres ton ami chinois au marché de Mokolo (Yaoundé). Il a l'air très fatigué. Il te demande : \zh{你怎么了？} Non — c'est toi qui lui poses la question, et il te répond. Rédige en chinois un court dialogue de \textbf{6 à 8 répliques} (environ \textbf{60 à 80 caractères}) dans lequel :

\begin{itemize}
\item tu demandes à ton ami ce qu'il a (\zh{怎么了}) ;
\item il explique la cause de sa fatigue avec \zh{因为……所以……} ;
\item il ajoute une information avec \zh{不但……而且……} ;
\item tu lui donnes un conseil pour terminer.
\end{itemize}

\textbf{Contraintes.} Utilise au moins une fois chacune des trois structures de la leçon ; écris uniquement en caractères chinois ; soigne l'ordre des mots et la ponctuation chinoise （，。？）.
\end{exercice}

\newpage

\coursec{Corrigés des exercices}

\coursec{Corrigés des exercices — Leçon 13 : Interrogation avec 怎么了 et 怎么 ; 不但……而且…… ; 因为……所以……}

\courssub{Exercice 1 — QCM : la bonne structure}

\begin{corrige}[Exercice 1 — QCM : la bonne structure]

\textbf{1. B — 所以}\\
\zh{你怎么了？—— 我头疼，所以我今天不能去上课。}\\
\emph{Nǐ zěnme le? —— Wǒ tóuténg, suǒyǐ wǒ jīntiān bù néng qù shàngkè.}\\
« Qu'est-ce que tu as ? — J'ai mal à la tête, donc je ne peux pas aller en cours aujourd'hui. »\\
Justification : la seconde partie exprime la \cle{conséquence} de la maladie ; seul \zh{所以} convient. \zh{因为} introduirait la cause (sens inverse) et \zh{而且} exige \zh{不但} dans la première proposition.

\textbf{2. B — 不但……而且……}\\
\zh{他不但会说汉语，而且会写汉字。}\\
\emph{Tā bùdàn huì shuō Hànyǔ, érqiě huì xiě Hànzì.}\\
« Non seulement il sait parler chinois, mais il sait aussi écrire les caractères. »\\
Justification : on additionne deux compétences du même sujet ; c'est la corrélation d'\cle{addition} \zh{不但……而且……}.

\textbf{3. B — 因为……所以……}\\
\zh{因为昨天下雨了，所以足球比赛取消了。}\\
\emph{Yīnwèi zuótiān xiàyǔ le, suǒyǐ zúqiú bǐsài qǔxiāo le.}\\
« Parce qu'il a plu hier, le match de football a été annulé. »\\
Justification : la pluie est la \cle{cause}, l'annulation la \cle{conséquence}.

\textbf{4. B — 怎么了}\\
\zh{你的脸色不太好，你怎么了？}\\
\emph{Nǐ de liǎnsè bù tài hǎo, nǐ zěnme le?}\\
« Tu n'as pas bonne mine, qu'est-ce que tu as ? »\\
Justification : \zh{怎么了} sert à demander ce qui ne va pas (état physique ou moral) ; \zh{怎么} seul signifie « comment » ; \zh{为什么了} n'existe pas (\zh{了} ne se combine pas avec \zh{为什么}).
\end{corrige}

\courssub{Exercice 2 — Vrai ou faux}

\begin{corrige}[Exercice 2 — Vrai ou faux]

\textbf{1. F — incorrect.} \zh{不但他学习努力，而且成绩很好。}\\
Quand les deux propositions ont le \cle{même sujet}, celui-ci se place \textbf{avant} \zh{不但} : \zh{他不但学习努力，而且成绩很好。} (\emph{Tā bùdàn xuéxí nǔlì, érqiě chéngjì hěn hǎo.} « Non seulement il étudie avec effort, mais ses résultats sont aussi très bons. ») La place de \zh{不但} avant le sujet n'est possible que si les deux propositions ont des sujets différents (ex. : \zh{不但他努力，而且老师也负责}).

\textbf{2. V — correct.} \zh{因为他病了，所以他没来上课。}\\
\emph{Yīnwèi tā bìng le, suǒyǐ tā méi lái shàngkè.} « Parce qu'il était malade, il n'est pas venu en cours. » La corrélation \zh{因为……所以……} est bien formée : cause puis conséquence, sujet répété dans la seconde proposition.

\textbf{3. V — correct.} \zh{你怎么了？你看起来很累。}\\
\emph{Nǐ zěnme le? Nǐ kànqǐlai hěn lèi.} « Qu'est-ce que tu as ? Tu as l'air très fatigué. » \zh{怎么了} est bien employé pour s'enquérir de l'état de quelqu'un.

\textbf{4. V — correct.} \zh{因为天气很热，我不想去市场。}\\
\emph{Yīnwèi tiānqì hěn rè, wǒ bù xiǎng qù shìchǎng.} « Parce qu'il fait très chaud, je ne veux pas aller au marché. » En chinois, contrairement au français, on peut utiliser \zh{因为} seul, sans \zh{所以} (l'inverse est aussi possible : \zh{所以} seul, la cause étant implicite ou énoncée précédemment).
\end{corrige}

\courssub{Exercice 3 — Vocabulaire : complète le tableau}

\begin{corrige}[Exercice 3 — Vocabulaire : complète le tableau]

\begin{center}
\begin{tabularx}{\linewidth}{|X|X|X|X|}
\hline
\rowcolor{popblueL} Caractères & Pinyin & Nature & Traduction \\
\hline
怎么 & zěnme & pronom interrogatif & comment ; comment se fait-il que \\
\hline
怎么了 & zěnme le & locution interrogative & qu'est-ce qu'il y a ? ; qu'est-ce qui ne va pas ? \\
\hline
而且 & érqiě & conjonction & et aussi, mais aussi, de plus \\
\hline
因为 & yīnwèi & conjonction & parce que \\
\hline
所以 & suǒyǐ & conjonction & donc, c'est pourquoi \\
\hline
取消 & qǔxiāo & verbe & annuler \\
\hline
\end{tabularx}
\end{center}

\textbf{Points de vigilance} :
\begin{itemize}
\item \zh{怎么} porte le 3\ieme{} ton suivi d'un ton léger (\emph{zěnme}) ; dans \zh{怎么了}, \zh{了} est ton neutre.
\item \zh{而且} s'écrit avec deux 3\iemes{} tons mais se prononce \emph{érqiě} (le premier \zh{且} devient 2\ieme{} ton par sandhi tonal).
\item \zh{因为} : attention, \zh{为} se lit ici \emph{wèi} (4\ieme{} ton), pas \emph{wéi} (2\ieme{} ton, sens « être / faire »).
\item \zh{取消} : \zh{取} (3\ieme{} ton) + \zh{消} (1\ieme{} ton).
\end{itemize}
\end{corrige}

\courssub{Exercice 4 — Associe le pinyin aux caractères}

\begin{corrige}[Exercice 4 — Associe le pinyin aux caractères]

1 — B : \zh{怎么了} = \emph{zěnme le} (« qu'est-ce qu'il y a ? »)\\
2 — D : \zh{不但} = \emph{bùdàn} (« non seulement »)\\
3 — C : \zh{而且} = \emph{érqiě} (« et aussi, mais aussi »)\\
4 — E : \zh{因为} = \emph{yīnwèi} (« parce que »)\\
5 — A : \zh{所以} = \emph{suǒyǐ} (« donc »)
\end{corrige}

\courssub{Exercice 5 — Traduis en chinois}

\begin{corrige}[Exercice 5 — Traduis en chinois]

\textbf{1.} \zh{你怎么了？你看起来病了。}\\
\emph{Nǐ zěnme le? Nǐ kànqǐlai bìng le.}\\
(On peut aussi écrire : \zh{你看起来很不舒服。} \emph{Nǐ kànqǐlai hěn bù shūfu.} « Tu as l'air très malade / pas bien. »)

\textbf{2.} \zh{他不但会说法语，而且会说英语。}\\
\emph{Tā bùdàn huì shuō Fǎyǔ, érqiě huì shuō Yīngyǔ.}\\
Même sujet dans les deux propositions : le sujet \zh{他} se place avant \zh{不但}.

\textbf{3.} \zh{因为昨天下了很大的雨，所以我没去上学。}\\
\emph{Yīnwèi zuótiān xià le hěn dà de yǔ, suǒyǐ wǒ méi qù shàngxué.}\\
Attention : pour une action passée non accomplie, la négation est \zh{没（有）}, jamais \zh{不} ; et pas de \zh{了} après \zh{没去}.

\textbf{4.} \zh{运动不但对健康好，而且能帮助更好地学习。}\\
\emph{Yùndòng bùdàn duì jiànkāng hǎo, érqiě néng bāngzhù gèng hǎo de xuéxí.}\\
Structure \zh{对……好} « être bon pour… » ; \zh{更好地学习} : ici \zh{地} (de) est le mot-outil qui relie l'adverbe \zh{更好} au verbe \zh{学习}.
\end{corrige}

\courssub{Exercice 6 — Combine les phrases}

\begin{corrige}[Exercice 6 — Combine les phrases]

\textbf{1.} Addition de deux compétences du même sujet :\\
\zh{他不但会说汉语，而且会写汉字。}\\
\emph{Tā bùdàn huì shuō Hànyǔ, érqiě huì xiě Hànzì.}\\
« Non seulement il sait parler chinois, mais il sait aussi écrire les caractères. »

\textbf{2.} Cause puis conséquence :\\
\zh{因为昨天下雨了，所以比赛取消了。}\\
\emph{Yīnwèi zuótiān xiàyǔ le, suǒyǐ bǐsài qǔxiāo le.}\\
« Parce qu'il a plu hier, le match a été annulé. »

\textbf{3.} Deux qualités du même sujet (contexte camerounais) :\\
\zh{雅温得的水果不但很便宜，而且很新鲜。}\\
\emph{Yǎwēndé de shuǐguǒ bùdàn hěn piányi, érqiě hěn xīnxiān.}\\
« Les fruits de Yaoundé sont non seulement très bon marché, mais aussi très frais. »

\textbf{4.} Cause puis conséquence :\\
\zh{因为我弟弟生病了，所以他今天不能去学校。}\\
\emph{Yīnwèi wǒ dìdi shēngbìng le, suǒyǐ tā jīntiān bù néng qù xuéxiào.}\\
« Parce que mon petit frère est tombé malade, il ne peut pas aller à l'école aujourd'hui. »
\end{corrige}

\courssub{Exercice 7 — Texte à trous : les mots-outils}

\begin{corrige}[Exercice 7 — Texte à trous : les mots-outils]

\textbf{Texte complété :}

\begin{textebox}[Dialogue chez le médecin]
医生：小王，你\cle{怎么了}？\\
\emph{Yīshēng: Xiǎo Wáng, nǐ zěnme le?}\\
小王：\cle{因为}我昨天晚上淋雨了，\cle{所以}今天头疼。\\
\emph{Xiǎo Wáng: Yīnwèi wǒ zuótiān wǎnshang lín yǔ le, suǒyǐ jīntiān tóuténg.}\\
医生：你\cle{不但}头疼，\cle{而且}发烧。你要多休息，多喝水。\\
\emph{Yīshēng: Nǐ bùdàn tóuténg, érqiě fāshāo. Nǐ yào duō xiūxi, duō hē shuǐ.}
\end{textebox}

\textbf{Traduction :} « Le médecin : Xiao Wang, qu'est-ce que tu as ? — Xiao Wang : Parce que j'ai été sous la pluie hier soir, j'ai mal à la tête aujourd'hui. — Le médecin : Non seulement tu as mal à la tête, mais tu as aussi de la fièvre. Tu dois beaucoup te reposer et boire beaucoup d'eau. »

\textbf{Justification :} le médecin s'enquiert de l'état du patient (\zh{怎么了}) ; le patient enchaîne cause (\zh{因为}) et conséquence (\zh{所以}) ; le médecin additionne deux symptômes du même sujet (\zh{不但……而且……}, sujet \zh{你} placé avant \zh{不但}).
\end{corrige}

\courssub{Exercice 8 — Remets les mots dans l'ordre}

\begin{corrige}[Exercice 8 — Remets les mots dans l'ordre]

\textbf{1.} \zh{他不但学习努力，而且成绩很好。}\\
\emph{Tā bùdàn xuéxí nǔlì, érqiě chéngjì hěn hǎo.}\\
« Non seulement il étudie avec effort, mais ses résultats sont aussi très bons. »\\
Piège évité : le sujet \zh{他} doit précéder \zh{不但}.

\textbf{2.} \zh{因为他病了，所以他没来上课。}\\
\emph{Yīnwèi tā bìng le, suǒyǐ tā méi lái shàngkè.}\\
« Parce qu'il était malade, il n'est pas venu en cours. »\\
Ordre : \zh{因为} + cause, \zh{所以} + conséquence ; le sujet est répété dans chaque proposition.

\textbf{3.} \zh{你怎么了？我不舒服。}\\
\emph{Nǐ zěnme le? Wǒ bù shūfu.}\\
« Qu'est-ce que tu as ? — Je ne me sens pas bien. »\\
\zh{了} se colle à \zh{怎么} : \zh{怎么了} ; le point d'interrogation chinois \zh{？} termine la question.

\textbf{4.} \zh{杜阿拉不但很大，而且很热闹。}\\
\emph{Dùālā bùdàn hěn dà, érqiě hěn rènao.}\\
« Douala est non seulement très grande, mais aussi très animée. »\\
Deux adjectifs qualifiant le même sujet : \zh{不但} et \zh{而且} se placent chacun devant le groupe \zh{很 + adjectif}.
\end{corrige}

\courssub{Exercice 9 — Compréhension de texte et questions de langue}

\begin{corrige}[Exercice 9 — Compréhension de texte et questions de langue]

\textbf{Questions de compréhension}

\textbf{1.} Parce que Xiao Li a l'air très fatigué ces derniers temps (\zh{最近他看起来很累}) ; le professeur s'inquiète de son état et lui demande ce qu'il a.

\textbf{2.} La cause de sa fatigue est qu'il se couche trop tard tous les soirs (\zh{每天晚上睡得太晚}) ; il n'a donc pas d'énergie en classe la journée.

\textbf{3.} Le professeur lui conseille de se coucher plus tôt (\zh{早点睡觉}) et de faire plus de sport (\zh{多做运动}), en expliquant que le sommeil est bon pour le corps et aide à étudier.

\textbf{4.} À la fin, Xiao Li suit les conseils : il se couche et se lève tôt chaque jour, il court souvent, et sa santé s'améliore de plus en plus (\zh{身体越来越好}).

\textbf{Questions de langue}

\textbf{5.} \zh{因为我每天晚上睡得太晚，所以白天上课没有精神。}\\
\emph{Yīnwèi wǒ měitiān wǎnshang shuì de tài wǎn, suǒyǐ báitiān shàngkè méiyǒu jīngshen.}\\
« Parce que je me couche trop tard tous les soirs, je n'ai pas d'énergie en classe la journée. »

\textbf{6.} Première phrase : \zh{睡眠不但对身体好，而且能帮助学习。}\\
\emph{Shuìmián bùdàn duì shēntǐ hǎo, érqiě néng bāngzhù xuéxí.}\\
« Le sommeil est non seulement bon pour le corps, mais il aide aussi à étudier. »\\
Deuxième phrase : \zh{现在他不但每天早睡早起，而且常常跑步。}\\
\emph{Xiànzài tā bùdàn měitiān zǎo shuì zǎo qǐ, érqiě chángcháng pǎobù.}\\
« Maintenant, non seulement il se couche et se lève tôt chaque jour, mais il court aussi souvent. »

\textbf{7.} Exemple : \zh{你怎么了？你看起来不太舒服。}\\
\emph{Nǐ zěnme le? Nǐ kànqǐlai bù tài shūfu.} « Qu'est-ce que tu as ? Tu n'as pas l'air bien. »

\textbf{Barème indicatif (sur 20)} : compréhension 8 pts (2 pts par question) ; questions de langue 9 pts (3 pts par question : relevé exact 1 pt, traduction fidèle 2 pts) ; qualité du français 3 pts.

\textbf{Points de vigilance} :
\begin{itemize}
\item Ne pas traduire \zh{越来越} par « de plus en plus » au mauvais endroit de la phrase française (place après le verbe/adjectif en chinois, avant en français : \zh{身体越来越好} → « la santé devient de mieux en mieux / s'améliore de plus en plus »).
\item \zh{没有精神} signifie « sans énergie, sans entrain, apathique », pas « sans esprit » (au sens intellectuel).
\item Dans \zh{睡得太晚}, le mot-outil \zh{得} introduit le complément de degré (manière/résultat) après le verbe.
\end{itemize}
\end{corrige}

\courssub{Exercice 10 — Thème et version}

\begin{corrige}[Exercice 10 — Thème et version]

\textbf{Thème (Français → Chinois)}

\textbf{1.} \zh{运动不但对健康好，而且能帮助更好地学习。}\\
\emph{Yùndòng bùdàn duì jiànkāng hǎo, érqiě néng bāngzhù gèng hǎo de xuéxí.}

\textbf{2.} \zh{因为昨天下了很大的雨，所以我没去上学。}\\
\emph{Yīnwèi zuótiān xià le hěn dà de yǔ, suǒyǐ wǒ méi qù shàngxué.}

\textbf{3.} \zh{你怎么了？今天早上你什么都没吃。}\\
\emph{Nǐ zěnme le? Jīntiān zǎoshang nǐ shénme dōu méi chī.}\\
(Structure \zh{什么……都……} « rien du tout » ; négation au passé : \zh{没}.)

\textbf{Version (Chinois → Français)}

\textbf{Texte source :}
\begin{textebox}[Texte à traduire]
小王怎么了？他今天没来上学。因为他感冒了，所以他在家休息。小王不但勤奋，而且喜欢帮助同学，我们都很喜欢他。
\end{textebox}

\textbf{Traduction attendue :}
« Qu'est-ce qu'a Xiao Wang ? Il n'est pas venu à l'école aujourd'hui. Parce qu'il est enrhumé, il se repose à la maison. Xiao Wang est non seulement travailleur, mais il aime aussi aider ses camarades ; nous l'aimons tous beaucoup. »

\textbf{Barème indicatif (sur 20)} : thème 10 pts (exactitude des structures 6 pts, caractères et ordre des mots 4 pts) ; version 8 pts (fidélité du sens 5 pts, qualité du français 3 pts) ; présentation 2 pts.

\textbf{Points de vigilance} :
\begin{itemize}
\item Au thème, ne jamais écrire \zh{不去了} pour « ne suis pas allé » — le passé négatif exige \zh{没（有）} sans \zh{了}.
\item \zh{感冒} (\emph{gǎnmào}) « être enrhumé / avoir un rhume » se traduit par un état, pas par « attraper froid » au sens littéral.
\item \zh{同学} = « camarades de classe » ; \zh{勤奋} = « travailleur, assidu ».
\end{itemize}
\end{corrige}

\courssub{Exercice 11 — Expression écrite : au marché, un ami fatigué}

\begin{corrige}[Exercice 11 — Expression écrite : au marché, un ami fatigué]

\textbf{Proposition de rédaction modèle} (7 répliques, environ 75 caractères chinois) :

\begin{textebox}[Au marché de Mokolo]
我：小王，你怎么了？你看起来很累。\\
\emph{Wǒ: Xiǎo Wáng, nǐ zěnme le? Nǐ kànqǐlai hěn lèi.}\\
小王：因为我昨天晚上睡得太晚，所以今天没有精神。\\
\emph{Xiǎo Wáng: Yīnwèi wǒ zuótiān wǎnshang shuì de tài wǎn, suǒyǐ jīntiān méiyǒu jīngshen.}\\
我：你昨天做什么了？\\
\emph{Wǒ: Nǐ zuótiān zuò shénme le?}\\
小王：我不但要准备考试，而且要帮妈妈做事。\\
\emph{Xiǎo Wáng: Wǒ bùdàn yào zhǔnbèi kǎoshì, érqiě yào bāng māma zuòshì.}\\
我：你太累了！你应该早点睡觉。\\
\emph{Wǒ: Nǐ tài lèi le! Nǐ yīnggāi zǎo diǎn shuìjiào.}\\
小王：你说得对，谢谢你。\\
\emph{Xiǎo Wáng: Nǐ shuō de duì, xièxie nǐ.}\\
我：不客气，我们走吧！\\
\emph{Wǒ: Bú kèqi, wǒmen zǒu ba!}
\end{textebox}

\textbf{Traduction :} « Moi : Xiao Wang, qu'est-ce que tu as ? Tu as l'air très fatigué. — Xiao Wang : Parce que je me suis couché trop tard hier soir, je n'ai pas d'énergie aujourd'hui. — Moi : Qu'as-tu fait hier ? — Xiao Wang : Non seulement je dois préparer l'examen, mais je dois aussi aider maman. — Moi : Tu es trop fatigué ! Tu devrais te coucher plus tôt. — Xiao Wang : Tu as raison, merci. — Moi : De rien, allons-y ! »

\textbf{Barème indicatif (sur 20)} : respect des contraintes (3 structures présentes, 6 à 8 répliques, 60 à 80 caractères) 6 pts ; correction grammaticale (ordre des mots, place de \zh{不但}, \zh{因为……所以……}) 6 pts ; exactitude des caractères 4 pts ; ponctuation chinoise et présentation en dialogue 2 pts ; richesse et naturel du vocabulaire 2 pts.

\textbf{Points de vigilance} :
\begin{itemize}
\item \zh{怎么了} doit apparaître dans la question d'ouverture, avec le point d'interrogation chinois \zh{？} ;
\item dans \zh{不但……而且……} avec un seul sujet, le sujet se place avant \zh{不但} : \zh{我不但要……而且要……} ;
\item ne pas écrire \zh{因为……} sans conséquence logique claire, ni mélanger \zh{因为} avec \zh{而且} dans la même corrélation ;
\item écrire uniquement en caractères : ni pinyin ni français dans la copie finale ;
\item compter les caractères : rester entre 60 et 80, hors ponctuation.
\end{itemize}
\end{corrige}

\newpage

\coursec{Fiche bilan}

\begin{popfigure}
\begin{tikzpicture}[mindmap, grow cyclic, every node/.style={concept, font=\small},
  root concept/.append style={concept color=popblue, font=\bfseries},
  level 1/.append style={level distance=3.2cm, sibling angle=60},
  level 1 concept/.append style={font=\footnotesize}]
\node[root concept, align=center]{Leçon 13\\Grammaire}
  child[concept color=poporange]{ node[align=center]{怎么了 / 怎么\\Interroger}
    child[concept color=poporangeL]{ node[font=\scriptsize, align=center]{怎么了？\\Qu'est-ce qu'il y a ?} }
    child[concept color=poporangeL]{ node[font=\scriptsize, align=center]{怎么 + V\\Comment ?} }
  }
  child[concept color=popgreen]{ node[align=center]{不但……而且……\\non seulement… mais aussi}
    child[concept color=popgreenL]{ node[font=\scriptsize, align=center]{même sujet :\\S 不但 V，而且 V} }
    child[concept color=popgreenL]{ node[font=\scriptsize, align=center]{sujets différents :\\不但 S1，而且 S2 也} }
  }
  child[concept color=poppurple]{ node[align=center]{因为……所以……\\cause → conséquence}
    child[concept color=poppurpleL]{ node[font=\scriptsize, align=center]{因为 + cause，\\所以 + résultat} }
    child[concept color=poppurpleL]{ node[font=\scriptsize, align=center]{所以 seul possible\\因为 seul possible} }
  }
  child[concept color=poppink]{ node[align=center]{Pièges\\francophones}
    child[concept color=poppinkL]{ node[font=\scriptsize]{pas de 因为…但是} }
    child[concept color=poppinkL]{ node[font=\scriptsize]{怎么 ≠ 怎么了} }
  }
  child[concept color=popgold]{ node[align=center]{Lexique\\santé / environnement}
    child[concept color=popgoldL]{ node[font=\scriptsize, align=center]{生病 医院 空气\\污染 锻炼 休息} }
  };
\end{tikzpicture}
\legende{Carte mentale de la leçon 13 : trois structures clés et leurs pièges.}
\end{popfigure}

\begin{bilan}
\textbf{1. Lexique clé à connaître}\par
\begin{center}
\begin{tabularx}{\linewidth}{|X|X|X|X|}
\hline
\rowcolor{popblueL} Caractères & Pinyin & Nature & Traduction \\
\hline
怎么 & zěnme & adv. interrogatif & comment ; pourquoi \\
\hline
怎么了 & zěnme le & expression & qu'est-ce qu'il y a ? \\
\hline
生病 & shēng bìng & v. & tomber malade \\
\hline
医院 & yīyuàn & n. & hôpital \\
\hline
空气 & kōngqì & n. & air \\
\hline
污染 & wūrǎn & n./v. & pollution ; polluer \\
\hline
锻炼 & duànliàn & v. & faire de l'exercice \\
\hline
休息 & xiūxi & v. & se reposer \\
\hline
不但……而且…… & búdàn… érqiě… & conj. & non seulement… mais aussi… \\
\hline
因为……所以…… & yīnwèi… suǒyǐ… & conj. & parce que… donc… \\
\hline
\end{tabularx}
\end{center}

\textbf{2. Règles de grammaire à connaître par cœur}\par
\begin{center}
\begin{tabularx}{\linewidth}{|X|X|X|}
\hline
\rowcolor{popblueL} Structure & Exemple & Traduction \\
\hline
怎么了 ? (inquiétude) & \zh{你怎么了？} \emph{Nǐ zěnme le ?} & Qu'est-ce que tu as ? \\
\hline
怎么 + verbe (manière) & \zh{去医院怎么走？} \emph{Qù yīyuàn zěnme zǒu ?} & Comment va-t-on à l'hôpital ? \\
\hline
S + 不但 + V/Adj，而且 + V/Adj & \zh{他不但聪明，而且努力。} \emph{Tā búdàn cōngming, érqiě nǔlì.} & Il est non seulement intelligent, mais aussi travailleur. \\
\hline
不但 + S1，而且 + S2 + 也 & \zh{不但老师来了，而且学生也来了。} & Non seulement le professeur est venu, mais les élèves aussi. \\
\hline
因为 + cause，所以 + résultat & \zh{因为空气不好，所以他生病了。} \emph{Yīnwèi kōngqì bù hǎo, suǒyǐ tā shēng bìng le.} & Parce que l'air est mauvais, il est tombé malade. \\
\hline
\end{tabularx}
\end{center}

\textbf{3. Méthodes express}\par
\begin{itemize}
  \item Pour interroger sur un \textbf{problème} (santé, situation) : \zh{怎么了} seul, en fin de phrase, avec \zh{了}.
  \item Pour interroger sur la \textbf{manière} : \zh{怎么} placé \textbf{avant le verbe}, sans \zh{了}.
  \item Pour ajouter une information : même sujet → \zh{不但} après le sujet ; sujets différents → \zh{不但} avant le premier sujet et \zh{也} dans la seconde proposition.
  \item Pour exprimer la cause : \zh{因为} introduit la cause, \zh{所以} le résultat ; en chinois, on peut utiliser les deux ensemble (contrairement au français).
\end{itemize}
\end{bilan}

\begin{aretenir}[Checklist avant le Bac]
\begin{itemize}
  \item $\square$ Je distingue \zh{怎么了} (qu'est-ce qu'il y a ?) et \zh{怎么} + verbe (comment ?).
  \item $\square$ Je place \zh{怎么} avant le verbe et jamais en fin de phrase.
  \item $\square$ Je sais construire \zh{不但……而且……} avec un même sujet.
  \item $\square$ Je sais construire \zh{不但……而且……也……} avec deux sujets différents.
  \item $\square$ Je place correctement la virgule entre les deux propositions de la corrélation.
  \item $\square$ Je sais exprimer la cause avec \zh{因为……所以……} sans hésiter sur l'ordre.
  \item $\square$ Je n'écris jamais \zh{因为……但是……} (cause + opposition = contradiction).
  \item $\square$ Je peux remplacer \zh{所以} par \zh{所以} seul ou \zh{因为} seul si la phrase reste claire.
  \item $\square$ Je connais le lexique santé / environnement : \zh{生病、医院、空气、污染、锻炼、休息}.
  \item $\square$ Je vérifie les tons : \zh{zěnme}, \zh{búdàn}, \zh{érqiě}, \zh{yīnwèi}, \zh{suǒyǐ}.
  \item $\square$ Je sais transformer une phrase simple en phrase avec \zh{不但……而且……} ou \zh{因为……所以……}.
  \item $\square$ Je relis mes productions pour vérifier l'ordre Sujet + corrélatif + verbe.
\end{itemize}
\end{aretenir}

\findecours

\end{document}
